% Chapter 3, Section 4 _Linear Algebra_ Jim Hefferon
%  http://joshua.smcvt.edu/linearalgebra
%  2001-Jun-12
\section{Operasi Matriks}
Bagian sebelumnya menunjukkan bagaimana matriks merepresentasikan peta linear.
Sekarang kita akan menelaah bagaimana representasi ini berinteraksi dengan hal-hal
yang telah kita ketahui.
Pertama-tama kita akan melihat bagaimana
representasi perkalian skalar $r\cdot f$ dari suatu peta linear berkaitan
dengan representasi $f$, dan
juga bagaimana representasi suatu
jumlah $f+g$ berkaitan dengan representasi kedua sukunya.
Kemudian kita akan melakukan perbandingan yang sama untuk operasi peta berupa
komposisi dan invers.




\subsection{Jumlah dan Perkalian Skalar}

\begin{example}
Misalkan \( \map{f}{V}{W} \) adalah fungsi linear
yang direpresentasikan terhadap basis-basis tertentu oleh matriks ini.
\begin{equation*}
  \rep{f}{B,D}
  =
    \begin{mat}[r]
      1  &0  \\
      1  &1
    \end{mat}
\end{equation*}
Perhatikan peta yang merupakan kelipatan skalar $\map{5f}{V}{W}$.
Kita akan mengaitkan representasi $\rep{5f}{B,D}$ dengan $\rep{f}{B,D}$.

Misalkan $f$ memetakan $\vec{v}\mapsto \vec{w}$
dengan representasi berikut.
\begin{equation*}
  \rep{\vec{v}}{B}
  =\colvec{v_1 \\ v_2}
  \qquad
  \rep{\vec{w}}{D}
  =\colvec{w_1 \\ w_2}
\end{equation*}
Karena basis kodomainnya adalah $D=\sequence{\vec{\delta}_1,\vec{\delta}_2}$,
representasi tersebut menyatakan
bahwa vektor keluarannya adalah $\vec{w}=w_1\vec{\delta}_1+w_2\vec{\delta}_2$.

Tindakan peta $5f$ adalah $\vec{v}\mapsto 5\vec{w}$ dan
$5\vec{w}=5\cdot(w_1\vec{\delta}_1+w_2\vec{\delta}_2)
=(5w_1)\vec{\delta}_1+(5w_2)\vec{\delta}_2$.
Jadi $5f$ memetakan vektor masukan $\vec{v}$ ke
vektor keluaran dengan
representasi berikut.
\begin{equation*}
  \rep{5\vec{w}}{D}
  =\colvec{5w_1 \\ 5w_2}
\end{equation*}
Beralih dari peta $f$ ke peta $5f$ mengakibatkan
setiap entri dalam representasi vektor keluaran
dikalikan dengan $5$.

Karena itu, \( \rep{5f}{B,D} \) adalah matriks berikut.
\begin{equation*}
  \rep{5f}{B,D}\cdot\colvec{v_1 \\ v_2}
  =
   \colvec{5v_1 \\ 5v_1+5v_2}
  \qquad  
  \rep{5f}{B,D}
  =
    \begin{mat}[r]
      5  &0  \\
      5  &5
    \end{mat}
\end{equation*}
Dengan demikian, beralih dari matriks yang merepresentasikan $f$ ke
matriks yang merepresentasikan $5f$ berarti mengalikan semua entri matriks
dengan $5$.
\end{example}

\begin{example}
Kita dapat melakukan penelaahan serupa untuk jumlah dua peta.
Misalkan dua peta linear dengan domain dan kodomain yang sama,
\( \map{f,g}{\Re^2}{\Re^2} \), direpresentasikan
terhadap basis $B$ dan~$D$ oleh matriks-matriks berikut.
\begin{equation*}
  \rep{f}{B,D}=
    \begin{mat}
       1  &3  \\
       2  &0 
    \end{mat}
  \qquad
  \rep{g}{B,D}=
    \begin{mat}[r]
      -2  &-1 \\
       2  &4
    \end{mat}
\end{equation*}
Ingat kembali definisi jumlah: jika
$f$~memetakan $\vec{v}\mapsto\vec{u}$ dan
$g$~memetakan $\vec{v}\mapsto\vec{w}$, maka
$f+g$ adalah fungsi yang bertindak sebagai $\vec{v}\mapsto\vec{u}+\vec{w}$.
Misalkan berikut ini representasi vektor masukan dan keluarannya.
\begin{equation*}
  \rep{\vec{v}}{B}=\colvec{v_1 \\ v_2}
  \qquad
  \rep{\vec{u}}{D}=\colvec{u_1 \\ u_2}
  \quad
  \rep{\vec{w}}{D}=\colvec{w_1 \\ w_2}
\end{equation*}
Karena $D=\sequence{\vec{\delta}_1,\vec{\delta}_2}$, kita memperoleh
$\vec{u}+\vec{w}=(u_1\vec{\delta}_1+u_2\vec{\delta}_2)
  +(w_1\vec{\delta}_1+w_2\vec{\delta}_2)
  =(u_1+w_1)\vec{\delta}_1+(u_2+w_2)\vec{\delta}_2$
sehingga berikut ini adalah representasi jumlah vektor tersebut.
\begin{equation*}
  \rep{\vec{u}+\vec{w}}{D}=\colvec{u_1+w_1 \\ u_2+w_2}
\end{equation*}
Jadi, karena persamaan berikut merepresentasikan tindakan peta $f$ dan~$g$ pada
masukan~$\vec{v}$,
\begin{equation*}
    \begin{mat}
       1  &3  \\
       2  &0 
    \end{mat}
    \colvec{v_1 \\ v_2}
    =\colvec{v_1+3v_2   \\ 2v_1}
  \qquad
    \begin{mat}
      -2  &-1 \\
       2  &4
    \end{mat}
    \colvec{v_1 \\ v_2}
    =\colvec{-2v_1-v_2   \\ 2v_1+4v_2}  
\end{equation*}
menjumlahkan entri-entrinya merepresentasikan tindakan peta $f+g$.
\begin{equation*}
  \rep{f+g}{B,D}\cdot\colvec{v_1 \\ v_2}
  =
  \colvec{-v_1+2v_2 \\ 4v_1+4v_2}
\end{equation*}
Dengan demikian, kita
menghitung matriks yang merepresentasikan jumlah fungsi dengan menjumlahkan entri-entri
matriks yang merepresentasikan fungsi-fungsi tersebut.
\begin{equation*} 
    \rep{f+g}{B,D}=
    \begin{mat}
       -1  &2  \\
        4  &4 
    \end{mat}
\end{equation*}
\end{example}

\begin{definition} \label{def:SumScalarProdMats}
%<*df:SumScalarProdMats>
\definend{Kelipatan skalar}\index{kelipatan skalar!matriks}%
\index{matriks!kelipatan skalar} suatu matriks adalah
hasil perkalian skalar entri demi entri.
\definend{Jumlah\/}\index{jumlah!matriks}\index{matriks!jumlah} dua
matriks berukuran sama adalah jumlah entri demi entrinya.
%</df:SumScalarProdMats>
\end{definition}

Operasi-operasi ini memperluas operasi penjumlahan dan perkalian skalar
vektor dari bab pertama.

Kita memerlukan suatu hasil yang membuktikan bahwa operasi-operasi
matriks ini benar-benar bekerja seperti yang ditunjukkan
oleh contoh-contoh tersebut.

\begin{theorem}  \label{th:MatOpsRepMapOps}
%<*th:MatOpsRepMapOps>
Misalkan \( \map{h,g}{V}{W} \) adalah peta-peta linear yang direpresentasikan terhadap
basis \( B,D \) oleh matriks \( H \) dan \( G \), serta misalkan $r$ adalah suatu skalar.
Maka, terhadap \( B,D \), peta
\( \map{r\cdot h}{V}{W} \) direpresentasikan oleh
\( rH \),
dan peta
\( \map{h+g}{V}{W} \) direpresentasikan
oleh \( H+G \).
%</th:MatOpsRepMapOps>
\end{theorem}

\begin{proof}
Generalisasikan contoh-contoh tersebut.
Ini adalah \nearbyexercise{exer:CorrspMapMatOps}.
\end{proof}

\begin{remark}
Dua operasi pada matriks ini sederhana,
tetapi kita tidak mendefinisikannya dengan cara ini karena kesederhanaannya.
Kita mendefinisikannya
demikian karena keduanya merepresentasikan penjumlahan fungsi
dan perkalian skalar fungsi.
Artinya, program kita adalah mendefinisikan operasi matriks dengan mengacu pada
operasi fungsi.
Kesederhanaannya merupakan bonus.

% We can express this program in another way.
% Recall Theorem~III.\ref{th:MatMultRepsFuncAppl},
% which says that matrix-vector 
% multiplication represents the application of a linear map.
% Following it, Remark~III.\ref{rem:NotSurprising} notes that 
% the theorem
% justifies the definition of matrix-vector multiplication
% and so in some sense the theorem must hold\Dash if  
% the theorem didn't hold
% then we would adjust the definition until it did.
% The above \nearbytheorem{th:MatOpsRepMapOps} 
% is another example of such a result.
% It justifies the given definition of the matrix operations. 

Kita akan menjumpai hal ini lagi
dalam subbagian berikutnya, tempat kita akan mendefinisikan operasi
perkalian matriks.
Karena kita baru saja mendefinisikan perkalian skalar matriks dan jumlah matriks
sebagai operasi entri demi entri,
gagasan yang naif adalah mendefinisikan perkalian matriks sebagai hasil kali
entri demi entri.
Secara teori kita dapat berbuat sesuka hati,
tetapi kita akan bersikap praktis dan
menggabungkan entri-entri dengan cara yang
merepresentasikan operasi komposisi fungsi.
\end{remark}

Kasus khusus perkalian skalar adalah perkalian dengan nol.
Untuk setiap peta, $0\cdot h$ adalah homomorfisme nol, dan untuk setiap matriks,
$0\cdot H$ adalah matriks yang semua entrinya nol.

\begin{definition} \label{df:ZeroMatrix}
%<*df:ZeroMatrix>
Sebuah \definend{matriks nol}\index{matriks nol}\index{matriks!nol} % 
memiliki semua entri bernilai $0$.
Kita menuliskannya sebagai $Z_{\nbym{n}{m}}$ atau cukup $Z$
(notasi lain yang umum adalah $0_{\nbym{n}{m}}$ atau cukup $0$).
%</df:ZeroMatrix>
\end{definition}

\begin{example}
Peta nol dari sebarang ruang berdimensi tiga ke sebarang ruang berdimensi dua
direpresentasikan oleh matriks nol \( \nbym{2}{3} \)
\begin{equation*}
   Z=\begin{mat}[r]
     0  &0  &0  \\
     0  &0  &0
   \end{mat}
\end{equation*}
apa pun basis domain dan kodomain yang kita gunakan.
\end{example}



\begin{exercises}
  \recommended \item
    Lakukan operasi yang ditunjukkan jika terdefinisi, atau nyatakan ``tidak terdefinisi.''
    \begin{exparts}
      \partsitem \( \begin{mat}[r]
                 5  &-1  &2  \\
                 6  &1   &1
               \end{mat}
               +
               \begin{mat}[r]
                 2  &1   &4  \\
                 3  &0   &5
               \end{mat}    \)
      \partsitem \( 6\cdot\begin{mat}[r]
                 2  &-1  &-1 \\
                 1  &2   &3
               \end{mat}   \)
      \partsitem \( \begin{mat}[r]
                 2  &1  \\
                 0  &3
               \end{mat}
               +
               \begin{mat}[r]
                 2  &1  \\
                 0  &3
               \end{mat} \)
      \partsitem \( 4\begin{mat}[r]
                 1  &2  \\
                 3  &-1
               \end{mat}
               +
               5\begin{mat}[r]
                -1  &4  \\
                -2  &1
               \end{mat} \)
      \partsitem \( 3\begin{mat}[r]
                 2  &1  \\
                 3  &0
               \end{mat}
               +2
               \begin{mat}[r]
                 1  &1  &4 \\
                 3  &0  &5
               \end{mat} \)
    \end{exparts}
    \begin{answer} 
      \begin{exparts*}
        \partsitem \( \begin{mat}[r]
                   7  &0  &6  \\
                   9  &1  &6
                 \end{mat}  \)
        \partsitem \( \begin{mat}[r]
                  12  &-6 &-6 \\
                   6  &12 &18
                 \end{mat}  \)
        \partsitem \( \begin{mat}[r]
                   4  &2  \\
                   0  &6
                 \end{mat}  \)
        \partsitem \( \begin{mat}[r]
                  -1  &28 \\
                   2  &1
                 \end{mat}  \)
        \partsitem Tidak terdefinisi.
      \end{exparts*}  
   \end{answer}
  \item Berikan matriks yang merepresentasikan peta nol dari $\Re^4$ ke $\Re^2$,
   terhadap basis-basis standar.
   \begin{answer}
     Basis-basisnya tidak berpengaruh.
     Satu-satunya hal yang penting adalah dimensinya tepat.
     \begin{equation*}
       \begin{mat}
         0 &0 &0 &0 \\
         0 &0 &0 &0
       \end{mat}
     \end{equation*}
   \end{answer}
  \item \label{exer:CorrspMapMatOps}
    Buktikan \nearbytheorem{th:MatOpsRepMapOps}.
    \begin{exparts}
      \partsitem Buktikan bahwa penjumlahan matriks merepresentasikan penjumlahan peta linear.
      \partsitem Buktikan bahwa perkalian skalar matriks merepresentasikan perkalian
         skalar peta linear.
    \end{exparts}
    \begin{answer}
      Representasikan vektor domain $\vec{v}\in V$ dan peta-peta
      $\map{g,h}{V}{W}$ terhadap basis $B,D$ dengan cara yang biasa.
      \begin{exparts}
      \partsitem Representasi
         $(g+h)\,(\vec{v})=g(\vec{v})+h(\vec{v})$
        \begin{multline*}
          \bigl( (g_{1,1}v_1+\dots+g_{1,n}v_n)\vec{\delta}_1  
                  +\cdots                     
                  +(g_{m,1}v_1+\dots+g_{m,n}v_n)\vec{\delta}_m \bigr) \\ 
           +\bigl( (h_{1,1}v_1+\cdots+h_{1,n}v_n)\vec{\delta}_1    
                   +\cdots                                  
                   +(h_{m,1}v_1+\dots+h_{m,n}v_n)\vec{\delta}_m \bigr) 
         \end{multline*}
         dapat dikelompokkan ulang menjadi
         \begin{multline*}
          =((g_{1,1}+h_{1,1})v_1+\dots+
                 (g_{1,n}+h_{1,n})v_n)\cdot\vec{\delta}_1              \\       
            +\cdots                            
           +((g_{m,1}+h_{m,1})v_1+\dots+
                 (g_{m,n}+h_{m,n})v_n)\cdot\vec{\delta}_m
        \end{multline*} 
        yaitu jumlah entri demi entri
        dari representasi \( g(\vec{v}) \) dan representasi
        \( h(\vec{v}) \).
      \partsitem Representasi
        $(r\cdot h)\,(\vec{v})=r\cdot \bigl(h(\vec{v})\bigr)$
        \begin{align*}
          r\cdot 
             &\bigl(
               (h_{1,1}v_1+h_{1,2}v_2+\dots+h_{1,n}v_n)\vec{\delta}_1   \\  
             &\quad +\dots      
             +(h_{m,1}v_1+h_{m,2}v_2+\dots+h_{m,n}v_n)\vec{\delta}_m\bigr)   \\
             &=(rh_{1,1}v_1+\dots+rh_{1,n}v_n)\cdot\vec{\delta}_1         \\ 
             &\quad+\dots                                    
              +(rh_{m,1}v_1+\dots+rh_{m,n}v_n)\cdot\vec{\delta}_m
        \end{align*}
        adalah hasil kali entri demi entri
        antara \( r \) dan representasi \( h \).
     \end{exparts}  
   \end{answer}
  \recommended \item
    Buktikan masing-masing pernyataan berikut dengan mengasumsikan bahwa operasinya terdefinisi,
    dengan \( G \), \( H \), dan
    \( J \) berupa matriks,
    \( Z \) berupa matriks nol, serta \( r \) dan \( s \) berupa skalar.
    \begin{exparts}
      \partsitem Penjumlahan matriks bersifat komutatif \( G+H=H+G \).
      \partsitem Penjumlahan matriks bersifat asosiatif \( G+(H+J)=(G+H)+J \).
      \partsitem Matriks nol adalah identitas aditif \( G+Z=G \).
      \partsitem \( 0\cdot G=Z \)
      \partsitem \( (r+s)G=rG+sG \)
      \partsitem Matriks memiliki invers aditif \( G+(-1)\cdot G=Z \).
      \partsitem \( r(G+H)=rG+rH \)
      \partsitem \( (rs)G=r(sG) \)
    \end{exparts}
    \begin{answer}
      Pertama, setiap sifat ini
      mudah diperiksa secara entri demi entri.
      Sebagai contoh, tuliskan
      \begin{equation*}
        G=
        \begin{mat}
          g_{1,1}  &\ldots  &g_{1,n}  \\
          \vdots   &        &\vdots   \\
          g_{m,1}  &\ldots  &g_{m,n}       
        \end{mat}
        \qquad
        H=
        \begin{mat}
          h_{1,1}  &\ldots  &h_{1,n}  \\
          \vdots   &        &\vdots   \\
          h_{m,1}  &\ldots  &h_{m,n}       
        \end{mat}
      \end{equation*}
      maka, menurut definisi, kita memperoleh
      \begin{equation*}
        G+H=
        \begin{mat}
          g_{1,1}+h_{1,1}  &\ldots  &g_{1,n}+h_{1,n}  \\
          \vdots           &        &\vdots           \\
          g_{m,1}+h_{m,1}  &\ldots  &g_{m,n}+h_{m,n}       
        \end{mat}
      \end{equation*}
      dan
      \begin{equation*}
        H+G=
        \begin{mat}
          h_{1,1}+g_{1,1}  &\ldots  &h_{1,n}+g_{1,n}  \\
          \vdots           &        &\vdots           \\
          h_{m,1}+g_{m,1}  &\ldots  &h_{m,n}+g_{m,n}       
        \end{mat}
      \end{equation*}
      dan keduanya sama karena entri-entrinya sama,
      $g_{i,j}+h_{i,j}=h_{i,j}+g_{i,j}$.
      Artinya, masing-masing pernyataan ini mudah diperiksa hanya dengan menggunakan
      \nearbydefinition{def:SumScalarProdMats}.

      Namun, setiap sifat
      juga mudah dipahami dari sudut pandang peta yang direpresentasikan,
      dengan menerapkan \nearbytheorem{th:MatOpsRepMapOps} beserta
      definisinya.
      \begin{exparts}
        \partsitem Kedua peta $g+h$ dan $h+g$ sama karena
          $g(\vec{v})+h(\vec{v})=h(\vec{v})+g(\vec{v})$, sebab penjumlahan
          bersifat komutatif dalam setiap ruang vektor.
          Karena kedua petanya sama, keduanya harus memiliki
          representasi yang sama.
        \partsitem Seperti jawaban sebelumnya, tetapi di sini kita menggunakan fakta bahwa
          penjumlahan ruang vektor bersifat asosiatif.
        \partsitem Seperti sebelumnya, tetapi di sini kita perhatikan bahwa
          $g(\vec{v})+z(\vec{v})=g(\vec{v})+\zero=g(\vec{v})$.
        \partsitem Gunakan fakta bahwa
          $0\cdot g(\vec{v})=\zero=z(\vec{v})$.
        \partsitem Gunakan fakta bahwa
          $(r+s)\cdot g(\vec{v})=r\cdot g(\vec{v})+s\cdot g(\vec{v})$.
        \partsitem Terapkan dua butir sebelumnya dengan $r=1$ dan $s=-1$.
        \partsitem Gunakan fakta bahwa
          $r\cdot (g(\vec{v})+h(\vec{v}))=r\cdot g(\vec{v})+r\cdot h(\vec{v})$.
        \partsitem Gunakan fakta bahwa
          $(rs)\cdot g(\vec{v})=r\cdot(s\cdot g(\vec{v}))$.
      \end{exparts}
    \end{answer}
  \item 
    Tetapkan ruang domain dan kodomain.
    Secara umum, satu
    matriks dapat merepresentasikan banyak peta berbeda terhadap basis yang berbeda.
    Namun, buktikan bahwa matriks nol hanya merepresentasikan peta nol.
    Adakah matriks lain yang memiliki sifat ini?
    \begin{answer}
      Untuk sebarang \( V,W \) dengan basis \( B,D \),
      matriks nol (dengan ukuran yang sesuai) merepresentasikan peta berikut.
      \begin{equation*}
        \vec{\beta}_1\mapsto 0\cdot\vec{\delta}_1+\dots+0\cdot\vec{\delta}_m
        \quad\cdots\quad
        \vec{\beta}_n\mapsto 0\cdot\vec{\delta}_1+\dots+0\cdot\vec{\delta}_m
      \end{equation*}
      Ini adalah peta nol.

      Tidak ada matriks lain yang hanya merepresentasikan satu peta.
      Untuk melihatnya, misalkan \( H \)
      bukan matriks nol.
      Matriks itu memiliki entri tak nol; anggap \( h_{i,j}\neq 0 \).
      Terhadap basis $B,D$, matriks tersebut merepresentasikan \( \map{h_1}{V}{W} \)
      yang memetakan
      \begin{equation*}
        \vec{\beta}_j\mapsto
        h_{1,j}\vec{\delta}_1+\dots+h_{i,j}\vec{\delta}_i
          +\dots+h_{m,j}\vec{\delta}_m
      \end{equation*}
      dan terhadap $B,2\cdot D$ matriks itu juga merepresentasikan
      \( \map{h_2}{V}{W} \) yang memetakan
      \begin{equation*}
        \vec{\beta}_j\mapsto
        h_{1,j}\cdot(2\vec{\delta}_1)+\dots+h_{i,j}\cdot(2\vec{\delta}_i)
          +\dots+h_{m,j}\cdot(2\vec{\delta}_m)
      \end{equation*}
      (notasi $2\cdot D$ berarti menggandakan
      semua anggota $D$).
      Mudah dilihat bahwa kedua peta ini tidak sama.
     \end{answer}
  \recommended \item \label{exer:LinMapsIsoMatSp}
    Misalkan \( V \) dan \( W \) adalah ruang vektor berdimensi \( n \) dan \( m \).
    Tunjukkan bahwa ruang \( \linmaps{V}{W} \) yang terdiri atas peta linear dari \( V \) ke
    \( W \) isomorfik dengan \( \matspace_{\nbym{m}{n}} \).
    \begin{answer}
      Tetapkan basis \( B \) dan \( D \) untuk \( V \) dan \( W \), lalu perhatikan
      \( \map{\mbox{Rep}_{B,D}}{\linmaps{V}{W}}{\matspace_{\nbym{m}{n}}} \)
      yang memasangkan setiap peta linear dengan matriks yang merepresentasikan peta itu,
      $h\mapsto \rep{h}{B,D}$.
      Dari bagian sebelumnya kita mengetahui bahwa (untuk basis yang tetap)
      matriks-matriks berkorespondensi dengan peta linear,
      sehingga peta representasi tersebut satu-satu dan pada.
      Bahwa peta ini mempertahankan operasi linear dinyatakan oleh
      \nearbytheorem{th:MatOpsRepMapOps}.
    \end{answer}
  \recommended \item 
    Tunjukkan bahwa dari soal sebelumnya mengikuti bahwa
    untuk sebarang enam transformasi
    \( \map{t_1,\dots,t_6}{\Re^2}{\Re^2} \)
    terdapat skalar \( c_1,\dots,c_6\in\Re \) sedemikian sehingga
    tidak semua $c_i$ sama dengan~$0$, tetapi
    \( c_1t_1+\dots+c_6t_6 \) adalah peta nol.
    (\textit{Petunjuk:}~angka enam agak menyesatkan.)
    \begin{answer}
      Tetapkan basis-basis dan representasikan transformasi-transformasi tersebut dengan
      matriks \( \nbyn{2} \).
      Ruang matriks \( \matspace_{\nbyn{2}} \) berdimensi empat,
      sehingga setiap himpunan beranggotakan enam matriks bergantung linear.
      Berdasarkan latihan sebelumnya, hal ini meluas menjadi kebergantungan peta.
      (Bagian yang menyesatkan hanyalah bahwa ada enam transformasi, bukan
      lima, sehingga jumlahnya lebih banyak daripada yang diperlukan untuk menunjukkan adanya
      kebergantungan.)
    \end{answer}
  \item 
    \definend{Teras}\index{teras}\index{matriks!teras}
    matriks persegi adalah jumlah entri-entri pada diagonal utama (entri
    \( 1,1 \)
    ditambah entri \( 2,2 \), dan seterusnya;
    kita akan melihat arti penting teras dalam Bab Lima).
    Tunjukkan bahwa \( \mbox{teras}(H+G)=\mbox{teras}(H)+\mbox{teras}(G)  \).
    Adakah hasil serupa untuk perkalian skalar?
    \begin{answer}
      Teras suatu jumlah sama dengan jumlah teras-terasnya
      karena baik \( \text{teras}(H+G) \) maupun
      \( \text{teras}(H)+\text{teras}(G) \) adalah jumlah dari
      \( h_{1,1}+g_{1,1} \), \( h_{2,2}+g_{2,2} \), dan seterusnya.
      Untuk perkalian skalar kita memperoleh
      \( \mbox{teras}(r\cdot H)=r\cdot\mbox{teras}(H) \); buktinya mudah.
      Jadi, peta teras adalah homomorfisme dari $\matspace_{\nbyn{n}}$ ke
      $\Re$.
    \end{answer}
  \item 
    Ingat kembali bahwa \definend{transpos}\index{matriks!transpos}%
    \index{transpos!interaksi dengan jumlah dan perkalian skalar}
    dari matriks $M$ adalah matriks lain yang entri $i,j$-nya merupakan
    entri $j,i$ dari $M$.
    Verifikasilah identitas-identitas berikut.
    \begin{exparts}
      \partsitem \( \trans{(G+H)}=\trans{G}+\trans{H} \)
      \partsitem \( \trans{(r\cdot H)}=r\cdot\trans{H} \)
    \end{exparts}
    \begin{answer} 
      \begin{exparts}
        \partsitem Entri \( i,j \) dari \( \trans{(G+H)} \) adalah
          \( g_{j,i}+h_{j,i} \).
          Ini juga merupakan entri \( i,j \) dari \( \trans{G}+\trans{H} \).
        \partsitem Entri \( i,j \) dari \( \trans{(r\cdot H)} \) adalah
          \( rh_{j,i} \),
          yang juga merupakan entri \( i,j \) dari \( r\cdot\trans{H} \).
      \end{exparts}  
   \end{answer}
  \recommended \item 
    Suatu matriks persegi disebut \definend{simetris}\index{matriks!simetris}%
    \index{matriks simetris} jika setiap entri \( i,j \) sama dengan
    entri \( j,i \), yakni jika matriks tersebut sama dengan transposnya.
    \begin{exparts}
      \partsitem Buktikan bahwa untuk sebarang matriks persegi~$H$,
        matriks \( H+\trans{H} \) bersifat simetris.
        Apakah setiap matriks simetris memiliki bentuk ini?
      \partsitem Buktikan bahwa himpunan matriks simetris \( \nbyn{n} \) merupakan
        subruang dari \( \matspace_{\nbyn{n}} \).
    \end{exparts}
    \begin{answer}  
      \begin{exparts}
        \partsitem Untuk \( H+\trans{H} \), entri \( i,j \)
          adalah \( h_{i,j}+h_{j,i} \), sedangkan
          entri \( j,i \) adalah \( h_{j,i}+h_{i,j} \).
          Keduanya sama, sehingga \( H+\trans{H} \) bersifat simetris.

          Setiap matriks simetris memang memiliki bentuk itu karena kita dapat menuliskan
          \( H=(1/2)\cdot(H+\trans{H}) \).
        \partsitem Himpunan matriks simetris tidak kosong karena
          memuat matriks nol.
          Jelas bahwa kelipatan skalar suatu matriks simetris juga simetris.
          Jumlah \( H+G \) dari dua matriks simetris
          bersifat simetris karena \( h_{i,j}+g_{i,j}=h_{j,i}+g_{j,i} \) (sebab
          \( h_{i,j}=h_{j,i} \) dan \( g_{i,j}=g_{j,i} \)).
          Jadi, subhimpunan tersebut tidak kosong dan tertutup terhadap operasi yang diwariskan,
          sehingga merupakan subruang.
      \end{exparts} 
    \end{answer}
  \recommended \item 
    \begin{exparts}
      \partsitem Bagaimana rank matriks berinteraksi dengan
        perkalian skalar\Dash dapatkah
        kelipatan skalar suatu matriks ber-rank~\( n \) memiliki rank kurang dari \( n \)?
        Lebih besar?
      \partsitem Bagaimana rank matriks berinteraksi dengan
        penjumlahan matriks\Dash dapatkah jumlah
        matriks-matriks ber-rank~\( n \) memiliki rank kurang dari \( n \)?
        Lebih besar?
    \end{exparts}
    \begin{answer}  
      \begin{exparts}
        \partsitem Perkalian skalar tidak mengubah rank matriks,
          kecuali bahwa perkalian dengan nol menghasilkan matriks
          ber-rank nol.
          (Hal ini mengikuti teorema pertama buku ini: mengalikan suatu
          baris dengan skalar tak nol
          tidak mengubah himpunan solusi sistem linear yang terkait.)
        \partsitem Jumlah matriks-matriks ber-rank \( n \) dapat memiliki rank
          kurang dari \( n \).
          Misalnya,
          untuk sebarang matriks \( H \), jumlah \( H+(-1)\cdot H \) ber-rank nol.

          Jumlah matriks-matriks ber-rank \( n \) juga dapat memiliki rank lebih besar dari \( n \).
          Berikut ini dua matriks ber-rank~satu yang jumlahnya merupakan matriks ber-rank~dua.
          \begin{equation*}
            \begin{mat}[r]
              1  &0  \\
              0  &0
            \end{mat}
            +\begin{mat}[r]
              0  &0  \\
              0  &1
            \end{mat}
            =\begin{mat}[r]
              1  &0  \\
              0  &1
            \end{mat}
          \end{equation*}
      \end{exparts}  
    \end{answer}
\end{exercises}











\subsection{Perkalian Matriks}
Setelah merepresentasikan penjumlahan dan perkalian skalar peta linear
dalam subbagian sebelumnya, operasi berikutnya yang wajar untuk ditinjau
adalah komposisi fungsi.

\begin{lemma} \label{lm:CompositionOfLinearMapsIsLinear}
%<*lm:CompositionOfLinearMapsIsLinear>
Komposisi peta-peta linear bersifat linear.
%</lm:CompositionOfLinearMapsIsLinear>
\end{lemma}

\begin{proof}
(\textit{Catatan:} argumen ini telah muncul sebelumnya,
sebagai bagian dari pembuktian Teorema~I.\ref{th:IsoEquivRel}.)
%<*pf:CompositionOfLinearMapsIsLinear>
Misalkan $\map{h}{V}{W}$ dan $\map{g}{W}{U}$ linear.
Perhitungan
\begin{multline*}
  \composed{g}{h}\,\bigl(c_1\cdot\vec{v}_1+c_2\cdot\vec{v}_2\bigr)
     =g\bigl(\,h(c_1\cdot \vec{v}_1+c_2\cdot\vec{v}_2)\,\bigr)        
     =g\bigl(\,c_1\cdot h(\vec{v}_1)+c_2\cdot h(\vec{v}_2)\,\bigr) \\ 
     =c_1\cdot g\bigl(h(\vec{v}_1))+c_2\cdot g(h(\vec{v}_2)\bigr)  
     =c_1\cdot (\composed{g}{h})(\vec{v}_1)
                   +c_2\cdot (\composed{g}{h})(\vec{v}_2)
\end{multline*}
menunjukkan bahwa $\map{\composed{g}{h}}{V}{U}$
mempertahankan kombinasi linear, sehingga bersifat linear.
%</pf:CompositionOfLinearMapsIsLinear>
\end{proof}

Seperti yang kita lakukan pada operasi penjumlahan matriks dan perkalian skalar,
kita akan melihat bagaimana representasi peta
komposisi berkaitan dengan representasi
peta-peta penyusunnya dengan terlebih dahulu meninjau sebuah contoh.

\begin{example} \label{ex:RepCompLinMaps}
Misalkan $\map{h}{\Re^4}{\Re^2}$ dan $\map{g}{\Re^2}{\Re^3}$,
tetapkan basis
\( B\subset\Re^4 \), \( C\subset\Re^2 \), \( D\subset\Re^3 \),
dan misalkan berikut ini representasinya.
\begin{equation*}
  H=\rep{h}{B,C}
  =\begin{mat}[r]
       4  &6  &8  &2  \\
       5  &7  &9  &3
    \end{mat}_{B,C}
  \quad
  G=\rep{g}{C,D}
  =\begin{mat}[r]
       1  &1  \\
       0  &1  \\
       1  &0
    \end{mat}_{C,D}
\end{equation*}
Untuk
merepresentasikan komposisi \( \map{\composed{g}{h}}{\Re^4}{\Re^3} \),
kita mulai dengan suatu $\vec{v}$, merepresentasikan $h$ dari $\vec{v}$,
lalu merepresentasikan $g$ dari hasil tersebut.
Representasi $h(\vec{v})$ adalah hasil kali matriks $h$ dan
vektor $\vec{v}$.
\begin{equation*}
  \rep{\,h(\vec{v})\,}{C}
  =
  \begin{mat}[r]
    4  &6  &8  &2  \\
    5  &7  &9  &3
  \end{mat}_{B,C}
  \colvec{v_1 \\ v_2 \\ v_3 \\ v_4}_B  
  =
  \colvec{4v_1+6v_2+8v_3+2v_4 \\ 5v_1+7v_2+9v_3+3v_4}_C
\end{equation*}
Representasi $g(\,h(\vec{v})\,)$ adalah hasil kali matriks $g$
dan vektor $h(\vec{v})$.
\begin{multline*}
  \rep{\,g(h(\vec{v}))\,}{D}\!
  =\begin{mat}[r]
      1  &1  \\
      0  &1  \\
      1  &0
   \end{mat}_{C,D}
   \colvec{4v_1+6v_2+8v_3+2v_4 \\ 5v_1+7v_2+9v_3+3v_4}_C    \\
  =\colvec{1\cdot(4v_1+6v_2+8v_3+2v_4)+1\cdot(5v_1+7v_2+9v_3+3v_4) \\ 
               0\cdot(4v_1+6v_2+8v_3+2v_4)+1\cdot(5v_1+7v_2+9v_3+3v_4) \\
               1\cdot(4v_1+6v_2+8v_3+2v_4)+0\cdot(5v_1+7v_2+9v_3+3v_4)}_D
\end{multline*}
Dengan mendistribusikan dan mengelompokkan ulang suku-suku menurut $v$, diperoleh
\begin{equation*}
  =
  \colvec{(1\cdot4+1\cdot5)v_1+
                (1\cdot6+1\cdot7)v_2+
                (1\cdot8+1\cdot9)v_3+
                (1\cdot2+1\cdot3)v_4 \\
                (0\cdot4+1\cdot5)v_1+
                (0\cdot6+1\cdot7)v_2+
                (0\cdot8+1\cdot9)v_3+
                (0\cdot2+1\cdot3)v_4 \\
                (1\cdot4+0\cdot5)v_1+
                (1\cdot6+0\cdot7)v_2+
                (1\cdot8+0\cdot9)v_3+
                (1\cdot2+0\cdot3)v_4}_D
\end{equation*}
yang merupakan hasil kali matriks-vektor berikut.
\begin{equation*}
  =
  \begin{mat}
    1\cdot4+1\cdot5  &1\cdot6+1\cdot7  &1\cdot8+1\cdot9  &1\cdot2+1\cdot3 \\
    0\cdot4+1\cdot5  &0\cdot6+1\cdot7  &0\cdot8+1\cdot9  &0\cdot2+1\cdot3 \\
    1\cdot4+0\cdot5  &1\cdot6+0\cdot7  &1\cdot8+0\cdot9  &1\cdot2+0\cdot3
  \end{mat}_{B,D}
  \colvec{v_1 \\ v_2 \\ v_3 \\ v_4}_B
\end{equation*}
Matriks yang merepresentasikan $\composed{g}{h}$ menggabungkan baris-baris $G$
dengan kolom-kolom $H$.
\end{example}

\begin{definition}  \label{def:MatMult}
%<*df:MatMult>
\definend{Hasil kali matriks}%
\index{matriks!perkalian}\index{perkalian!matriks-matriks}
dari matriks \( \nbym{m}{r} \)~\( G \)
dan matriks \( \nbym{r}{n} \)~\( H \) adalah
matriks \( \nbym{m}{n} \)~\( P \), dengan
\begin{equation*}
  p_{i,j}
  =
  g_{i,1}h_{1,j}+g_{i,2}h_{2,j}+\dots+g_{i,r}h_{r,j}
\end{equation*}
sehingga
entri ke-\( i,j \) dari hasil kali tersebut adalah hasil kali titik antara
baris ke-\( i \) matriks pertama dan kolom ke-\( j \) matriks kedua.
\begin{equation*}
    GH=
    \begin{mat}
              &\vdots                     \\
      g_{i,1} &g_{i,2} &\cdots  &g_{i,r}  \\
              &\vdots
    \end{mat}
    \begin{mat}
              &h_{1,j}           \\
      \cdots  &h_{2,j} &\cdots   \\
              &\vdots            \\
              &h_{r,j}
    \end{mat}
  =
    \begin{mat}
              &\vdots            \\
      \cdots  &p_{i,j}  &\cdots  \\
              &\vdots
    \end{mat}
\end{equation*}
%</df:MatMult>
\end{definition}

\begin{example}
\begin{equation*}
  \begin{mat}[r]
     2  &0  \\
     4  &6  \\
     8  &2
  \end{mat}
  \begin{mat}[r]
    1  &3  \\
    5  &7
  \end{mat}
  =
  \begin{mat}
   2\cdot 1+0\cdot 5  &2\cdot 3+0\cdot 7  \\
   4\cdot 1+6\cdot 5  &4\cdot 3+6\cdot 7  \\
   8\cdot 1+2\cdot 5  &8\cdot 3+2\cdot 7
  \end{mat}
  =
  \begin{mat}[r]
    2  &6  \\
   34  &54 \\
   18  &38
  \end{mat}
\end{equation*}
\end{example}


\begin{example}  \label{ex:PairTwoByTwoMult}
Beberapa hasil kali tidak terdefinisi, misalnya hasil kali matriks
$\nbym{2}{3}$ dengan matriks $\nbyn{2}$, karena banyaknya kolom
dalam matriks pertama harus sama dengan banyaknya baris dalam matriks kedua.
Namun, hasil kali dua matriks $\nbyn{n}$ selalu terdefinisi.
Berikut ini dua matriks $\nbyn{2}$.
\begin{equation*}
  \begin{mat}[r]
     1  &2  \\
     3  &4  
  \end{mat}
  \begin{mat}[r]
    -1  &0  \\
    2   &-2
  \end{mat}
  =
  \begin{mat}
   1\cdot (-1)+2\cdot 2  &1\cdot 0+2\cdot (-2)  \\
   3\cdot (-1)+4\cdot 2  &3\cdot 0+4\cdot (-2)  
  \end{mat}
  =
  \begin{mat}[r]
    3  &-4  \\
    5  &-8 
  \end{mat}
\end{equation*}
\end{example}


\begin{example}
Matriks-matriks dari \nearbyexample{ex:RepCompLinMaps} digabungkan sebagai berikut.
\begin{multline*}
  \begin{mat}[r]
       1  &1  \\
       0  &1  \\
       1  &0
    \end{mat}               
  \begin{mat}[r]
       4  &6  &8  &2  \\
       5  &7  &9  &3
    \end{mat}  \\
   \begin{aligned}
     &=\begin{mat}
      1\cdot4+1\cdot5  &1\cdot6+1\cdot7  &1\cdot8+1\cdot9  &1\cdot2+1\cdot3 \\
      0\cdot4+1\cdot5  &0\cdot6+1\cdot7  &0\cdot8+1\cdot9  &0\cdot2+1\cdot3 \\
      1\cdot4+0\cdot5  &1\cdot6+0\cdot7  &1\cdot8+0\cdot9  &1\cdot2+0\cdot3
      \end{mat}               \\
     &=\begin{mat}[r]
        9  &13  &17  &5  \\
        5  &7   &9   &3  \\
        4  &6   &8   &2
      \end{mat}
    \end{aligned}
\end{multline*}
\end{example}

% We next verify that our
% definition of the matrix-matrix multiplication operation
% does what we intend.

\begin{theorem}  \label{th:MatMultRepComp}
%<*th:MatMultRepComp>
\index{fungsi!komposisi}\index{homomorfisme!komposisi}
Komposisi peta-peta linear direpresentasikan oleh hasil kali matriks
dari representasi peta-peta tersebut.
%</th:MatMultRepComp>
\end{theorem}

\begin{proof}
Argumen ini menggeneralisasi \nearbyexample{ex:RepCompLinMaps}.
%<*pf:MatMultRepComp0>
Misalkan \( \map{h}{V}{W} \) dan \( \map{g}{W}{X} \) direpresentasikan oleh
\( H \) dan \( G \) terhadap basis
\( B\subset V \), \( C\subset W \), dan \( D\subset X \), yang masing-masing berukuran
\( n \), \( r \), dan \( m \).
Untuk sebarang \( \vec{v}\in V \), komponen ke-\( k \) dari
\( \rep{\,h(\vec{v})\,}{C} \) adalah
\begin{equation*}
  h_{k,1}v_1+\cdots+h_{k,n}v_n
\end{equation*}
sehingga komponen ke-\( i \) dari
\( \rep{\,\composed{g}{h}\,(\vec{v})\,}{D} \) adalah sebagai berikut.
\begin{multline*}
  g_{i,1}\cdot(h_{1,1}v_1+\dots+h_{1,n}v_n)
  +g_{i,2}\cdot(h_{2,1}v_1+\dots+h_{2,n}v_n)  \\
  +\dots                                 
  +g_{i,r}\cdot(h_{r,1}v_1+\dots+h_{r,n}v_n)
\end{multline*}
Distribusikan dan kelompokkan ulang suku-suku menurut \( v \).
\begin{multline*}
  =(g_{i,1} h_{1,1}+g_{i,2} h_{2,1}+\dots+g_{i,r}h_{r,1})\cdot v_1    \\
   +\dots
   +(g_{i,1} h_{1,n}+g_{i,2} h_{2,n}
           +\dots+g_{i,r} h_{r,n})\cdot v_n
\end{multline*}
%</pf:MatMultRepComp0>
%<*pf:MatMultRepComp1>
Akhiri pembuktian dengan mengamati bahwa koefisien setiap \( v_j \),
\begin{equation*}
  g_{i,1}h_{1,j}+g_{i,2}h_{2,j}+\dots+g_{i,r}h_{r,j}
\end{equation*}
sesuai dengan definisi entri \( i,j \) dari hasil kali \( GH \).
%</pf:MatMultRepComp1>
\end{proof}

%<*MatMultArrowDiag0>
\definend{Diagram panah}\index{diagram panah} ini
menggambarkan hubungan antara peta dan matriks
(`thd' menyingkat `terhadap').
%</MatMultArrowDiag0>
%\medskip
\begin{center}
  \includegraphics{map/mp/ch3.20}
  %\quad
  %\includegraphics{map/mp/ch3.19}  
\end{center}
%\medskip
%<*MatMultArrowDiag1>
Di atas anak panah, peta-petanya menunjukkan bahwa dua cara berpindah dari
\( V \) ke \( X \),
langsung melalui komposisi atau dalam dua langkah melalui \( W \),
memberikan hasil yang sama,
\begin{equation*}
  \vec{v}\mapsunder{g\circ h}g(h(\vec{v}))
  \qquad
  \vec{v}\mapsunder{h}h(\vec{v})\mapsunder{g}g(h(\vec{v}))
\end{equation*}
(ini hanyalah definisi komposisi).
Di bawah anak panah, matriks-matriksnya menunjukkan bahwa mengalikan
$GH$ dengan vektor kolom $\rep{\vec{v}}{B}$ memberikan
hasil yang sama dengan mula-mula mengalikan vektor kolom itu dengan $H$,
lalu mengalikan hasilnya dengan $G$.
\begin{equation*}
   \rep{\composed{g}{h}}{B,D}
   =GH
   \qquad
   \rep{g}{C,D}\;\rep{h}{B,C}=GH
\end{equation*}
%</MatMultArrowDiag1>

Seperti disebutkan dalam \nearbyexample{ex:PairTwoByTwoMult},
karena banyaknya kolom di sebelah kiri
tidak sama dengan banyaknya baris di sebelah kanan,
hasil kali matriks $\nbym{2}{3}$ dengan matriks
$\nbyn{2}$ seperti berikut tidak terdefinisi.
\begin{equation*}
     \begin{mat}[r]
       -1  &2  &0  \\
        0  &10 &1.1
     \end{mat}
     \begin{mat}[r]
        0  &0  \\
        0  &2
     \end{mat}
\end{equation*}
Definisi tersebut mengharuskan ukurannya bersesuaian
karena kita menginginkan agar
komposisi fungsi yang mendasarinya dapat dilakukan.
\begin{equation*}
  \text{ruang berdimensi \( n \)}
  \;\stackrel{h}{\longrightarrow}\;
  \text{ruang berdimensi \( r \)}
  \;\stackrel{g}{\longrightarrow}\;
  \text{ruang berdimensi \( m \)}
  \tag{$*$}
\end{equation*}
Jadi, perkalian matriks menggabungkan matriks $\nbym{m}{r}$~$G$
dengan matriks $\nbym{r}{n}$~$F$ untuk menghasilkan
$GF$ berukuran $\nbym{m}{n}$.
Singkatnya:
$\nbym{m}{r}\text{\ dikali\ }\nbym{r}{n}\text{\ sama dengan\ }\nbym{m}{n}$.

\begin{remark}
Urutan dimensi dapat membingungkan.
Dalam `$\nbym{m}{r}\text{\ dikali\ }\nbym{r}{n}\text{\ sama dengan\ }\nbym{m}{n}$',
bilangan yang dituliskan pertama adalah~$m$.
Namun,~$m$ muncul terakhir dalam
baris keterangan dimensi peta~($*$) di atas, dan dimensi-dimensi lainnya
juga muncul dalam urutan terbalik.
Penjelasannya adalah bahwa meskipun $h$ dikerjakan terlebih dahulu,
kemudian diikuti oleh $g$,
kita menuliskan komposisinya sebagai $\composed{g}{h}$, dengan $g$ di sebelah kiri
(berasal dari notasi $g(h(\vec{v}))$).
% (Some people try to lessen confusion by reading 
% `$\composed{g}{h}$' aloud as ``$g$ following $h$.'')
Hal ini terbawa ke matriks, sehingga $\composed{g}{h}$
direpresentasikan oleh $GH$.
\end{remark}

Kita dapat memperoleh wawasan tentang operasi perkalian matriks-matriks
dengan mempelajari bagaimana entri-entrinya digabungkan.
Misalnya, cara lain untuk memahami mengapa kita mengharuskan ukuran-ukuran di atas
bersesuaian
adalah bahwa baris matriks sebelah kiri harus memiliki banyak entri yang sama
dengan kolom matriks sebelah kanan; jika tidak, ada entri
yang tidak memiliki pasangan dari matriks lainnya.

Aspek lain dari kombinatorika
perkalian matriks
dalam jumlah yang mendefinisikan entri $i,j$
ditonjolkan di sini dengan memberi kotak pada subskrip-subskrip yang sama.
\begin{equation*} % \setlength{\fboxsep}{.15em}
  p_{i,j}
  =
  g_{i,\text{\highlight{$1 $}}}h_{\text{\highlight{$1 $}},j}
   +g_{i,\text{\highlight{$2 $}}}h_{\text{\highlight{$2 $}},j}
   +\dots+g_{i,\text{\highlight{$r $}}}h_{\text{\highlight{$r $}},j}
\end{equation*}
Subskrip yang disorot pada $g$ adalah indeks kolom, sedangkan yang pada
$h$ adalah indeks baris.
Artinya, penjumlahan berlangsung pada kolom-kolom $G$
tetapi pada baris-baris $H$\Dash
definisi tersebut memperlakukan sisi kiri dan kanan secara berbeda.
Karena itu, kita patut menduga bahwa \( GH \)
dapat tidak sama dengan \( HG \).

\begin{example}   \label{ex:MatMultNoCommute}
Perkalian matriks tidak bersifat komutatif.
\begin{equation*}
    \begin{mat}[r]
      1  &2  \\
      3  &4
    \end{mat}
    \begin{mat}[r]
      5  &6  \\
      7  &8
    \end{mat}
  =
    \begin{mat}[r]
      19  &22  \\
      43  &50
    \end{mat}
  \qquad
    \begin{mat}[r]
      5  &6  \\
      7  &8
    \end{mat}
    \begin{mat}[r]
      1  &2  \\
      3  &4
    \end{mat}
  =
    \begin{mat}[r]
      23  &34  \\
      31  &46
    \end{mat}
\end{equation*}
% In fact, matrix multiplication hardly ever commutes, in that if we use
% multiply randomly chosen matrices both ways.
\end{example}

\begin{example}  \label{ex:MatMultNoCommuteII}
Sifat komutatif bahkan dapat gagal secara lebih mencolok:
\begin{equation*}
    \begin{mat}[r]
      5  &6  \\
      7  &8
    \end{mat}
    \begin{mat}[r]
      1  &2  &0 \\
      3  &4  &0
    \end{mat}
  =
    \begin{mat}[r]
      23  &34  &0  \\
      31  &46  &0
    \end{mat}
\end{equation*}
sedangkan
\begin{equation*}
    \begin{mat}[r]
      1  &2  &0 \\
      3  &4  &0
    \end{mat}
    \begin{mat}[r]
      5  &6  \\
      7  &8
    \end{mat}
\end{equation*}
bahkan tidak terdefinisi.
\end{example}

\begin{remark}
Fakta bahwa
perkalian matriks tidak bersifat komutatif mungkin terasa ganjil pada awalnya,
barangkali karena kebanyakan
operasi matematika dalam mata kuliah sebelumnya bersifat komutatif.
Namun, perkalian matriks merepresentasikan komposisi fungsi,
dan komposisi fungsi
tidak bersifat komutatif: jika
\( f(x)=2x \) dan \( g(x)=x+1 \), maka \( \composed{g}{f}(x)=2x+1 \)
sedangkan \( \composed{f}{g}(x)=2(x+1)=2x+2 \).
% (True, this \( g \) is not linear and we might have hoped that linear functions
% would commute but this shows that the failure of commutativity for
% matrix multiplication fits into a larger context.)
\end{remark}

Selain ketidakkomutatifannya,
perkalian matriks berperilaku baik secara aljabar.
Hasil berikut memberikan beberapa sifat yang baik; sifat lainnya terdapat dalam
\nearbyexercise{exer:NicePropsMatMult} dan
\nearbyexercise{exer:MoreNicePropsMatMult}.

\begin{theorem}  \label{th:MatMultWellBehaved}
%<*th:MatMultWellBehaved>
Jika \( F \), \( G \), dan \( H \) adalah matriks,
dan hasil kali matriksnya terdefinisi,
maka perkalian tersebut bersifat asosiatif,
$(FG)H=F(GH)$
serta distributif terhadap penjumlahan matriks,
$F(G+H)=FG+FH$ dan $(G+H)F=GF+HF$.
%</th:MatMultWellBehaved>
\end{theorem}

\begin{proof}
%<*pf:MatMultWellBehaved0>
Sifat asosiatif berlaku karena perkalian matriks merepresentasikan komposisi
fungsi, yang bersifat asosiatif:~peta
\( \composed{(\composed{f}{g})}{h} \) dan
\( \composed{f}{(\composed{g}{h})} \) sama
karena keduanya memetakan \( \vec{v} \) ke \( f(g(h(\vec{v}))) \).
%</pf:MatMultWellBehaved0>

%<*pf:MatMultWellBehaved1>
Sifat distributif serupa.
Misalnya, sifat pertama memberikan
\( \composed{f}{(g+h)}\,(\vec{v})
    =f\bigl(\,(g+h)(\vec{v})\,\bigr)  
    =f\bigl(\,g(\vec{v})+h(\vec{v})\,\bigr)
    =f(g(\vec{v}))+f(h(\vec{v}))
    =\composed{f}{g}(\vec{v})+\composed{f}{h}(\vec{v}) \) 
(kesamaan ketiga menggunakan linearitas $f$).
Distributivitas kanan dibuktikan dengan cara yang sama.
%</pf:MatMultWellBehaved1>
\end{proof}

\begin{remark} 
Sebagai alternatif, kita dapat membuktikan hasil tersebut dengan bersusah payah menelusuri indeks.
Sebagai contoh, untuk sifat asosiatif,
entri \( i,j \) dari \( (FG)H \) adalah
\begin{align*}
   \MoveEqLeft (f_{i,1}g_{1,1}+f_{i,2}g_{2,1}+\dots+f_{i,r}g_{r,1})h_{1,j}         \\
   &+(f_{i,1}g_{1,2}+f_{i,2}g_{2,2}+\dots+f_{i,r}g_{r,2})h_{2,j}   \\
   &\vdotswithin{+}                                       \\
   &+(f_{i,1}g_{1,s}+f_{i,2}g_{2,s}+\dots+f_{i,r}g_{r,s})h_{s,j}
\end{align*}
dengan \( F \), \( G \), dan \( H \) masing-masing merupakan matriks \( \nbym{m}{r} \),
\( \nbym{r}{s} \), dan \( \nbym{s}{n} \).
Distribusikan menjadi
\begin{align*}
   \MoveEqLeft f_{i,1}g_{1,1}h_{1,j}+f_{i,2}g_{2,1}h_{1,j}+
            \dots+f_{i,r}g_{r,1}h_{1,j}                  \\ 
   &+f_{i,1}g_{1,2}h_{2,j}+f_{i,2}g_{2,2}h_{2,j}+
          \dots+f_{i,r}g_{r,2}h_{2,j}                      \\
   &\vdotswithin{+}                                           \\
   &+f_{i,1}g_{1,s}h_{s,j}+f_{i,2}g_{2,s}h_{s,j}+\dots
        +f_{i,r}g_{r,s}h_{s,j}
\end{align*}
dan kelompokkan ulang menurut \( f \)
\begin{align*}
   \MoveEqLeft f_{i,1}(g_{1,1}h_{1,j}+g_{1,2}h_{2,j}+\dots+g_{1,s}h_{s,j})     \\
   &+f_{i,2}(g_{2,1}h_{1,j}+g_{2,2}h_{2,j}+\dots+g_{2,s}h_{s,j})   \\
   &\vdotswithin{+}                                       \\
   &+f_{i,r}(g_{r,1}h_{1,j}+g_{r,2}h_{2,j}+\dots+g_{r,s}h_{s,j})
\end{align*}
untuk memperoleh entri \( i,j \) dari \( F(GH) \).

Bandingkan kedua pembuktian tersebut.
Argumen yang sarat indeks sulit dipahami karena meskipun
perhitungannya mudah diperiksa, aritmetikanya tampak tidak terhubung dengan
gagasan apa pun.
Argumen dalam pembuktian
lebih singkat dan juga menjelaskan mengapa sifat ini ``sebenarnya'' berlaku.
Hal ini menggambarkan komentar yang dibuat pada awal
bab tentang ruang vektor\Dash setidaknya kadang-kadang,
argumen yang menggunakan konstruksi tingkat tinggi lebih jelas.
\end{remark}

Sekarang kita telah melihat cara
merepresentasikan komposisi peta linear.
Subbagian berikutnya akan melanjutkan
penelaahan operasi ini.

\begin{exercises}
  \recommended \item 
    Hitunglah, atau nyatakan ``tidak terdefinisi.''
    \begin{exparts*}
      \partsitem
        $\begin{mat}[r]
          3  &1  \\
          -4 &2 
        \end{mat}
        \begin{mat}[r]
          0  &5  \\
          0  &0.5
        \end{mat}$
      \partsitem \(
        \begin{mat}[r]
          1  &1  &-1  \\
          4  &0  &3
        \end{mat}
        \begin{mat}[r]
          2  &-1 &-1  \\
          3  &1  &1   \\
          3  &1  &1
        \end{mat}   \)
      \partsitem \(
        \begin{mat}[r]
          2  &-7 \\
          7  &4
        \end{mat}
        \begin{mat}[r]
          1  &0  &5   \\
         -1  &1  &1   \\
          3  &8  &4
        \end{mat}   \)
      \partsitem \(
        \begin{mat}[r]
          5  &2  \\
          3  &1
        \end{mat}
        \begin{mat}[r]
          -1  &2   \\
           3  &-5
        \end{mat} \)
    \end{exparts*}
    \begin{answer} 
      \begin{exparts*}
        \partsitem
          $\begin{mat}[r]
            0  &15.5  \\
            0  &-19
          \end{mat}$
        \partsitem \(
          \begin{mat}[r]
            2  &-1  &-1  \\
           17  &-1  &-1
          \end{mat} \)
        \partsitem Tidak terdefinisi.
        \partsitem \( \begin{mat}[r]
                   1  &0  \\
                   0  &1
                 \end{mat}  \)
      \end{exparts*}  
    \end{answer}
  \recommended \item  
    Dengan
    \begin{equation*}
      A=
       \begin{mat}[r]
         1  &-1  \\
         2  &0   \\
       \end{mat}
      \quad
      B=
       \begin{mat}[r]
         5  &2   \\
         4  &4   \\
       \end{mat}
      \quad
      C=
       \begin{mat}[r]
        -2  &3   \\
        -4  &1   \\
       \end{mat}
    \end{equation*}
    hitunglah, atau nyatakan ``tidak terdefinisi.''
    \begin{exparts*}
      \partsitem \( AB \)
      \partsitem \( (AB)C \)
      \partsitem \( BC \)
      \partsitem \( A(BC) \)
    \end{exparts*}
    \begin{answer}  
      \begin{exparts*}
        \partsitem \(
          \begin{mat}[r]
            1  &-2   \\
            10 &4
          \end{mat}  \)
        \partsitem \( 
          \begin{mat}[r]
            1  &-2   \\
            10 &4
          \end{mat}
          \begin{mat}[r]
            -2 &3    \\
            -4 &1
          \end{mat}=
          \begin{mat}[r]
            6  &1    \\
           -36 &34
          \end{mat}  \)
        \partsitem \(
          \begin{mat}[r]
           -18 &17   \\
           -24 &16
          \end{mat}  \)
        \partsitem \(
          \begin{mat}[r]
            1  &-1  \\
            2  &0
          \end{mat}
          \begin{mat}[r]
           -18 &17   \\
           -24 &16
          \end{mat}=
          \begin{mat}[r]
            6  &1    \\
           -36 &34
          \end{mat}  \)
      \end{exparts*}  
    \end{answer}
  \item  
    Hasil kali mana yang terdefinisi?
    \begin{exparts*}
      \partsitem \( \nbym{3}{2} \)~dikali~\( \nbym{2}{3} \)
      \partsitem \( \nbym{2}{3} \)~dikali~\( \nbym{3}{2} \)
      \partsitem \( \nbym{2}{2} \)~dikali~\( \nbym{3}{3} \)
      \partsitem \( \nbym{3}{3} \)~dikali~\( \nbym{2}{2} \)
    \end{exparts*}
    \begin{answer}
      \begin{exparts*}
        \partsitem Ya.
        \partsitem Ya.
        \partsitem Tidak.
        \partsitem Tidak.
      \end{exparts*}
    \end{answer}
  \recommended \item 
    Berikan ukuran hasil kalinya, atau nyatakan ``tidak terdefinisi.''
    \begin{exparts}
      \partsitem matriks \( \nbym{2}{3} \) dikali matriks \( \nbym{3}{1} \)
      \partsitem matriks \( \nbym{1}{12} \) dikali matriks \( \nbym{12}{1} \)
      \partsitem matriks \( \nbym{2}{3} \) dikali matriks \( \nbym{2}{1} \)
      \partsitem matriks \( \nbym{2}{2} \) dikali matriks \( \nbym{2}{2} \)
    \end{exparts}
    \begin{answer}  
      \begin{exparts*}
        \partsitem \( \nbym{2}{1} \)
        \partsitem \( \nbym{1}{1} \)
        \partsitem Tidak terdefinisi.
        \partsitem \( \nbym{2}{2} \)
      \end{exparts*}  
     \end{answer}
  \recommended \item 
    Tentukan sistem persamaan yang dihasilkan dengan memulai dari
    \begin{equation*}
      \begin{linsys}{3}
        h_{1,1}x_1  &+  &h_{1,2}x_2  &+  &h_{1,3}x_3  &=  &d_1  \\
        h_{2,1}x_1  &+  &h_{2,2}x_2  &+  &h_{2,3}x_3  &=  &d_2  
       \end{linsys}
    \end{equation*}
    dan melakukan perubahan variabel (yakni substitusi) berikut.
    \begin{equation*}
      \begin{linsys}{2}
        x_1  &=  &g_{1,1}y_1  &+  &g_{1,2}y_2  \\
        x_2  &=  &g_{2,1}y_1  &+  &g_{2,2}y_2  \\
        x_3  &=  &g_{3,1}y_1  &+  &g_{3,2}y_2  
      \end{linsys}
    \end{equation*}
    \begin{answer}
      Kita memperoleh
      \begin{equation*}
        \begin{linsys}{3}
            h_{1,1}\cdot(g_{1,1}y_1+g_{1,2}y_2)  
               &+  &h_{1,2}\cdot (g_{2,1}y_1+g_{2,2}y_2)  
               &+  &h_{1,3}\cdot (g_{3,1}y_1+g_{3,2}y_2)  
               &=  &d_1  \\
            h_{2,1}\cdot(g_{1,1}y_1+g_{1,2}y_2)  
               &+  &h_{2,2}\cdot (g_{2,1}y_1+g_{2,2}y_2)  
               &+  &h_{2,3}\cdot (g_{3,1}y_1+g_{3,2}y_2)  
               &=  &d_2  
         \end{linsys}
      \end{equation*}
      yang, setelah dikembangkan dan dikelompokkan ulang menurut $y$, menghasilkan berikut ini.
      \begin{equation*}
        \begin{linsys}{2}
            (h_{1,1}g_{1,1}+h_{1,2}g_{2,1}+h_{1,3}g_{3,1})y_1
              &+  &(h_{1,1}g_{1,2}+h_{1,2}g_{2,2}+h_{1,3}g_{3,2})y_2  
               &=  &d_1  \\
            (h_{2,1}g_{1,1}+h_{2,2}g_{2,1}+h_{2,3}g_{3,1})y_1
              &+  &(h_{2,1}g_{1,2}+h_{2,2}g_{2,2}+h_{2,3}g_{3,2})y_2  
               &=  &d_2
         \end{linsys}
      \end{equation*}
      Kita dapat menyatakan sistem awal
      dan sistem yang digunakan untuk substitusi
dalam bahasa matriks sebagai
      \begin{equation*}
        \begin{mat}
          h_{1,1}  &h_{1,2}  &h_{1,3}  \\
          h_{2,1}  &h_{2,2}  &h_{2,3}
        \end{mat}
        \colvec{x_1 \\ x_2 \\ x_3}
        =H\colvec{x_1 \\ x_2 \\ x_3}=\colvec{d_1 \\ d_2}
      \end{equation*}
      dan
      \begin{equation*}
        \begin{mat}
          g_{1,1}  &g_{1,2}  \\
          g_{2,1}  &g_{2,2}  \\
          g_{3,1}  &g_{3,2}
        \end{mat}
        \colvec{y_1 \\ y_2}
        =G\colvec{y_1 \\ y_2}=\colvec{x_1 \\ x_2 \\ x_3}
      \end{equation*}
      dan dengan ini, substitusinya adalah
      $\vec{d}=H\vec{x}=H(G\vec{y})=(HG)\vec{y}$.
    \end{answer}
  \recommended \item   Perhatikan dua fungsi linear
    $\map{h}{\Re^3}{\polyspace_2}$
    dan
    $\map{g}{\polyspace_2}{\matspace_{\nbyn{2}}}$
    yang diberikan sebagai berikut.
    \begin{equation*}
      \colvec{a \\ b \\ c}\mapsto (a+b)x^2+(2a+2b)x+c
      \qquad
      px^2+qx+r\mapsto
      \begin{mat}
        p &p-2q \\
        q &0
      \end{mat}
    \end{equation*}
    Gunakan basis-basis berikut untuk ruang-ruang tersebut.
    \begin{gather*}
      B=\sequence{\colvec{1 \\ 1 \\ 1}, 
                  \colvec{0 \\ 1 \\ 1}, 
                  \colvec{0 \\ 0 \\ 1}}
     \qquad
     C=\sequence{1+x,1-x,x^2}
     \\
     D=\sequence{
       \begin{mat}
         1 &0 \\
         0 &0
       \end{mat},
       \begin{mat}
         0 &2 \\
         0 &0
       \end{mat},
       \begin{mat}
         0 &0 \\
         3 &0
       \end{mat},
       \begin{mat}
         0 &0 \\
         0 &4
       \end{mat}}
    \end{gather*}
    \begin{exparts}
      \partsitem Berikan rumus untuk peta komposisi
        $\map{\composed{g}{h}}{\Re^3}{\matspace_{\nbyn{2}}}$ yang diturunkan
        langsung dari definisi di atas.
      \partsitem Representasikan $h$ dan~$g$ terhadap basis-basis yang sesuai.
      \partsitem Representasikan peta $\composed{g}{h}$ yang dihitung pada bagian pertama
        terhadap basis-basis yang sesuai.
      \partsitem Periksa bahwa hasil kali kedua matriks dari bagian kedua
        adalah matriks dari bagian ketiga.
    \end{exparts}
    \begin{answer}
      \begin{exparts}
        \partsitem
          Dengan mengikuti definisi, diperoleh berikut ini.
          \begin{align*}
            \colvec{a \\ b \\ c}
            &\mapsto 
            (a+b)x^2+(2a+2b)x+c                   \\
            &\mapsto
            \begin{mat}
              a+b &(a+b)-2(2a+2b) \\
              2a+2b &0
            \end{mat}    
            =
            \begin{mat}
              a+b   &-3a-3b  \\
              2a+2b &0
            \end{mat}    
          \end{align*}
        \partsitem
          Karena
          \begin{equation*}
            \colvec{1 \\ 1 \\ 1}\mapsto 2x^2+4x+1
            \quad
            \colvec{0 \\ 1 \\ 1}\mapsto x^2+2x+1
            \quad
            \colvec{0 \\ 0 \\ 1}\mapsto 0x^2+0x+1   
          \end{equation*}
          kita memperoleh representasi berikut untuk~$h$.
          \begin{equation*}
            \rep{h}{B,C}
            =
            \begin{mat}
              5/2 &3/2  &1/2  \\ 
             -3/2 &-1/2 &1/2  \\
              2   &1    &0
            \end{mat}
          \end{equation*}
          Demikian pula, karena
          \begin{equation*}
            1+x\mapsto
            \begin{mat}
              0 &-2 \\
              1 &0
            \end{mat}
            \quad
            1-x\mapsto
            \begin{mat}
              0 &2 \\
              -1 &0
            \end{mat}
            \quad
            x^2\mapsto
            \begin{mat}
              1 &1 \\
              0 &0
            \end{mat}
          \end{equation*}
          berikut ini adalah representasi~$g$.
          \begin{equation*}
            \rep{g}{C,D}=
            \begin{mat}
              0  &0    &1  \\
             -1  &1    &1/2 \\
             1/3 &-1/3 &0   \\
              0  &0    &0
            \end{mat}
          \end{equation*}
        \partsitem
          Tindakan $\composed{g}{h}$ pada basis domain adalah sebagai berikut.
          \begin{equation*}
            \colvec{1 \\ 1 \\ 1}\mapsto 
              \begin{mat}
                2 &-6 \\
                4 &0
              \end{mat}
            \quad
            \colvec{0 \\ 1 \\ 1}\mapsto
              \begin{mat}
                1 &-3 \\
                2 &0
              \end{mat}
            \quad
            \colvec{0 \\ 0 \\ 1}\mapsto
              \begin{mat}
                0 &0 \\
                0 &0
              \end{mat}
          \end{equation*}
          Kita memperoleh berikut ini.
          \begin{equation*}
            \rep{\composed{g}{h}}{B,D}=
            \begin{mat}
               2   &1    &0  \\  
              -3   &-3/2 &0  \\
             4/3   &2/3  &0  \\
               0   &0    &0
            \end{mat}
          \end{equation*}
        \partsitem 
          Perkalian matriksnya rutin; perhatikan saja urutannya.
          \begin{equation*}
            \begin{mat}
              0  &0    &1  \\
             -1  &1    &1/2 \\
             1/3 &-1/3 &0   \\
              0  &0    &0
            \end{mat}
            \begin{mat}
              5/2 &3/2  &1/2  \\ 
             -3/2 &-1/2 &1/2  \\
              2   &1    &0
            \end{mat}
            =        
            \begin{mat}
               2   &1    &0  \\  
              -3   &-3/2 &0  \\
             4/3   &2/3  &0  \\
               0   &0    &0
           \end{mat}
          \end{equation*}
      \end{exparts}
    \end{answer}
  \item 
    Seperti ditunjukkan oleh \nearbydefinition{def:MatMult}, operasi perkalian
    matriks menggeneralisasi hasil kali titik.
    Apakah hasil kali titik antara vektor baris \( \nbym{1}{n} \)
    dan vektor kolom \( \nbym{n}{1} \) sama dengan
    hasil kali matriks keduanya?
    \begin{answer}
      Secara teknis, tidak.
      Operasi hasil kali titik menghasilkan suatu skalar, sedangkan hasil kali matriks
      menghasilkan matriks $\nbyn{1}$.
      Namun, biasanya kita akan mengabaikan perbedaan ini.
    \end{answer}
  \recommended \item 
    Representasikan peta turunan pada \( \polyspace_n \)
    terhadap \( B,B \), dengan \( B \) sebagai basis alami
    \( \sequence{1,x,\ldots,x^n} \).
    Tunjukkan bahwa hasil kali matriks ini dengan dirinya sendiri terdefinisi;
    peta apa yang direpresentasikannya?
    \begin{answer}
      Tindakan $d/dx$ pada $B$ adalah
      $1\mapsto0$, $x\mapsto 1$, $x^2\mapsto 2x$, \ldots sehingga
      berikut ini representasi matriks $\nbyn{(n+1)}$-nya.
      \begin{equation*}
        \rep{\frac{d}{dx}}{B,B}=
        \begin{mat}
          0  &1 &0  &  &0  \\
          0  &0 &2  &  &0  \\
             &  &   &\ddots\\
          0  &0 &0  &  &n  \\
          0  &0 &0  &  &0
        \end{mat}
      \end{equation*}
      Hasil kali matriks ini dengan dirinya sendiri terdefinisi karena matriks tersebut
      persegi.
      \begin{equation*}
        \begin{mat}
          0  &1 &0  &  &0  \\
          0  &0 &2  &  &0  \\
             &  &   &\ddots\\
          0  &0 &0  &  &n  \\
          0  &0 &0  &  &0
        \end{mat}^2
        =
        \begin{mat}
          0  &0 &2  &0  &        &0       \\
          0  &0 &0  &6  &        &0       \\
             &  &   &   &\ddots  &        \\
          0  &0 &0  &  &         &n(n-1)  \\
          0  &0 &0  &  &         &0       \\
          0  &0 &0  &  &         &0
        \end{mat}
      \end{equation*}
      Peta yang direpresentasikan adalah komposisi
      \begin{equation*}
        p\;\mapsunder{\frac{d}{dx}}\;\frac{d\,p}{dx}
         \;\mapsunder{\frac{d}{dx}}\;\frac{d^2\,p}{dx^2}
      \end{equation*}
      yang merupakan operasi turunan kedua.
     \end{answer}
  \item 
    \cite{Cleary}
    Pasangkan setiap jenis matriks dengan semua deskripsi berikut yang mungkin sesuai;
    nyatakan `Tidak ada' jika tidak satu pun berlaku:
    (i)~dapat dikalikan dengan transposnya untuk menghasilkan matriks $\nbyn{1}$,
    % (ii)~is similar to the $\nbyn{3}$ matrix of all zeros,
    (ii)~dapat merepresentasikan peta linear dari $\Re^3$ ke $\Re^2$ yang tidak pada,
    (iii)~dapat merepresentasikan isomorfisme dari $\Re^3$ ke $\polyspace^2$.
    \begin{exparts}
      \item matriks $\nbym{2}{3}$ ber-rank~$1$
      \item matriks $\nbyn{3}$ yang nonsingular
      \item matriks $\nbyn{2}$ yang singular
      \item vektor kolom $\nbym{n}{1}$
    \end{exparts}
    \begin{answer}
      \begin{exparts}
        \item ii
        \item iii
        \item Tidak ada
        \item Tidak ada, atau (i) jika kita menyertakan perkalian dari kiri.
      \end{exparts}
    \end{answer}
  \item  
    Tunjukkan bahwa komposisi transformasi linear pada \( \Re^1 \)
    bersifat komutatif.
    Apakah ini berlaku untuk sebarang ruang berdimensi satu?
    \begin{answer}
      Ini berlaku untuk semua ruang berdimensi satu.
      Misalkan $f$ dan $g$ adalah transformasi suatu ruang berdimensi satu.
      Kita harus menunjukkan bahwa
      $\composed{g}{f}\,(\vec{v})=\composed{f}{g}\,(\vec{v})$
      untuk semua vektor.
      Tetapkan suatu basis $B$ untuk ruang tersebut; transformasi-transformasinya kemudian
      direpresentasikan oleh matriks $\nbyn{1}$.
      \begin{equation*}
        F=\rep{f}{B,B}=\begin{mat}
                        f_{1,1}
                     \end{mat}
        \qquad
        G=\rep{g}{B,B}=\begin{mat}
                        g_{1,1}
                     \end{mat}
      \end{equation*}
      Dengan demikian, komposisi-komposisinya dapat direpresentasikan sebagai $GF$ dan $FG$.
      \begin{equation*}
        GF=\rep{\composed{g}{f}}{B,B}=\begin{mat}
                        g_{1,1}f_{1,1}
                     \end{mat}
        \qquad
        FG=\rep{\composed{f}{g}}{B,B}=\begin{mat}
                        f_{1,1}g_{1,1}
                     \end{mat}
      \end{equation*}
      Kedua matriks ini sama, sehingga komposisi-komposisinya memberikan
      hasil yang sama pada setiap vektor di ruang tersebut.
     \end{answer}
  \item  
    Mengapa perkalian matriks tidak didefinisikan sebagai perkalian entri demi entri?
    Cara itu akan lebih mudah dan juga komutatif.
    \begin{answer}
      Operasi itu tidak akan merepresentasikan komposisi peta linear;
      \nearbytheorem{th:MatMultRepComp} tidak akan berlaku.
    \end{answer}
  \item \label{exer:NicePropsMatMult} 
    \begin{exparts}
      \partsitem Buktikan bahwa $H^pH^q=H^{p+q}$ dan $(H^p)^q=H^{pq}$
        untuk bilangan bulat positif \( p,q \).
      \partsitem Buktikan bahwa $(rH)^p=r^p\cdot H^p$
        untuk sebarang bilangan bulat positif \( p \) dan skalar \( r\in\Re \).
    \end{exparts}
    \begin{answer}
      Masing-masing mudah mengikuti fakta tentang peta yang bersesuaian.
      Misalnya, $p$ penerapan transformasi $h$ setelah $q$
      penerapan hanyalah $p+q$ penerapan.
    \end{answer}
  \recommended \item \label{exer:MoreNicePropsMatMult} 
    \begin{exparts}
      \partsitem Bagaimana perkalian matriks berinteraksi dengan
        perkalian skalar:~apakah \( r(GH)=(rG)H \)?
        Apakah \( G(rH)=r(GH) \)?
      \partsitem Bagaimana perkalian matriks berinteraksi dengan
        kombinasi linear:~apakah \( F(rG+sH)=r(FG)+s(FH) \)?
        Apakah $(rF+sG)H=rFH+sGH$?
    \end{exparts}
    \begin{answer}  
      Meskipun kita dapat mengerjakannya dengan menelusuri indeks, pernyataan-pernyataan ini
      paling baik dipahami melalui peta yang direpresentasikan.
      Artinya, tetapkan ruang dan basis sedemikian sehingga matriks-matriks tersebut
      merepresentasikan peta linear $f,g,h$.
       \begin{exparts}
        \partsitem Ya; kita memiliki baik
          $r\cdot (\composed{g}{h})\,(\vec{v})
           =r\cdot g(\,h(\vec{v})\,)
           =\composed{(r\cdot g)}{h}\,(\vec{v})$
          maupun
          $\composed{g}{(r\cdot h)}\,(\vec{v})=g(\,r\cdot h(\vec{v})\,)
            =r\cdot g(h(\vec{v}))
            =r\cdot (\composed{g}{h})\,(\vec{v})$
          (kesamaan kedua berlaku karena linearitas $g$).
        \partsitem Jawaban untuk keduanya adalah ya.
          Pertama, $\composed{f}{(rg+sh)}$ dan
          $r\cdot(\composed{f}{g})+s\cdot(\composed{f}{h})$ sama-sama memetakan
          $\vec{v}$ ke $r\cdot f(g(\vec{v}))+s\cdot f(h(\vec{v}))$;
          perhitungannya seperti pada butir sebelumnya
          (dengan menggunakan linearitas $f$ untuk yang pertama).
          Untuk yang lain,
          $\composed{(rf+sg)}{h}$ dan
          $r\cdot(\composed{f}{h})+s\cdot(\composed{g}{h})$ sama-sama memetakan
          $\vec{v}$ ke $r\cdot f(h(\vec{v}))+s\cdot g(h(\vec{v}))$.
      \end{exparts}  
    \end{answer}
  \item \label{exer:TranspAndMult} 
    Kita dapat menanyakan bagaimana operasi perkalian
    matriks berinteraksi dengan operasi transpos.
    \begin{exparts}  
      \partsitem Tunjukkan bahwa \( \trans{(GH)}=\trans{H}\trans{G} \).
      \partsitem Suatu matriks persegi disebut
        \definend{simetris}\index{matriks simetris}%
        \index{matriks!simetris} jika setiap
        entri \( i,j \) sama dengan
        entri \( j,i \), yakni jika matriks itu sama dengan transposnya sendiri.
        Tunjukkan bahwa
        matriks \( H\trans{H} \) dan \( \trans{H}H \) bersifat simetris.
      %\partsitem Is every symmetric matrix of that form?
    \end{exparts}
    \begin{answer}
      Kita belum melihat interpretasi operasi transpos sebagai peta, jadi
      kita akan memverifikasi pernyataan-pernyataan ini dengan meninjau entri-entrinya.
      \begin{exparts}
        \partsitem Entri \( i,j \) dari \( \trans{GH} \) adalah entri $j,i$
          dari $GH$, yaitu hasil kali titik antara
          baris ke-\( j \) dari \( G \) dan kolom ke-\( i \) dari \( H \).
          Entri \( i,j \) dari \( \trans{H}\trans{G} \) adalah hasil kali titik antara
          baris ke-\( i \) dari \( \trans{H} \) dan kolom ke-\( j \) dari
          \( \trans{G} \), yaitu
          hasil kali titik antara kolom ke-\( i \) dari \( H \) dan
          baris ke-\( j \) dari \( G \).
          Hasil kali titik bersifat komutatif, sehingga keduanya sama.
        \partsitem Berdasarkan butir sebelumnya, masing-masing sama dengan transposnya, misalnya,
          $\trans{(H\trans{H})}=\trans{\trans{H}}\trans{H}=H\trans{H}$.
      \end{exparts}  
    \end{answer}
  \recommended \item
    Rotasi vektor dalam \( \Re^3 \) terhadap suatu sumbu adalah peta linear.
    Tunjukkan bahwa peta linear tidak bersifat komutatif dengan
    menunjukkan secara geometris bahwa rotasi tidak bersifat komutatif.
    \begin{answer}
      Perhatikan \( \map{r_x,r_y}{\Re^3}{\Re^3} \) yang merotasikan semua vektor
      sebesar \( \pi/2 \)~radian
      berlawanan arah jarum jam terhadap sumbu \( x \) dan \( y \)
      (berlawanan arah jarum jam dalam arti bahwa seseorang yang kepalanya berada di
      \( \vec{e}_1 \) atau \( \vec{e}_2 \) dan kakinya di titik asal,
      ketika melihat ke arah titik asal, melihat rotasi tersebut sebagai
      berlawanan arah jarum jam).
      \begin{center}  \small
        \includegraphics{map/mp/ch3.78}
        \qquad
        \includegraphics{map/mp/ch3.79}
      \end{center}
      Melakukan rotasi $r_x$ terlebih dahulu lalu $r_y$ berbeda dengan
      melakukan rotasi $r_y$ terlebih dahulu lalu $r_x$.
      Khususnya, $r_x(\vec{e}_3)=-\vec{e}_2$,
      sehingga $\composed{r_y}{r_x}(\vec{e}_3)=-\vec{e}_2$,
      sedangkan $r_y(\vec{e}_3)=\vec{e}_1$, sehingga
      $\composed{r_x}{r_y}(\vec{e}_3)=\vec{e}_1$,
      dan karenanya peta-peta tersebut tidak bersifat komutatif.
    \end{answer}
  \item \label{exer:MatProdPropsSpAndBas}
    Dalam pembuktian \nearbytheorem{th:MatMultWellBehaved}, kita menggunakan beberapa peta.
    Apakah domain dan kodomainnya?
    \begin{answer}
      Tidak menjadi masalah (asalkan ruang-ruangnya memiliki
      dimensi yang sesuai).

      Untuk sifat asosiatif,
      misalkan $F$ berukuran $\nbym{m}{r}$, $G$ berukuran $\nbym{r}{n}$, dan
      $H$ berukuran $\nbym{n}{k}$.
      Kita dapat mengambil sebarang ruang berdimensi~$r$,
      sebarang ruang berdimensi~$m$, sebarang ruang berdimensi~$n$, dan sebarang
      ruang berdimensi~$k$\Dash misalnya,
      $\Re^r$, $\Re^m$, $\Re^n$, dan $\Re^k$ dapat digunakan.
      Kita dapat mengambil sebarang basis $A$, $B$, $C$, dan $D$ untuk ruang-ruang tersebut.
      Kemudian,
      terhadap $C,D$, matriks $H$ merepresentasikan peta linear $h$;
      terhadap $B,C$, matriks $G$ merepresentasikan $g$;
      dan terhadap $A,B$, matriks $F$ merepresentasikan $f$.
      Kita dapat menggunakan peta-peta tersebut dalam pembuktian.
      
      Paruh kedua serupa, kecuali bahwa kita menjumlahkan $G$ dan $H$,
      sehingga keduanya harus merepresentasikan peta dengan domain
      dan kodomain yang sama.
    \end{answer}
  \item \label{exer:RankProdLeqRankFacts}
    Bagaimana rank matriks berinteraksi dengan perkalian matriks?
    \begin{exparts}
      \partsitem Dapatkah hasil kali matriks-matriks ber-rank \( n \) memiliki rank kurang
        dari \( n \)?
        Lebih besar?
      \partsitem Tunjukkan bahwa rank hasil kali dua matriks kurang dari
        atau sama dengan nilai minimum rank masing-masing faktor.
    \end{exparts}
    \begin{answer}
     \begin{exparts}
      \partsitem Hasil kali matriks-matriks ber-rank \( n \) dapat memiliki rank kurang
        dari atau sama dengan \( n \), tetapi tidak lebih besar dari \( n \).

        Untuk melihat bahwa rank dapat turun,
        perhatikan peta \( \map{\pi_x,\pi_y}{\Re^2}{\Re^2} \) yang memproyeksikan ke
        sumbu-sumbu koordinat.
        Masing-masing ber-rank satu, tetapi komposisinya,
        $\composed{\pi_x}{\pi_y}$, yang merupakan peta nol, ber-rank nol.
        Hal ini diterjemahkan ke matriks yang merepresentasikan peta-peta tersebut
        sebagai berikut.
        \begin{equation*}
          \rep{\pi_x}{\stdbasis_2,\stdbasis_2}\cdot\rep{\pi_y}{\stdbasis_2,\stdbasis_2}=
          \begin{mat}[r]
            1  &0  \\
            0  &0
          \end{mat}
          \begin{mat}[r]
            0  &0  \\
            0  &1
          \end{mat}=
          \begin{mat}[r]
            0  &0  \\
            0  &0
          \end{mat}
        \end{equation*}

        Untuk membuktikan bahwa hasil kali matriks-matriks ber-rank \( n \) tidak dapat memiliki rank
        lebih besar dari \( n \), kita dapat menerapkan hasil tentang peta bahwa bayangan suatu
        himpunan yang bergantung linear juga bergantung linear.
        Artinya, jika \( \map{h}{V}{W} \) dan \( \map{g}{W}{X} \) sama-sama ber-rank
        \( n \), maka suatu himpunan dalam daerah hasil
        \( \rangespace{\composed{g}{h}} \) yang berukuran
        lebih besar dari \( n \) adalah bayangan oleh \( g \) dari suatu himpunan dalam \( W \) yang
        berukuran lebih besar dari \( n \), sehingga bergantung linear
        (karena rank \( h \) adalah \( n \)).
        Bayangan himpunan yang bergantung linear juga bergantung, jadi setiap himpunan
        berukuran lebih besar dari \( n \) dalam daerah hasil bergantung linear.
        (Perhatikan bahwa rank \( g \) tidak disebutkan.
        Lihat bagian berikutnya.)
       \partsitem Tetapkan ruang dan basis, lalu perhatikan peta linear terkait
          \( f \) dan \( g \).
          Ingat bahwa dimensi bayangan suatu peta (rank peta tersebut)
          kurang dari atau sama dengan dimensi domainnya, lalu perhatikan
          diagram panah berikut.
          \begin{equation*}
            \begin{matrix}
              V &\mapsunder{f} &\rangespace{f} &\mapsunder{g}
                &\rangespace{\composed{g}{f}}
            \end{matrix}
          \end{equation*}
          Pertama, bayangan dari \( \rangespace{f} \) harus memiliki dimensi
          kurang dari atau sama dengan dimensi \( \rangespace{f} \),
          berdasarkan kalimat sebelumnya.
          Di sisi lain, \( \rangespace{f} \) merupakan subhimpunan dari
          domain \( g \), sehingga bayangannya berdimensi kurang dari
          atau sama dengan dimensi domain \( g \).
          Dengan menggabungkan keduanya,
          rank suatu komposisi kurang dari atau sama dengan nilai minimum
          kedua rank tersebut.

          Fakta tentang matriks segera mengikuti.
     \end{exparts}
    \end{answer}
  \item 
    Apakah relasi `komutatif dengan' merupakan relasi ekuivalensi di antara
    matriks $\nbyn{n}$?
    \begin{answer}
      Relasi `komutatif dengan' bersifat refleksif dan simetris.
      Namun, relasi ini tidak transitif:~misalnya, dengan
      \begin{equation*}
        G=\begin{mat}[r]
          1  &2  \\
          3  &4
        \end{mat}
        \quad
        H=\begin{mat}[r]
          1  &0  \\
          0  &1
        \end{mat}
        \quad
        J=\begin{mat}[r]
          5  &6  \\
          7  &8
        \end{mat}
      \end{equation*}
      $G$ komutatif dengan $H$ dan $H$ komutatif dengan $J$, tetapi $G$ tidak
      komutatif dengan $J$.
    \end{answer}
  \item \label{exer:ZeroDivisor}
    \textit{(Kita akan menggunakan latihan ini dalam latihan Invers Matriks.)}
    Berikut ini sifat lain perkalian matriks yang mungkin
    membingungkan pada pandangan pertama.
    \begin{exparts}
      \partsitem Buktikan bahwa komposisi proyeksi
        \( \map{\pi_x,\pi_y}{\Re^3}{\Re^3} \)
        ke sumbu \( x \) dan \( y \)
        adalah peta nol meskipun tidak satu pun dari keduanya merupakan peta nol.
      \partsitem Buktikan bahwa komposisi turunan
        \( \map{d^2/dx^2,\,d^3/dx^3}{\polyspace_4}{\polyspace_4} \)
        adalah peta nol meskipun tidak satu pun merupakan peta nol.
      \partsitem Berikan persamaan matriks yang merepresentasikan fakta pertama.
      \partsitem Berikan persamaan matriks yang merepresentasikan fakta kedua.
    \end{exparts}
    Jika dua objek dikalikan dan menghasilkan nol meskipun keduanya tidak nol,
kita menyebut masing-masing sebagai
    \definend{pembagi nol}.\index{pembagi nol}%
    % \index{zero!divisor}
    \begin{answer}
      \begin{exparts}
        \partsitem Salah satu dari berikut ini.
          \begin{equation*}
            \colvec{x \\ y \\ z}
              \mapsunder{\pi_x}
            \colvec{x \\ 0 \\ 0}
              \mapsunder{\pi_y}
            \colvec[r]{0 \\ 0 \\ 0}
            \qquad
            \colvec{x \\ y \\ z}
              \mapsunder{\pi_y}
            \colvec{0 \\ y \\ 0}
              \mapsunder{\pi_x}
            \colvec[r]{0 \\ 0 \\ 0}
          \end{equation*}
        \partsitem Komposisinya adalah peta turunan kelima
          $d^5/dx^5$ pada ruang polinomial berderajat empat.
        \partsitem Terhadap basis alami,
          \begin{equation*}
            \rep{\pi_x}{\stdbasis_3,\stdbasis_3}=\begin{mat}[r]
              1  &0  &0  \\
              0  &0  &0  \\
              0  &0  &0
            \end{mat}
            \qquad
            \rep{\pi_y}{\stdbasis_3,\stdbasis_3}=\begin{mat}[r]
              0  &0  &0  \\
              0  &1  &0  \\
              0  &0  &0
            \end{mat}
          \end{equation*}
          dan hasil kali keduanya (dalam urutan mana pun) adalah matriks nol.
        \partsitem Untuk \( B=\sequence{1,x,x^2,x^3,x^4} \),
          \begin{equation*}
            \rep{\frac{d^2}{dx^2}}{B,B}=\begin{mat}[r]
              0  &0  &2  &0  &0  \\
              0  &0  &0  &6  &0  \\
              0  &0  &0  &0  &12 \\
              0  &0  &0  &0  &0  \\
              0  &0  &0  &0  &0
            \end{mat}
            \qquad
            \rep{\frac{d^3}{dx^3}}{B,B}=\begin{mat}[r]
              0  &0  &0  &6  &0  \\
              0  &0  &0  &0  &24 \\
              0  &0  &0  &0  &0  \\
              0  &0  &0  &0  &0  \\
              0  &0  &0  &0  &0
            \end{mat}
          \end{equation*}
          dan hasil kali keduanya (dalam urutan mana pun) adalah matriks nol.
      \end{exparts}  
    \end{answer}
  \item 
    Tunjukkan bahwa, untuk matriks persegi, \( (S+T)(S-T) \) tidak harus sama dengan
    \( S^2-T^2 \).
    \begin{answer}
      Perhatikan bahwa \( (S+T)(S-T)=S^2-ST+TS-T^2 \), jadi masuk akal untuk mencari
      matriks-matriks yang tidak komutatif agar $-ST$ dan $TS$ tidak 
      saling meniadakan:~dengan 
      \begin{equation*}
        S=\begin{mat}[r]
            1  &2  \\
            3  &4
          \end{mat}
        \quad
        T=\begin{mat}[r]
            5  &6  \\
            7  &8
          \end{mat}
      \end{equation*}
      kita memperoleh ketaksamaan yang diinginkan.
      \begin{equation*}
        (S+T)(S-T)=\begin{mat}[r]
            -56  &-56  \\
            -88  &-88
          \end{mat}
        \qquad
        S^2-T^2=\begin{mat}[r]
            -60  &-68  \\
            -76  &-84
          \end{mat}
      \end{equation*}    
    \end{answer}
  \recommended \item \label{exer:IdMat} 
    Representasikan transformasi identitas 
    $\map{\identity}{V}{V}$ terhadap $B,B$ untuk sebarang
    basis $B$.
    Inilah \definend{matriks identitas}\index{matriks!identitas} $I$.
    Tunjukkan bahwa dalam perkalian matriks, matriks ini memainkan peran yang
    dimainkan bilangan $1$ dalam perkalian bilangan real:~$HI=IH=H$ (untuk semua matriks
    $H$ yang hasil kalinya terdefinisi).
    \begin{answer}
      Karena peta identitas bekerja pada basis $B$ sebagai
      $\vec{\beta}_1\mapsto\vec{\beta}_1$, \ldots,
      $\vec{\beta}_n\mapsto\vec{\beta}_n$,
      representasinya adalah sebagai berikut.
      \begin{equation*}
        \begin{mat}
          1  &0  &0  &  &0 \\
          0  &1  &0  &  &0 \\
          0  &0  &1  &  &0 \\
             &   &   &\ddots  \\
          0  &0  &0  &  &1
        \end{mat}
      \end{equation*}
      Bagian kedua soal ini langsung mengikuti 
      \nearbytheorem{th:MatMultRepComp}.
    \end{answer}
  \item 
    \begin{exparts}
      \partsitem Buktikan bahwa untuk setiap matriks \( \nbyn{2} \) \( T \) 
        terdapat skalar-skalar
        \( c_0,\dots,c_4 \) yang tidak semuanya $0$ sehingga kombinasi
        $c_4T^4+c_3T^3+c_2T^2+c_1T+c_0I$
        merupakan matriks nol
        (dengan $I$ matriks identitas $\nbyn{2}$, yang entri
        $1,1$ dan $2,2$-nya bernilai $1$ dan entri lainnya bernilai nol;~lihat
        \nearbyexercise{exer:IdMat}).
      \partsitem Misalkan \( p(x) \) adalah polinomial 
        \( p(x)=c_nx^n+\dots+c_1x+c_0 \).
        Jika \( T \) adalah matriks persegi, kita definisikan \( p(T) \) sebagai matriks
        \( c_nT^n+\dots+c_1T+c_0I \)
        (dengan $I$ matriks identitas berukuran sesuai).
        Buktikan bahwa untuk setiap matriks persegi terdapat suatu polinomial sehingga
        \( p(T) \) merupakan matriks nol.
      \partsitem \definend{Polinomial minimal}\index{polinomial minimal}%
        \index{matriks!polinomial minimal}
        $m(x)$ dari suatu matriks persegi adalah
        polinomial berderajat paling rendah dan berkoefisien utama \( 1 \),
        sedemikian sehingga \( m(T) \) merupakan matriks nol.
        Carilah polinomial minimal matriks berikut.
        \begin{equation*}
          \begin{mat}[r]
            \sqrt{3}/2  &-1/2       \\
            1/2         &\sqrt{3}/2
          \end{mat}
        \end{equation*}
        (Ini adalah representasi, terhadap basis standar
        $\stdbasis_2,\stdbasis_2$, dari rotasi sejauh $\pi/6$~radian
        berlawanan arah jarum jam.)
    \end{exparts}
    \begin{answer} 
      \begin{exparts}
         \partsitem Ruang vektor \( \matspace_{\nbyn{2}} \) berdimensi
           empat. 
           Himpunan \( \set{T^4,\dots,T,I} \) memiliki lima
           elemen sehingga bergantung linear.
         \partsitem Jika \( T \) berukuran \( \nbyn{n} \), 
           generalisasi argumen dari butir
           sebelumnya menunjukkan bahwa terdapat polinomial semacam itu yang berderajat 
           paling tinggi \( n^2 \),
           karena \( \set{T^{n^2},\dots,T,I} \) adalah subhimpunan beranggota \( n^2+1 \)
           dari ruang berdimensi $n^2$, $\matspace_{\nbyn{n}}$.
         \partsitem Pertama, hitung pangkat-pangkatnya
           \begin{equation*}
              T^2=
              \begin{mat}[r]
                 1/2         &-\sqrt{3}/2  \\
                 \sqrt{3}/2  &1/2
              \end{mat} 
              \qquad
              T^3=
              \begin{mat}[r]
                 0           &-1           \\
                 1           &0
              \end{mat}
             \qquad
              T^4=
              \begin{mat}[r]
                 -1/2        &-\sqrt{3}/2  \\
                 \sqrt{3}/2  &-1/2
              \end{mat}
           \end{equation*}
           (perhatikan bahwa rotasi sebesar $\pi/6$ sebanyak tiga kali menghasilkan 
           rotasi sebesar $\pi/2$, yang memang direpresentasikan oleh $T^3$).
           Kemudian samakan \( c_4T^4+c_3T^3+c_2T^2+c_1T+c_0I \) 
           dengan matriks nol 
           \begin{multline*}
             \begin{mat}[r]
                 -1/2        &-\sqrt{3}/2  \\
                 \sqrt{3}/2  &-1/2
              \end{mat}c_4
              +\begin{mat}[r]
                 0           &-1           \\
                 1           &0
              \end{mat}c_3
              +\begin{mat}[r]
                 1/2         &-\sqrt{3}/2  \\
                 \sqrt{3}/2  &1/2
              \end{mat}c_2                                     \\
              +\begin{mat}[r]
                 \sqrt{3}/2  &-1/2       \\
                 1/2         &\sqrt{3}/2
               \end{mat}c_1
              +\begin{mat}[r]
                 1  &0  \\
                 0  &1 
               \end{mat}c_0                          
               =\begin{mat}[r]
                  0  &0  \\
                  0  &0 
               \end{mat}                     
           \end{multline*}
           untuk memperoleh sistem linear berikut.
           \begin{equation*}
             \begin{linsys}{5}
                -(1/2)c_4         &  &  &+ &(1/2)c_2
                    &+ &(\sqrt{3}/2)c_1  &+  &c_0  &=  &0      \\
                -(\sqrt{3}/2)c_4  &- &c_3 &- &(\sqrt{3}/2)c_2
                    &- &(1/2)c_1         &   &          &=  &0        \\
                 (\sqrt{3}/2)c_4  &+ &c_3 &+ &(\sqrt{3}/2)c_2
                    &+ &(1/2)c_1  &   &            &=  &0      \\
                -(1/2)c_4              &  &  &+ &(1/2)c_2
                    &+ &(\sqrt{3}/2)c_1  &+  &c_0  &=  &0
              \end{linsys} 
           \end{equation*}
           Terapkan reduksi Gauss.
           \begin{align*}
             &\grstep{-\rho_1+\rho_4}
             \repeatedgrstep{\rho_2+\rho_3}
             \begin{linsys}{5}
                -(1/2)c_4         &  &  &+ &(1/2)c_2
                    &+ &(\sqrt{3}/2)c_1  &+  &c_0  &=  &0      \\
                -(\sqrt{3}/2)c_4  &- &c_3 &- &(\sqrt{3}/2)c_2
                    &- &(1/2)c_1         &   &          &=  &0   \\
                                  &  &  &  &     
                    &  &                 &   &0    &=  &0      \\
                                  &  &  &  &     
                    &  &                 &   &0    &=  &0      
             \end{linsys}                                               \\
             &\grstep{-\sqrt{3}\rho_1+\rho_2}                   
             \begin{linsys}{5}
                -(1/2)c_4         &  &  &+ &(1/2)c_2
                    &+ &(\sqrt{3}/2)c_1  &+  &c_0  &=  &0      \\
                                  &- &c_3 &- &\sqrt{3}c_2
                    &- &2c_1             &-  &\sqrt{3}c_0  &=  &0   \\
                                  &  &  &  &     
                    &  &                 &   &0    &=  &0      \\
                                  &  &  &  &     
                    &  &                 &   &0    &=  &0      
              \end{linsys} 
           \end{align*}
           Menetapkan \( c_4 \), \( c_3 \), dan \( c_2 \) sama dengan nol 
           membuat \( c_1 \) dan \( c_0 \) juga bernilai nol, sehingga
           tidak ada polinomial berderajat satu ataupun nol yang memenuhi.
           Menetapkan \( c_4 \) dan \( c_3 \) sama dengan nol (serta $c_2$ sama dengan satu)
           menghasilkan sistem linear 
           \begin{equation*}
             \begin{linsys}{3}
               (1/2)     &+  &(\sqrt{3}/2)c_1  &+  &c_0          &=  &0 \\
               -\sqrt{3} &-  &2c_1             &-  &\sqrt{3}c_0  &=  &0
             \end{linsys}
           \end{equation*}
           dengan solusi $c_1=-\sqrt{3}$ dan $c_0=1$.
           Kesimpulan:~polinomial 
           $m(x)=x^2-\sqrt{3}x+1$
           merupakan polinomial minimal bagi matriks $T$.
      \end{exparts}  
    \end{answer}
  \item 
    Ruang berdimensi takhingga $\polyspace$ 
    yang terdiri dari semua polinomial berderajat hingga
    memberikan contoh yang mudah diingat tentang ketidakkomutatifan
    peta-peta linear.
    Misalkan \( \map{d/dx}{\polyspace}{\polyspace} \) adalah turunan biasa
    dan \( \map{s}{\polyspace}{\polyspace} \) adalah peta \definend{geser}.
    \begin{equation*}
      a_0+a_1x+\dots+a_nx^n
      \;\mapsunder{s}\;
      0+a_0x+a_1x^2+\dots+a_nx^{n+1}
    \end{equation*}
    Tunjukkan bahwa kedua peta tersebut tidak komutatif,
    \( \composed{d/dx}{s}\neq \composed{s}{d/dx} \); bahkan, 
    \( (\composed{d/dx}{s})-(\composed{s}{d/dx}) \) bukan saja bukan
    peta nol, melainkan peta identitas.
    \begin{answer}
      Pemeriksaannya rutin:
      \begin{equation*}
        a_0+a_1x+\dots+a_nx^n\mapsunder{s}a_0x+a_1x^2+\dots+a_nx^{n+1}
                  \mapsunder{d/dx} a_0+2a_1x+\dots+(n+1)a_nx^n
      \end{equation*}
      sedangkan
      \begin{equation*}
        a_0+a_1x+\dots+a_nx^n
                  \mapsunder{d/dx} a_1+\dots+na_nx^{n-1}
                  \mapsunder{s}a_1x+\dots+a_nx^n
      \end{equation*}
      sehingga di bawah peta $(\composed{d/dx}{s})-(\composed{s}{d/dx})$ 
      kita memperoleh 
        $a_0+a_1x+\dots+a_nx^n
                  \mapsto
                  a_0+a_1x+\dots+a_nx^n$.
    \end{answer}
  \item 
    Ingat kembali notasi untuk jumlah barisan bilangan
    \( a_1, a_2, \dots, a_n \).
    \begin{equation*}
      \sum_{i=1}^{n}a_i=a_1+a_2+\dots+a_n
    \end{equation*}
    Dalam notasi ini, entri \( i,j \) dari hasil kali \( G \) dan
    \( H \) adalah sebagai berikut.
    \begin{equation*}
       p_{i,j}=\sum_{k=1}^{r} g_{i,k}h_{k,j}
    \end{equation*}
    Dengan menggunakan notasi ini,
    \begin{exparts}
       \partsitem buktikan kembali bahwa perkalian matriks bersifat asosiatif;
       \partsitem buktikan kembali \nearbytheorem{th:MatMultRepComp}.
    \end{exparts}
    \begin{answer}
       \begin{exparts}
         \partsitem Menelusuri uraian di akhir subbagian
           menunjukkan bahwa entri \( i,j \) dari \( (FG)H \) adalah
           \begin{multline*}
             \sum_{t=1}^s\bigl(\sum_{k=1}^{r} f_{i,k}g_{k,t}\bigr)h_{t,j}
             =\sum_{t=1}^s\sum_{k=1}^{r} (f_{i,k}g_{k,t})h_{t,j} 
             =\sum_{t=1}^s\sum_{k=1}^{r} f_{i,k}(g_{k,t}h_{t,j})  \\ 
             =\sum_{k=1}^{r}\sum_{t=1}^s f_{i,k}(g_{k,t}h_{t,j}) 
             =\sum_{k=1}^{r}f_{i,k}\bigl(\sum_{t=1}^s g_{k,t}h_{t,j}\bigr)
           \end{multline*}
           (kesamaan pertama diperoleh dengan menggunakan 
           hukum distributif untuk mengalikan setiap suku dengan
           $h$, kesamaan kedua menggunakan hukum asosiatif bilangan
           real, kesamaan ketiga menggunakan hukum komutatif bilangan real,
           dan kesamaan keempat diperoleh dengan menggunakan hukum distributif untuk
           memfaktorkan $f$),
           yang merupakan entri \( i,j \) dari \( F(GH) \).
         \partsitem Komponen ke-\( k \) dari \( h(\vec{v})\) adalah
           \begin{equation*}
             \sum_{j=1}^n h_{k,j}v_j
           \end{equation*}
           sehingga komponen ke-\( i \) dari 
           \( \composed{g}{h}\,(\vec{v}) \) adalah
           \begin{multline*}
              \sum_{k=1}^r g_{i,k}\bigl(\sum_{j=1}^n h_{k,j}v_j\bigr) 
              =\sum_{k=1}^r \sum_{j=1}^n g_{i,k}h_{k,j}v_j   
              =\sum_{k=1}^r \sum_{j=1}^n (g_{i,k}h_{k,j})v_j
                                                              \\ 
              =\sum_{j=1}^n \sum_{k=1}^r (g_{i,k}h_{k,j})v_j 
              =\sum_{j=1}^n (\sum_{k=1}^r g_{i,k}h_{k,j})\,v_j
           \end{multline*}
           (kesamaan pertama berlaku karena hukum distributif digunakan untuk mengalikan
           setiap suku dengan $g$, kesamaan kedua menggunakan 
           sifat asosiatif bilangan real, kesamaan ketiga mengikuti dari sifat komutatif
           bilangan real, dan kesamaan keempat diperoleh dengan menggunakan 
           hukum distributif untuk memfaktorkan $v$).
       \end{exparts} 
    \end{answer}
\end{exercises}





















\subsection{Mekanika Perkalian Matriks}
Kita dapat memandang perkalian matriks sebagai suatu proses mekanis, 
sementara mengesampingkan segala implikasi tentang peta-peta yang mendasarinya.

Hal yang mencolok dari operasi ini adalah
cara baris dan kolom dipadukan.
Entri \( i,j \) pada hasil kali matriks merupakan hasil kali titik antara
baris~$i$ matriks kiri dan kolom~$j$ matriks kanan.
Sebagai contoh,
di sini baris kedua dan kolom ketiga dipadukan untuk menghasilkan entri~$2,3$.
\begin{equation*}
\setlength{\fboxsep}{1.5pt}
    \begin{mat}
       \begin{array}{@{}cc@{}} 1  &1 \end{array}                         \\ 
       \text{\highlight{\( \begin{array}{@{}cc@{}}  0  &1  \end{array} \)}}   \\  
       \begin{array}{@{}cc@{}} 1  &0 \end{array}
    \end{mat}
    \begin{mat}
      \begin{array}{@{}c@{}}  4  \\  5  \end{array}
      &\begin{array}{@{}c@{}}  6  \\  7  \end{array}
      &\text{\highlight{ \(\begin{array}{@{}c@{}}  8  \\  9  \end{array}\) }}
      &\begin{array}{@{}c@{}}  2  \\  3  \end{array}
    \end{mat}
  =
    \begin{mat}[r]
      9  &13   &17                      &5  \\
      5  &7    &\text{\highlight{ \(9\) }}   &3  \\
      4  &6    &8                       &2
    \end{mat}
\end{equation*}
Kita dapat memandangnya sebagai matriks kiri yang bekerja
dengan mengalikan baris-barisnya dengan kolom-kolom matriks kanan.
Atau, sebagai
matriks kanan yang menggunakan kolom-kolomnya untuk
bekerja pada baris-baris matriks kiri.
Di bawah ini kita akan menelaah aksi dari kiri dan dari kanan untuk beberapa
matriks sederhana.

Yang paling sederhana adalah matriks nol.

\begin{example}
Mengalikan dengan matriks nol
dari kiri ataupun dari kanan
menghasilkan matriks nol.
\begin{equation*}
    \begin{mat}[r]
       0  &0  \\
       0  &0
    \end{mat}
    \begin{mat}[r]
       1  &3  &2   \\
      -1  &1  &-1
    \end{mat}
  =
    \begin{mat}[r]
       0  &0  &0   \\
       0  &0  &0
    \end{mat}
    \qquad
    \begin{mat}[r]
       2  &3  \\
       1  &4
    \end{mat}
    \begin{mat}[r]
       0  &0  \\
       0  &0
    \end{mat}
  =
    \begin{mat}[r]
       0  &0  \\
       0  &0
    \end{mat}
\end{equation*}
\end{example}

Matriks berikutnya yang paling sederhana
adalah matriks dengan satu entri taknol.

\begin{definition}  \label{df:UnitMatrix}
%<*df:UnitMatrix>
Matriks yang semua entrinya $0$, kecuali sebuah $1$ pada entri \( i,j \),
adalah matriks \definend{satuan}\index{matriks!satuan}\index{matriks satuan} 
\( i,j \)
(atau \definend{satuan matriks}).\index{matriks!satuan}
%</df:UnitMatrix>
\end{definition}

\begin{example}
Berikut adalah matriks satuan \( 1,2\, \) dengan tiga baris dan dua kolom,
yang mengalikan dari kiri.
\begin{equation*}
    \begin{mat}[r]
       0  &1  \\
       0  &0  \\
       0  &0
    \end{mat}
    \begin{mat}[r]
       5  &6  \\
       7  &8
    \end{mat}
  =
    \begin{mat}[r]
       7  &8  \\
       0  &0  \\
       0  &0
    \end{mat}
\end{equation*}
Ketika bekerja
dari kiri, matriks satuan \( i,j \) menyalin baris~\( j \) dari
matriks yang dikalikan ke baris~\( i \) hasilnya.
Dari kanan, matriks satuan \( i,j \) memilih kolom~\( i \) dari
matriks yang dikalikan dan menyalinnya ke kolom~\( j \) hasilnya.
\begin{equation*}
    \begin{mat}[r]
       1  &2  &3  \\
       4  &5  &6  \\
       7  &8  &9
    \end{mat}
    \begin{mat}[r]
       0  &1  \\
       0  &0  \\
       0  &0
    \end{mat}
  =
    \begin{mat}[r]
       0  &1  \\
       0  &4  \\
       0  &7
    \end{mat}
\end{equation*}
\end{example}

\begin{example}
Mengalikan matriks satuan dengan suatu skalar cukup mengalikan hasilnya dengan skalar itu. 
Berikut adalah aksi dari kiri oleh matriks yang besarnya dua kali matriks pada
contoh sebelumnya.
\begin{equation*}
    \begin{mat}[r]
       0  &2  \\
       0  &0  \\
       0  &0
    \end{mat}
    \begin{mat}[r]
       5  &6  \\
       7  &8
    \end{mat}
  =
    \begin{mat}[r]
      14  &16 \\
       0  &0  \\
       0  &0
    \end{mat}
\end{equation*}
% And this is the action of the matrix that is $-3$~times the one 
% above.
% \begin{equation*}
%     \begin{mat}[r]
%        1  &2  &3  \\
%        4  &5  &6  \\
%        7  &8  &9
%     \end{mat}
%     \begin{mat}[r]
%        0  &-3 \\
%        0  &0  \\
%        0  &0
%     \end{mat}
%   =
%     \begin{mat}[r]
%        0  &-3  \\
%        0  &-12 \\
%        0  &-21
%     \end{mat}
% \end{equation*}
\end{example}

Tingkat kerumitan berikutnya adalah matriks dengan dua entri taknol.

\begin{example}
Ada dua kasus.
Jika sebuah pengali-kiri memiliki entri-entri pada baris yang berbeda, aksi keduanya 
tidak saling berinteraksi.
\begin{align*}
    \begin{mat}[r]
       1  &0  &0  \\
       0  &0  &2  \\
       0  &0  &0
    \end{mat}
    \begin{mat}[r]
       1  &2  &3  \\
       4  &5  &6  \\
       7  &8  &9
    \end{mat}
  &=(  \begin{mat}[r]
          1  &0  &0  \\
          0  &0  &0  \\
          0  &0  &0
       \end{mat}
    +
       \begin{mat}[r]
          0  &0  &0  \\
          0  &0  &2  \\
          0  &0  &0
       \end{mat}     )
      \begin{mat}[r]
         1  &2  &3  \\
         4  &5  &6  \\
         7  &8  &9
      \end{mat}                      \\
  &=\begin{mat}[r]
       1  &2  &3  \\
       0  &0  &0  \\
       0  &0  &0
    \end{mat}
  +
    \begin{mat}[r]
       0  &0  &0  \\
      14  &16 &18 \\
       0  &0  &0
    \end{mat}                        \\
  &=\begin{mat}[r]
       1  &2  &3  \\
      14  &16 &18 \\
       0  &0  &0
    \end{mat}
\end{align*}
Namun, jika entri-entri taknol pengali-kiri berada pada baris yang sama,
baris hasil tersebut merupakan suatu kombinasi.
\begin{align*}
    \begin{mat}[r]
       1  &0  &2  \\
       0  &0  &0  \\
       0  &0  &0
    \end{mat}
    \begin{mat}[r]
       1  &2  &3  \\
       4  &5  &6  \\
       7  &8  &9
    \end{mat}
  &=(  \begin{mat}[r]
          1  &0  &0  \\
          0  &0  &0  \\
          0  &0  &0
       \end{mat}
    +
       \begin{mat}[r]
          0  &0  &2  \\
          0  &0  &0  \\
          0  &0  &0
       \end{mat}     )
      \begin{mat}[r]
         1  &2  &3  \\
         4  &5  &6  \\
         7  &8  &9
      \end{mat}                     \\
  &=\begin{mat}[r]
       1  &2  &3  \\
       0  &0  &0  \\
       0  &0  &0
    \end{mat}
  +
    \begin{mat}[r]
      14  &16 &18 \\
       0  &0  &0  \\
       0  &0  &0
    \end{mat}                      \\
  &=\begin{mat}[r]
      15  &18 &21 \\
       0  &0  &0  \\
       0  &0  &0
    \end{mat}
\end{align*}
\end{example}

\par\noindent Perkalian dari kanan bekerja dengan cara yang sama, tetapi pada kolom.


\begin{lemma} \label{lm:ColsAndRowsInMatrixMult}
%<*lm:ColsAndRowsInMatrixMult>
Dalam hasil kali dua matriks $G$ dan $H$,
kolom-kolom $GH$ dibentuk dengan mengalikan $G$ dengan kolom-kolom $H$
    \begin{equation*}
      G\cdot \begin{pmat}{c@{\hspace{1em}}c@{\hspace{1em}}c}
         % \MTFlushSpaceAbove
         % \vdotswithin{\vec{h}_1}       &        &\vdotswithin{\vec{h}_n}   
         % \MTFlushSpaceBelow
         \vdots       &        &\vdots   \\
         \vec{h}_1    &\cdots  &\vec{h}_n \\ 
         \vdots       &        &\vdots 
       \end{pmat}
      =\begin{pmat}{c@{\hspace{1em}}c@{\hspace{1em}}c}
         \vdots             &        &\vdots    \\
         G\cdot \vec{h}_1   &\cdots  &G\cdot\vec{h}_n \\ 
         \vdots             &        &\vdots 
       \end{pmat}
    \end{equation*}
dan baris-baris $GH$ dibentuk dengan mengalikan baris-baris $G$ dengan $H$
    \begin{equation*}
      \begin{pmat}{c}
        \cdots\; \vec{g}_1 \;\cdots \rule[-1.25ex]{0pt}{2ex} \\ 
         % \hline
         \vdots \\[.5ex]  
         % \hline
        \rule{0pt}{2.5ex}\cdots\; \vec{g}_r \;\cdots 
       \end{pmat}\cdot H
      =\begin{pmat}{c}
         \cdots\; \vec{g}_1\cdot H \;\cdots\rule[-1.25ex]{0pt}{2ex} \\ 
        % \hline
         \vdots\rule[-1ex]{0pt}{2ex} \\[.5ex]   
        % \hline
         \rule{0pt}{2.5ex}\cdots\; \vec{g}_r\cdot H \;\cdots
       \end{pmat}
    \end{equation*}
% (ignoring the extra parentheses).
%</lm:ColsAndRowsInMatrixMult>
\end{lemma}

\begin{proof}
%<*pf:ColsAndRowsInMatrixMult>
Kita akan memeriksa bahwa dalam hasil kali matriks-matriks $\nbyn{2}$, 
baris-baris hasil kalinya sama dengan
hasil kali baris-baris~$G$ dengan seluruh matriks~$H$; kasus-kasus
lain bekerja dengan cara yang sama.
\begin{align*}
  \begin{mat}
    g_{1,1}  &g_{1,2}    \\
    g_{2,1}  &g_{2,2}    
  \end{mat}
  \begin{mat}
    h_{1,1}  &h_{1,2}    \\
    h_{2,1}  &h_{2,2}    
  \end{mat}
  &=\begin{mat}
    \rowvec{g_{1,1}  &g_{1,2}}H  \\ 
    \rowvec{g_{2,1}  &g_{2,2}}H    
  \end{mat}          \\
  &=
  \begin{mat}
    \rowvec{g_{1,1}h_{1,1}+g_{1,2}h_{2,1}  &g_{1,1}h_{1,2}+g_{1,2}h_{2,2}}    \\ 
    \rowvec{g_{2,1}h_{1,1}+g_{2,2}h_{2,1}  &g_{2,1}h_{1,2}+g_{2,2}h_{2,2}} 
  \end{mat}
\end{align*}
(Kita abaikan tanda kurung tambahan.)
%</pf:ColsAndRowsInMatrixMult>
\end{proof}

\begin{example}
Perhatikan kolom-kolom hasil kali dua matriks~$\nbyn{2}$.
\begin{equation*}
  \begin{mat}
    g_{1,1}  &g_{1,2}    \\
    g_{2,1}  &g_{2,2}    
  \end{mat}
  \begin{mat}
    h_{1,1}  &h_{1,2}    \\
    h_{2,1}  &h_{2,2}    
  \end{mat}   
  =
  \begin{mat}
    g_{1,1}h_{1,1}+g_{1,2}h_{2,1}  &g_{1,1}h_{1,2}+g_{1,2}h_{2,2} \\
    g_{2,1}h_{1,1}+g_{2,2}h_{2,1}  &g_{2,1}h_{1,2}+g_{2,2}h_{2,2} 
  \end{mat}
\end{equation*}
Setiap kolom merupakan hasil perkalian $G$ 
dengan kolom $H$ yang bersesuaian.
\begin{equation*}
  G\colvec{h_{1,1} \\ h_{2,1}}  
  =
  \colvec{g_{1,1}h_{1,1}+g_{1,2}h_{2,1} \\ g_{2,1}h_{1,1}+g_{2,2}h_{2,1}}
  \quad
  G\colvec{h_{1,2}  \\ h_{2,2}}    
  =\colvec{g_{1,1}h_{1,2}+g_{1,2}h_{2,2} \\ g_{2,1}h_{1,2}+g_{2,2}h_{2,2}} 
\end{equation*}
\end{example}

Salah satu penerapan pengamatan tersebut adalah adanya suatu matriks
yang sekadar menyalin baris dan kolom.

\begin{definition}  \label{df:MainDiagonal}
%<*df:MainDiagonal>
\definend{Diagonal utama}\index{matriks!diagonal utama}
(atau \definend{diagonal prinsipal}, atau cukup \definend{diagonal}) 
suatu matriks persegi membentang dari kiri atas ke kanan bawah.
%</df:MainDiagonal>
\end{definition}

\begin{definition} \label{df:IdentityMatrix}
%<*df:IdentityMatrix>
Sebuah \definend{matriks identitas\/}\index{identitas!matriks}\index{matriks!identitas}
berbentuk persegi dan setiap entrinya bernilai $0$, kecuali entri-entri $1$ pada diagonal utama.
\begin{equation*}
   I_{\nbyn{n}}=
      \begin{mat}
        1  &0  &\ldots  &0  \\
        0  &1  &\ldots  &0  \\
           &\vdots          \\
        0  &0  &\ldots  &1
      \end{mat}
\end{equation*}
%</df:IdentityMatrix>
\end{definition}

\begin{example}
Berikut adalah matriks identitas \( \nbyn{2} \) yang membiarkan matriks yang dikalikannya tidak berubah
ketika bekerja dari kanan.
\begin{equation*}
    \begin{mat}[r]
       1  &-2 \\
       0  &-2 \\
       1  &-1 \\
       4  &3
    \end{mat}
    \begin{mat}[r]
       1  &0  \\
       0  &1  \\
    \end{mat}
  =
    \begin{mat}[r]
       1  &-2 \\
       0  &-2 \\
       1  &-1 \\
       4  &3
    \end{mat}
\end{equation*}
\end{example}

\begin{example}
Di sini identitas \( \nbyn{3} \)
membiarkan matriks yang dikalikannya tidak berubah, baik dari kiri
\begin{equation*}
    \begin{mat}[r]
       1  &0  &0  \\
       0  &1  &0  \\
       0  &0  &1
    \end{mat}
    \begin{mat}[r]
       2  &3  &6  \\
       1  &3  &8  \\
      -7  &1  &0
    \end{mat}
  =
    \begin{mat}[r]
       2  &3  &6  \\
       1  &3  &8  \\
      -7  &1  &0
    \end{mat}
\end{equation*}
dan dari kanan.
\begin{equation*}
    \begin{mat}[r]
       2  &3  &6  \\
       1  &3  &8  \\
      -7  &1  &0
    \end{mat}
    \begin{mat}[r]
       1  &0  &0  \\
       0  &1  &0  \\
       0  &0  &1
    \end{mat}
  =
    \begin{mat}[r]
       2  &3  &6  \\
       1  &3  &8  \\
      -7  &1  &0
    \end{mat}
\end{equation*}
\end{example}

\noindent Singkatnya, matriks identitas merupakan elemen identitas 
dari himpunan matriks $\nbyn{n}$ terhadap
operasi perkalian matriks.

Kita dapat memperumum matriks identitas
dengan mengganti bilangan satu pada diagonalnya dengan bilangan real sebarang. 
Matriks yang dihasilkan mengalikan seluruh baris atau kolom dengan skalar.

\begin{definition}  \label{df:DiagonalMatrix}
%<*df:DiagonalMatrix>
Sebuah \definend{matriks diagonal\/}\index{matriks!diagonal}\index{matriks diagonal}
berbentuk persegi dan memiliki entri $0$ di luar diagonal utama.
\begin{equation*}
    \begin{mat}
       a_{1,1}  &0        &\ldots     &0     \\
       0        &a_{2,2}  &\ldots     &0     \\
                &\vdots                      \\
       0        &0        &\ldots     &a_{n,n}
    \end{mat}
\end{equation*}
%</df:DiagonalMatrix>
\end{definition}

\begin{example}
Dari kiri, 
aksi perkalian oleh matriks diagonal adalah
mengalikan baris-baris dengan skalar.
\begin{equation*}
    \begin{mat}[r]
       2  &0   \\
       0  &-1
    \end{mat}
    \begin{mat}[r]
       2  &1  &4  &-1  \\
      -1  &3  &4  &4
    \end{mat}
  =
    \begin{mat}[r]
       4  &2  &8  &-2  \\
       1  &-3 &-4 &-4
    \end{mat}
\end{equation*}
Dari kanan, matriks semacam itu mengalikan kolom-kolom dengan skalar.
\begin{equation*}
     \begin{mat}[r]
        1  &2  &1  \\
        2  &2  &2
     \end{mat}
    \begin{mat}[r]
      3  &0  &0  \\
      0  &2  &0  \\
      0  &0  &-2
    \end{mat}
  =
     \begin{mat}[r]
        3  &4  &-2 \\
        6  &4  &-4
     \end{mat}
\end{equation*}
\end{example}

Kita juga dapat memperumum matriks identitas dengan 
menempatkan tepat satu bilangan satu pada setiap baris dan kolom
dengan susunan selain di sepanjang diagonal.

\begin{definition}  \label{df:PermutationMatrix}
%<*df:PermutationMatrix>
Sebuah \definend{matriks permutasi}%
\index{permutasi!matriks}\index{matriks!permutasi}
berbentuk persegi dan semua entrinya $0$, kecuali
tepat satu~$1$ pada setiap baris dan kolom.
%</df:PermutationMatrix>
\end{definition}

\begin{example}
Dari kiri, matriks-matriks ini mempermutasikan baris.
\begin{equation*}
     \begin{mat}[r]
        0  &0  &1  \\
        1  &0  &0  \\
        0  &1  &0
     \end{mat}
     \begin{mat}[r]
        1  &2  &3  \\
        4  &5  &6  \\
        7  &8  &9
     \end{mat}
  =
     \begin{mat}[r]
        7  &8  &9  \\
        1  &2  &3  \\
        4  &5  &6
     \end{mat}
\end{equation*}
Dari kanan, matriks-matriks ini mempermutasikan kolom.
\begin{equation*}
     \begin{mat}[r]
        1  &2  &3  \\
        4  &5  &6  \\
        7  &8  &9
     \end{mat}
     \begin{mat}[r]
        0  &0  &1  \\
        1  &0  &0  \\
        0  &1  &0
     \end{mat}
  =
     \begin{mat}[r]
        2  &3  &1  \\
        5  &6  &4  \\
        8  &9  &7
     \end{mat}
\end{equation*}
\end{example}

Kita menutup subbagian ini dengan menerapkan pengamatan tersebut untuk memperoleh matriks-matriks 
yang menjalankan Metode Gauss dan reduksi Gauss--Jordan.
Kita telah
melihat cara menghasilkan matriks yang mengalikan baris dengan skalar dan matriks penukar baris.

\begin{example}
Mengalikan dengan matriks ini mengalikan baris kedua 
dengan tiga.
\begin{equation*}
    \begin{mat}[r]
      1  &0  &0  \\
      0  &3  &0  \\
      0  &0  &1
    \end{mat}
    \begin{mat}[r]
      0  &2    &1  &1  \\
      0  &1/3  &1  &-1 \\
      1  &0    &2  &0
    \end{mat}
  =
    \begin{mat}[r]
      0  &2    &1  &1  \\
      0  &1    &3  &-3 \\
      1  &0    &2  &0
    \end{mat}
\end{equation*}
\end{example}

\begin{example}
Perkalian ini menukar baris pertama dan ketiga.
\begin{equation*}
    \begin{mat}[r]
      0  &0  &1  \\
      0  &1  &0  \\
      1  &0  &0
    \end{mat}
    \begin{mat}[r]
      0  &2    &1  &1  \\
      0  &1    &3  &-3 \\
      1  &0    &2  &0
    \end{mat}
  =
    \begin{mat}[r]
      1  &0    &2  &0  \\
      0  &1    &3  &-3 \\
      0  &2    &1  &1
    \end{mat}
\end{equation*}
\end{example}

Untuk melihat cara melakukan kombinasi baris, 
kita amati sesuatu tentang kedua contoh tersebut.
Matriks yang mengalikan baris kedua dengan faktor tiga diperoleh dengan cara 
berikut dari identitas.
\begin{equation*}
    \begin{mat}[r]
      1  &0  &0  \\
      0  &1  &0  \\
      0  &0  &1
    \end{mat}
  \grstep{3\rho_2}
    \begin{mat}[r]
      1  &0  &0  \\
      0  &3  &0  \\
      0  &0  &1
    \end{mat}
\end{equation*}
Demikian pula, matriks yang menukar baris pertama dan ketiga diperoleh dengan cara berikut.
\begin{equation*}
    \begin{mat}[r]
      1  &0  &0  \\
      0  &1  &0  \\
      0  &0  &1
    \end{mat}
  \grstep{\rho_1\leftrightarrow\rho_3}
    \begin{mat}[r]
      0  &0  &1  \\
      0  &1  &0  \\
      1  &0  &0
    \end{mat}
\end{equation*}

\begin{example}
Matriks \( \nbyn{3} \) yang diperoleh sebagai
\begin{equation*}
    \begin{mat}[r]
      1  &0  &0  \\
      0  &1  &0  \\
      0  &0  &1
    \end{mat}
  \grstep{-2\rho_2+\rho_3}
    \begin{mat}[r]
      1  &0  &0  \\
      0  &1  &0  \\
      0  &-2 &1
    \end{mat}
\end{equation*}
ketika bekerja dari kiri,
akan melakukan operasi kombinasi \( -2\rho_2+\rho_3 \).
\begin{equation*}
    \begin{mat}[r]
      1  &0  &0  \\
      0  &1  &0  \\
      0  &-2 &1
    \end{mat}
    \begin{mat}[r]
      1  &0    &2  &0  \\
      0  &1    &3  &-3 \\
      0  &2    &1  &1
    \end{mat}
  =
    \begin{mat}[r]
      1  &0    &2  &0  \\
      0  &1    &3  &-3 \\
      0  &0    &-5 &7
    \end{mat}
\end{equation*}
\end{example}

\begin{definition}   \label{df:ElementaryReductionMatrices}
%<*df:ElementaryReductionMatrices>
\definend{Matriks reduksi elementer}%
\index{matriks!reduksi elementer}\index{matriks reduksi elementer}
(atau cukup \definend{matriks elementer})\index{matriks elementer}%
\index{matriks!elementer}
dihasilkan dengan menerapkan satu operasi Gauss pada matriks identitas.
\begin{enumerate}
  \item \( I\grstep{k\rho_i}M_i(k) \) untuk \( k\neq 0 \)
  \item \( I\grstep{\rho_i\leftrightarrow\rho_j}P_{i,j} \) untuk
          \( i\neq j \)
  \item \( I\grstep{k\rho_i+\rho_j}C_{i,j}(k) \) untuk \( i\neq j \)
\end{enumerate}
%</df:ElementaryReductionMatrices>
\end{definition}

\begin{lemma}   \label{GrByMatMult}
%<*lm:GrByMatMult>
Perkalian matriks dapat melakukan reduksi Gauss.\index{operasi reduksi elementer!melalui perkalian matriks}\index{operasi baris elementer!melalui perkalian matriks}\index{Metode Gauss!melalui perkalian matriks}
\begin{enumerate}
  \item Jika \( H\grstep{k\rho_i}G \), maka \( M_i(k)H=G \).
  \item Jika \( H\grstep{\rho_i\leftrightarrow\rho_j}G \),
         maka \( P_{i,j}H=G \).
  \item Jika \( H\grstep{k\rho_i+\rho_j}G \), maka \( C_{i,j}(k)H=G \).
\end{enumerate}
%</lm:GrByMatMult>
\end{lemma}

\begin{proof}
Jelas.
\end{proof}

\begin{example}
Berikut adalah sistem pertama, dari bab pertama, yang kita
selesaikan dengan Metode Gauss.
\begin{equation*}
  \begin{linsys}{3}
             &   &      &   &3x_3  &=  &9  \\
        x_1  &+  &5x_2  &-  &2x_3  &=  &2  \\
   (1/3)x_1  &+  &2x_2  &   &      &=  &3 
  \end{linsys}
\end{equation*}
Kita dapat mereduksinya dengan perkalian matriks.
Tukar baris pertama dan ketiga,
\begin{equation*}
    \begin{mat}[r]
       0  &0  &1  \\
       0  &1  &0  \\
       1  &0  &0
    \end{mat}
    \begin{amat}[r]{3}
       0    &0    &3   &9   \\
       1    &5    &-2  &2   \\
      1/3   &2    &0   &3
    \end{amat}
  =
    \begin{amat}[r]{3}
      1/3   &2    &0   &3   \\
       1    &5    &-2  &2   \\
       0    &0    &3   &9
    \end{amat}
\end{equation*}
kalikan baris pertama dengan tiga,
\begin{equation*}
    \begin{mat}[r]
       3  &0  &0  \\
       0  &1  &0  \\
       0  &0  &1
    \end{mat}
    \begin{amat}[r]{3}
      1/3   &2    &0   &3   \\
       1    &5    &-2  &2   \\
       0    &0    &3   &9
    \end{amat}
  =
    \begin{amat}[r]{3}
       1    &6    &0   &9   \\
       1    &5    &-2  &2   \\
       0    &0    &3   &9
    \end{amat}
\end{equation*}
kemudian tambahkan \( -1 \) kali baris pertama ke baris kedua.
\begin{equation*}
    \begin{mat}[r]
       1  &0  &0  \\
      -1  &1  &0  \\
       0  &0  &1
    \end{mat}
    \begin{amat}[r]{3}
       1    &6    &0   &9   \\
       1    &5    &-2  &2   \\
       0    &0    &3   &9
    \end{amat}
  =
    \begin{amat}[r]{3}
       1    &6    &0   &9   \\
       0    &-1   &-2  &-7  \\
       0    &0    &3   &9
    \end{amat}
\end{equation*}
Sekarang substitusi balik akan memberikan solusinya.
\end{example}

\begin{example}
Reduksi Gauss--Jordan bekerja dengan cara yang sama.
Untuk matriks pada akhir contoh sebelumnya, pertama-tama ubah entri-entri utama
menjadi satu,
\begin{equation*}
    \begin{mat}[r]
       1  &0  &0  \\
       0  &-1 &0  \\
       0  &0  &1/3
    \end{mat}
    \begin{amat}[r]{3}
       1    &6    &0   &9   \\
       0    &-1   &-2  &-7  \\
       0    &0    &3   &9
    \end{amat}
  =
    \begin{amat}[r]{3}
       1    &6    &0   &9   \\
       0    &1    &2   &7   \\
       0    &0    &1   &3
    \end{amat}
\end{equation*}
kemudian nolkan kolom ketiga, lalu kolom kedua.
\begin{equation*}
    \begin{mat}[r]
       1  &-6 &0  \\
       0  &1  &0  \\
       0  &0  &1
    \end{mat}
    \begin{mat}[r]
       1  &0  &0  \\
       0  &1  &-2 \\
       0  &0  &1
    \end{mat}
    \begin{amat}[r]{3}
       1    &6    &0   &9   \\
       0    &1    &2   &7   \\
       0    &0    &1   &3
    \end{amat}
  =
    \begin{amat}[r]{3}
       1    &0    &0   &3   \\
       0    &1    &0   &1   \\
       0    &0    &1   &3
    \end{amat}
\end{equation*}
\end{example}

% We have observed the following result
% that we shall use in the next subsection.

\begin{corollary} \label{cor:ReducViaMatrices}
%<*co:ReducViaMatrices>
Untuk setiap matriks \( H \), terdapat matriks-matriks reduksi elementer
\( R_1 \), \ldots, \( R_r \) sedemikian sehingga
\( R_r\cdot R_{r-1}\cdots R_1\cdot H \)
berada dalam bentuk eselon tereduksi.
%</co:ReducViaMatrices>
\end{corollary}

Sampai saat ini kita telah mengambil sudut pandang bahwa objek utama yang kita kaji
adalah ruang vektor dan peta di antaranya, dan seolah-olah kita 
menggunakan matriks hanya demi kemudahan komputasi.
Subbagian ini menunjukkan bahwa ceritanya tidak sesederhana itu.

Memahami operasi matriks melalui pemahaman 
tentang mekanika penggabungan entri juga berguna. 
Dalam sisa buku ini kita akan tetap berfokus pada peta sebagai objek
utama, tetapi kita akan bersikap pragmatis\Dash jika sudut pandang matriks memberi
gagasan yang lebih jelas, kita akan menggunakannya.


\begin{exercises}
  \recommended \item
    Prediksikan hasil setiap perkalian dengan matriks permutasi,
    lalu periksa dengan mengerjakan perkaliannya.
    \begin{exparts*}
      \partsitem 
        $\begin{mat}
           0 &1 \\
           1 &0
        \end{mat}
        \begin{mat}
           1 &2 \\
           3 &4
        \end{mat}$
      \partsitem 
        $\begin{mat}
          1 &2  \\
          3 &4
        \end{mat}
        \begin{mat}
          0 &1 \\
          1 &0
        \end{mat}$
      \partsitem 
        $\begin{mat}
           1 &0 &0 \\
           0 &0 &1 \\
           0 &1 &0
        \end{mat}
        \begin{mat}
           1 &2 &3 \\
           4 &5 &6 \\
           7 &8 &9
        \end{mat}$
    \end{exparts*}
    \begin{answer}
      \begin{exparts}
        \partsitem  Ketika bekerja dari kiri, matriks $P_{1,2}$ menukar 
          baris pertama dan kedua.
          $\begin{mat}
             3 &4 \\
             1 &2
           \end{mat}$
        \partsitem 
          Ketika bekerja dari kanan, $P_{1,2}$ menukar kolom pertama dan kedua.
          $\begin{mat}
             2 &1 \\
             4 &3
           \end{mat}$
        \partsitem Dari kiri, $P_{2,3}$ menukar baris kedua dan ketiga.
          $\begin{mat}
             1 &2 &3 \\
             7 &8 &9 \\
             4 &5 &6    
           \end{mat}$
      \end{exparts}
    \end{answer}
  \recommended \item
    Prediksikan hasil setiap perkalian dengan matriks
    reduksi elementer, lalu periksa dengan mengerjakan perkaliannya.
    \begin{exparts*}
      \partsitem \( \begin{mat}[r]
                 3  &0  \\
                 0  &1
               \end{mat}
               \begin{mat}[r]
                 1  &2  \\
                 3  &4
               \end{mat} \)
      \partsitem \( \begin{mat}[r]
                 1  &0  \\
                 0  &2
               \end{mat}
               \begin{mat}[r]
                 1  &2  \\
                 3  &4
               \end{mat} \)
      \partsitem \( \begin{mat}[r]
                 1  &0  \\
                -2  &1
               \end{mat}
               \begin{mat}[r]
                 1  &2  \\
                 3  &4
               \end{mat} \)
      \partsitem \( \begin{mat}[r]
                 1  &2  \\
                 3  &4
               \end{mat}
               \begin{mat}[r]
                 1  &-1 \\
                 0  &1
               \end{mat} \)
      \partsitem \( \begin{mat}[r]
                 1  &2  \\
                 3  &4
               \end{mat}
               \begin{mat}[r]
                 0  &1  \\
                 1  &0
               \end{mat} \)
    \end{exparts*}
    \begin{answer}
      \begin{exparts}
        \partsitem Baris pertama matriks kedua dikalikan 
          dengan \( 3 \).
          \begin{equation*}
            \begin{mat}[r]
              3  &6  \\
              3  &4
            \end{mat}
          \end{equation*}
        \partsitem Baris kedua matriks kedua dikalikan dengan \( 2 \).
          \begin{equation*}
            \begin{mat}[r]
              1  &2  \\
              6  &8
            \end{mat}
          \end{equation*}
        \partsitem Matriks kedua dikenai operasi kombinasi
          yang mengganti baris kedua dengan \( -2 \) kali baris pertama ditambah
          baris kedua.
          \begin{equation*}
            \begin{mat}[r]
              1  &2  \\
              1  &0
            \end{mat}
          \end{equation*}
        \partsitem Matriks pertama dikenai operasi kolom:~ganti
          kolom
          kedua dengan $-1$ kali kolom pertama ditambah kolom
          kedua.
          \begin{equation*}
            \begin{mat}[r]
              1  &1  \\
              3  &1
            \end{mat}
          \end{equation*}
        \partsitem Kolom-kolom matriks pertama ditukar.
          \begin{equation*}
            \begin{mat}[r]
              2  &1  \\
              4  &3
            \end{mat}
          \end{equation*}
      \end{exparts}   
    \end{answer}
  \item
    Prediksikan hasil setiap perkalian dengan
    matriks diagonal, lalu periksa dengan mengerjakan perkaliannya.
    \begin{exparts*}
      \partsitem \( \begin{mat}[r]
                 -3  &0  \\
                 0  &0
               \end{mat}
               \begin{mat}[r]
                 1  &2  \\
                 3  &4
               \end{mat} \)
      \partsitem \( \begin{mat}[r]
                 4  &0  \\
                 0  &2
               \end{mat}
               \begin{mat}[r]
                 1  &2  \\
                 3  &4
               \end{mat} \)
    \end{exparts*}
    \begin{answer}
      \begin{exparts}
        \partsitem Baris pertama matriks kedua dikalikan 
          dengan \( -3 \), dan
          baris keduanya dikalikan dengan \( 0 \).
          \begin{equation*}
            \begin{mat}[r]
              -3  &-6  \\
               0  &0
            \end{mat}
          \end{equation*}
        \partsitem Baris pertama matriks kedua dikalikan 
          dengan \( 4 \), dan
          baris keduanya dikalikan dengan \( 2 \).
          \begin{equation*}
            \begin{mat}[r]
              4  &8  \\
              6  &8
            \end{mat}
          \end{equation*}
      \end{exparts}   
    \end{answer}
  \item Buatlah masing-masing matriks berikut.
    \begin{exparts}
      \partsitem matriks $\nbyn{3}$ yang, ketika bekerja dari kiri,
         menukar baris pertama dan kedua
      \partsitem matriks $\nbyn{2}$ yang, ketika bekerja dari kanan,
         menukar kolom pertama dan kedua
    \end{exparts}
    \begin{answer}
      \begin{exparts}
        \partsitem Matriks ini menukar baris pertama dan kedua.
          \begin{equation*}
            \begin{mat}
              0  &1 &0 \\
              1  &0 &0 \\
              0  &0 &1
            \end{mat}
            \begin{mat}
              a &b &c \\
              d &e &f \\
              g &h &i
            \end{mat}
            =
            \begin{mat}
              d &e &f \\
              a &b &c \\
              g &h &i
            \end{mat}         
          \end{equation*}
        \partsitem Matriks ini menukar kolom pertama dan kedua.
          \begin{equation*}
            \begin{mat}
              a  &b  \\
              c  &d
            \end{mat}
            \begin{mat}
              0  &1  \\
              1  &0   
            \end{mat}
            =
            \begin{mat}
              b  &a  \\
              d  &c
            \end{mat}
          \end{equation*}
      \end{exparts}
    \end{answer}
  \recommended \item  Tunjukkan cara menggunakan perkalian matriks untuk membawa matriks ini
    ke bentuk eselon.
    \begin{equation*}
      \begin{mat}
        1 &2 &1  &0   \\
        2 &3 &1  &-1  \\
        7 &11 &4 &-3
      \end{mat}
    \end{equation*}
    \begin{answer}
      Kalikan dengan $C_{1,2}(-2)$, kemudian dengan~$C_{1,3}(-7)$, lalu dengan~$C_{2,3}(-3)$,
      dengan memperhatikan urutan dari kanan ke kiri.
      \begin{multline*}
        \begin{mat}
          1  &0 &0 \\
          0  &1 &0 \\
          0  &-3 &1    
        \end{mat}
        \begin{mat}
          1  &0 &0 \\
          0  &1 &0 \\
          -7 &0 &1    
        \end{mat}
        \begin{mat}
          1 &0 &0 \\
         -2 &1 &0 \\
          0 &0 &1
        \end{mat}
          \begin{mat}
            1 &2 &1  &0   \\
            2 &3 &1  &-1  \\
            7 &11 &4 &-3
          \end{mat}                 \\
        =
          \begin{mat}
            1 &2 &1  &0   \\
            0 &-1 &-1  &-1  \\
            0 &0  &0 &0
          \end{mat}
      \end{multline*}     
    \end{answer}
  \item 
    Carilah hasil kali matriks ini dengan transposenya.
    \begin{equation*}
      \begin{mat}
        \cos\theta  &-\sin\theta  \\
        \sin\theta  &\cos\theta
      \end{mat}
    \end{equation*}
    \begin{answer}
      Hasil kalinya adalah matriks identitas (ingat bahwa
      $\cos^2\theta+\sin^2\theta =1$).
      Penjelasannya adalah bahwa, terhadap basis standar, matriks yang diberikan
      merepresentasikan rotasi sebesar \( \theta \) radian di \( \Re^2 \),
      sedangkan transposenya merepresentasikan rotasi sebesar \( -\theta \)
      radian.
      Keduanya saling meniadakan.  
     \end{answer}
  \item 
    Kebutuhan untuk mengambil kombinasi linear baris dan kolom dalam tabel 
    bilangan sering muncul dalam praktik.
    Sebagai contoh, berikut adalah peta sebagian wilayah Vermont dan New York.
    %\typeout{Ignore the overfull box here; it is bogus.}
    \begin{center}
      \parbox{2in}{Sebagian karena adanya Danau Champlain,
                      beberapa pasangan kota tidak terhubung langsung oleh jalan.
                      Sebagai contoh, tidak ada cara untuk pergi dari
                      Winooski ke Grand Isle tanpa melewati
                      Colchester.
                      (Untuk menyederhanakan graf, banyak jalan dan kota lain
                      dihilangkan.
                      Jarak dari atas ke bawah peta ini
                      sekitar empat puluh mil.)}
      \quad
      \parbox{2in}{\includegraphics{map/mp/ch3.99}}
    \end{center}
    \begin{exparts}
      \partsitem \definend{Matriks ketetanggaan}\index{matriks ketetanggaan}%
        \index{matriks!ketetanggaan} 
        suatu peta adalah matriks persegi
        yang entri \( i,j \)-nya menyatakan banyaknya jalan dari kota \( i \)
        ke kota \( j \) (semua entri $(i,i)$ bernilai~$0$).
        Buatlah matriks ketetanggaan peta ini dengan mengurutkan kota-kotanya
        secara alfabetis.
      \partsitem Suatu matriks disebut 
        \definend{simetris}\index{matriks!simetris} 
        jika sama dengan transposenya.
        Tunjukkan bahwa matriks ketetanggaan bersifat simetris.
        (Semua jalan ini dapat dilalui dalam dua arah.
        Vermont tidak memiliki banyak jalan satu arah.)
      \partsitem Apa makna kuadrat dari 
        matriks ketetanggaan?
        Bagaimana dengan pangkat tiganya?
    \end{exparts}
    \begin{answer}
      \begin{exparts}
        \partsitem Matriks ketetanggaannya adalah sebagai berikut
          (misalnya, baris pertama menunjukkan bahwa hanya ada satu jalan
          yang terhubung dengan Burlington, yaitu jalan menuju Winooski).
          \begin{equation*}
            \begin{mat}[r]
              0  &0  &0  &0  &1  \\
              0  &0  &1  &1  &1  \\
              0  &1  &0  &1  &0  \\
              0  &1  &1  &0  &0  \\
              1  &1  &0  &0  &0
            \end{mat}
          \end{equation*}
        \partsitem Karena jalan-jalan ini dapat dilalui dalam dua arah, 
          setiap jalan yang menghubungkan kota~\( i \)
          ke kota~\( j \) juga memberikan hubungan antara kota~\( j \) dan 
          kota~\( i \).
        \partsitem Kuadrat matriks ketetanggaan menunjukkan bagaimana kota-kota
          terhubung
          melalui perjalanan yang melibatkan dua jalan.
      \end{exparts}  
   \end{answer}
  \recommended \item
    Tabel ini memberikan jumlah jam untuk setiap 
    jenis pekerjaan yang dilakukan tiap pekerja, beserta tarif upahnya.
    Gunakan matriks untuk menghitung upah yang harus dibayarkan.
    \begin{center}
      \begin{tabular}[t]{l|cc}
        \multicolumn{1}{c}{\ }  
         &\multicolumn{1}{c}{\textit{reguler}} 
         &\multicolumn{1}{c}{\textit{lembur}}  \\
        \cline{2-3}
        Alan      &40        &12        \\
        Betty     &35        &6         \\  
        Catherine &40        &18         \\  
        Donald    &28        &0         %\\  \cline{2-3}
      \end{tabular}
      \qquad
      \begin{tabular}[t]{l|c}
        \multicolumn{1}{c}{\ }   &\multicolumn{1}{c}{\textit{upah}}   \\
        \cline{2-2}
        reguler   &$\$ 25.00$  \\
        lembur  &$\$ 45.00$  %\\  \cline{2-2}
      \end{tabular}
    \end{center}
    \textit{Catatan.}
    Ini menggambarkan bahwa dalam praktik
    kita sering ingin menghitung kombinasi linear baris dan
    kolom dalam konteks yang sama sekali tidak membuat kita tertarik pada
    peta linear terkait.
    \begin{answer}
      Upah yang harus dibayarkan kepada setiap orang muncul dalam hasil kali matriks kedua
      susunan tersebut. 
    \end{answer}
  \item Nyatakan matriks nonsingular ini sebagai hasil kali matriks-matriks reduksi
  elementer.
  \begin{equation*}
    T=\begin{mat}
      1 &2  &0 \\
      2 &-1 &0 \\
      3 &1 &2
    \end{mat}
  \end{equation*}
  \begin{answer}
    Reduksi Gauss--Jordannya rutin.
    \begin{equation*}
        \begin{mat}
          1 &2  &0 \\
          2 &-1 &0 \\
          3 &1 &2
        \end{mat}
      \grstep[-3\rho_1+\rho_3]{-2\rho_1+\rho_2}
      \grstep{-\rho_2+\rho_3}
      \grstep[(1/2)\rho_3]{-(1/5)\rho_2}
      \grstep{-2\rho_2+\rho_1}
      \begin{mat}
        1 &0 &0 \\
        0 &1 &0 \\
        0 &0 &1
      \end{mat}
    \end{equation*}
    Jadi kita mengetahui matriks-matriks reduksi elementer $R_1,\ldots,R_6$ sedemikian sehingga
    $R_6\cdot R_5\cdots R_1\cdot T=I$.
    Pindahkan matriks-matriks itu ke ruas lain dengan memperhatikan urutannya.
    Sebagai contoh, pertama kalikan kedua ruas dari kiri dengan $R_6^{-1}$
    untuk memperoleh $(R_6^{-1}R_6)\cdot R_5\cdots R_1\cdot T=R_6^{-1}I$, 
    yang menyederhana menjadi $R_5\cdots R_1\cdot T=R_6^{-1}$, dan seterusnya.
    \begin{align*}
      T=
      &\begin{mat}
        1 &0 &0 \\
        -2 &1 &0 \\
        0 &0 &1 
      \end{mat}^{-1}
      \begin{mat}
        1 &0 &0 \\
        0 &1 &0 \\
        -3 &0 &1 
      \end{mat}^{-1}
      \begin{mat}
        1 &0 &0 \\
        0 &1 &0 \\
        0 &-1 &1 
      \end{mat}^{-1}          \\
      &\qquad\cdot
      \begin{mat}
        1 &0 &0 \\
        0 &-1/5 &0 \\
        0 &0 &1 
      \end{mat}^{-1}
      \begin{mat}
        1 &0 &0 \\
        0 &1 &0 \\
        0 &0 &1/2 
      \end{mat}^{-1}
      \begin{mat}
        1 &-2 &0 \\
        0 &1 &0 \\
        0 &0 &1 
      \end{mat}^{-1}             
    \end{align*}
    Mengambil invers matriks reduksi elementer itu mudah.
    Sebagai contoh, untuk membatalkan penambahan
    $k$ kali baris~$i$ ke baris~$j$, ambil $-k$ kali
    baris~$i$ dan tambahkan ke baris~$j$.
    Dengan demikian, 
    Lema~\ref{GrByMatMult} menyatakan bahwa $C_{i,j}(k)^{-1}=C_{i,j}(-k)$.
    \begin{equation*}
      T=
      \begin{mat}
        1 &0 &0 \\
        2 &1 &0 \\
        0 &0 &1
      \end{mat}
      \begin{mat}
        1 &0 &0 \\
        0 &1 &0 \\
        3 &0 &1
      \end{mat}
      \begin{mat}
        1 &0 &0 \\
        0 &1 &0 \\
        0 &1 &1
      \end{mat}
      \begin{mat}
        1 &0 &0 \\
        0 &-5 &0 \\
        0 &0 &1
      \end{mat}
      \begin{mat}
        1 &0 &0 \\
        0 &1 &0 \\
        0 &0 &2
      \end{mat}
      \begin{mat}
        1 &2 &0 \\
        0 &1 &0 \\
        0 &0 &1
      \end{mat}
    \end{equation*}    
    \end{answer}
  \item 
    Nyatakan
    \begin{equation*}
      \begin{mat}[r]
        1   &0  \\
        -3  &3
      \end{mat}
    \end{equation*}
    sebagai hasil kali dua matriks reduksi elementer.
    \begin{answer}
      Salah satu cara menghasilkan matriks ini dari identitas adalah menggunakan 
      operasi kolom:
      pertama kalikan kolom kedua dengan tiga, kemudian tambahkan 
      negatif kolom kedua yang dihasilkan ke kolom pertama.
      \begin{equation*}
        \begin{mat}[r]
          1  &0  \\
          0  &1
        \end{mat}
        \grstep{}
        \begin{mat}[r]
          1  &0  \\
          0  &3
        \end{mat}
        \grstep{}
        \begin{mat}[r]
          1  &0  \\
          -3  &3
        \end{mat}
      \end{equation*}
      Berbeda dengan operasi baris, operasi kolom dituliskan dari
      kiri ke kanan, sehingga hasil kali matriks ini menyatakan
      pelaksanaan kedua operasi di atas.
      \begin{equation*}
        \begin{mat}[r]
          1  &0  \\
          0  &3
        \end{mat}
        \begin{mat}[r]
          1  &0  \\
         -1  &1
        \end{mat}
      \end{equation*}
      \textit{Catatan.}
      Sebagai alternatif, kita dapat memperoleh matriks yang diminta dengan operasi baris.
      Mulai dari identitas, pertama tambahkan negatif baris pertama 
      ke baris kedua, kemudian kalikan baris kedua dengan tiga.
      Karena kita menuliskan operasi baris berturut-turut sebagai hasil kali matriks dari
      kanan ke kiri, pelaksanaan kedua operasi baris ini dinyatakan oleh:~hasil kali 
      matriks yang sama.    
    \end{answer}
  \recommended \item 
    Buktikan bahwa matriks-matriks diagonal membentuk subruang dari
    \( \matspace_{\nbyn{n}} \).
    Berapakah dimensinya?
    \begin{answer}
      Himpunan matriks diagonal tidak kosong karena matriks nol
      bersifat diagonal.
      Jelas bahwa himpunan ini tertutup terhadap perkalian skalar dan penjumlahan.
      Oleh karena itu, himpunan ini merupakan subruang.
      Dimensinya adalah \( n \); berikut adalah salah satu basisnya.
      \begin{equation*}
        \set{\begin{mat}
               1  &0  &\ldots    \\
               0  &0             \\
                  &   &\ddots    \\
               0  &0  &      &0
             \end{mat},\ldots,
             \begin{mat}
               0  &0  &\ldots    \\
               0  &0             \\
                  &   &\ddots    \\
               0  &0  &      &1
             \end{mat}  }
      \end{equation*} 
    \end{answer}
  \item 
    Apakah matriks identitas merepresentasikan peta identitas jika kedua basisnya
    berbeda?
    \begin{answer}
      Tidak.
      Dalam \( \polyspace_1 \), terhadap basis-basis berbeda
      \( B=\sequence{1,x} \) dan \( D=\sequence{1+x,1-x} \),
      transformasi identitas direpresentasikan oleh matriks berikut.
      \begin{equation*}
         \rep{\text{id}}{B,D}=
         \begin{mat}[r]
           1/2  &1/2  \\
           1/2  &-1/2
         \end{mat}_{B,D}
      \end{equation*}   
    \end{answer}
  \item 
    Tunjukkan bahwa setiap kelipatan identitas komutatif dengan setiap
    matriks persegi.
    Adakah matriks lain yang komutatif dengan semua matriks persegi?
    \begin{answer}
      Untuk setiap skalar \( r \) dan matriks persegi \( H \), kita memiliki
      \( (rI)H=r(IH)=rH=r(HI)=(Hr)I=H(rI) \).

      Tidak ada matriks lain semacam itu; berikut argumen untuk matriks 
      $\nbyn{2}$ yang mudah diperluas ke $\nbyn{n}$.
      Jika suatu matriks komutatif dengan semua matriks lain, matriks itu komutatif dengan
      matriks satuan berikut.
      \begin{equation*}
        \begin{mat}
          0  &a  \\
          0  &c 
        \end{mat}
        =\begin{mat}
          a  &b  \\
          c  &d   
        \end{mat}
        \begin{mat}[r]
          0  &1  \\
          0  &0
        \end{mat}
        =\begin{mat}[r]
          0  &1  \\
          0  &0
        \end{mat}
        \begin{mat}
          a  &b  \\
          c  &d   
        \end{mat}
        =\begin{mat}
          c  &d  \\
          0  &0 
        \end{mat}
      \end{equation*}
      Dari sini pertama-tama kita simpulkan bahwa entri kiri atas~$a$ 
      harus sama dengan entri kanan bawah~$d$.
      Kita juga menyimpulkan bahwa entri kiri bawah~$c$ bernilai nol.
      Argumen untuk entri kanan atas~$b$ serupa.
    \end{answer}
  \item 
    Buktikan atau sangkal: matriks-matriks nonsingular bersifat komutatif.
    \begin{answer}
      Pernyataan ini salah; kedua matriks berikut tidak komutatif.
      \begin{equation*}
         \begin{mat}[r]
           1  &2  \\
           3  &4
         \end{mat}
         \qquad
         \begin{mat}[r]
           5  &6  \\
           7  &8
         \end{mat}
      \end{equation*}  
    \end{answer}
  \recommended \item \label{exer:PermTimesTransEqId}
    Tunjukkan bahwa hasil kali suatu matriks permutasi dengan transposenya
    adalah matriks identitas.
    \begin{answer}
      Matriks permutasi memiliki tepat satu bilangan satu pada setiap baris dan kolom,
      dan semua entri lainnya bernilai nol.
      Tetapkan sebuah matriks semacam itu.
      Misalkan bilangan satu pada baris ke-\( i \) berada di kolom ke-\( j \).
      Maka tidak ada baris lain yang memiliki bilangan satu pada kolom ke-\( j \); setiap
      baris lain memiliki nol pada kolom ke-\( j \).
      Jadi hasil kali titik baris ke-\( i \) dengan sebarang baris lain bernilai nol.

      Baris ke-\( i \) hasil kali tersusun atas hasil kali titik antara
      baris ke-\( i \) matriks dan kolom-kolom transposenya.
      Berdasarkan paragraf sebelumnya, semua hasil kali titik tersebut bernilai nol kecuali
      yang ke-\( i \), yang bernilai satu.  
    \end{answer}
  \item 
    Tunjukkan bahwa jika baris pertama dan kedua \( G \) sama, maka
    baris pertama dan kedua \( GH \) juga sama.
    Perumumlah.
    \begin{answer}
      Perumumannya adalah mengganti baris pertama dan kedua dengan
      baris ke-$i_1$ dan ke-$i_2$.
      Baris~$i$ dari \( GH \) tersusun atas hasil kali titik antara
      baris~$i$ dari \( G \) dan kolom-kolom \( H \).
      Jadi, jika baris \( i_1 \) dan \( i_2 \) dari \( G \) sama, maka
      baris \( i_1 \) dan \( i_2 \) dari \( GH \) juga sama.  
    \end{answer}
  \item 
    Deskripsikan hasil kali dua matriks diagonal.
    \begin{answer}
      Jika hasil kali dua matriks diagonal terdefinisi\Dash yaitu jika 
      keduanya berukuran $\nbyn{n}$\Dash maka 
      entri diagonal hasil kalinya adalah hasil kali entri-entri
      diagonalnya:~jika \( G,H \) adalah matriks diagonal berukuran sama,
      semua entri \( GH \) bernilai nol kecuali entri \( i,i \)-nya, yaitu
      \( g_{i,i}h_{i,i} \).
    \end{answer}
  \recommended \item 
    Tunjukkan bahwa jika \( G \) memiliki satu baris nol, maka \( GH \)
    (jika terdefinisi) memiliki satu baris nol.
    Apakah hal itu berlaku untuk kolom?
    \begin{answer}
      Baris ke-\( i \) dari \( GH \) tersusun atas hasil kali titik antara
      baris ke-\( i \) dari \( G \) dan kolom-kolom \( H \).
      Hasil kali titik baris nol dengan suatu kolom bernilai nol.

      Hal itu berlaku untuk kolom jika dinyatakan dengan tepat:~jika \( H \) memiliki satu kolom
      nol, maka \( GH \) (jika terdefinisi) memiliki satu kolom nol.
      Pembuktiannya mudah.  
    \end{answer}
  \item 
    Tunjukkan bahwa himpunan matriks satuan membentuk basis bagi 
    \( \matspace_{\nbym{n}{m}} \).
    \begin{answer}
      Mungkin cara termudah adalah menunjukkan bahwa setiap matriks \( \nbym{n}{m} \)
      merupakan kombinasi linear matriks-matriks satuan dengan tepat satu cara:
      \begin{equation*}
        c_1\begin{mat}
             1  &0  &\ldots  \\
             0  &0           \\
             \vdots
           \end{mat}
        +\dots+
        c_{n,m}\begin{mat}
             0  &0      &\ldots  \\
             \vdots              \\
             0  &\ldots &       &1
           \end{mat}
        =\begin{mat}
           a_{1,1} &a_{1,2} &\ldots          \\
           \vdots                            \\
           a_{n,1} &\ldots  &      &a_{n,m} 
         \end{mat}
      \end{equation*}
      memiliki solusi tunggal \( c_1=a_{1,1} \), \( c_2=a_{1,2} \),
      dan seterusnya.  
    \end{answer}
  \item  
    Carilah rumus untuk pangkat ke-$n$ matriks berikut.
    \begin{equation*}
      \begin{mat}[r]
        1  &1  \\
        1  &0
      \end{mat}
    \end{equation*}
    \begin{answer}
      Namai matriks tersebut \( F \).
      Kita memperoleh
      \begin{equation*}
        F^2=\begin{mat}[r]
          2  &1  \\
          1  &1
        \end{mat}
        \quad
        F^3=\begin{mat}[r]
          3  &2  \\
          2  &1
        \end{mat}
        \quad
        F^4=\begin{mat}[r]
          5  &3  \\
          3  &2
        \end{mat}
      \end{equation*}
      Secara umum, 
      \begin{equation*}
        F^n=\begin{mat}
          f_{n+1} &f_n  \\
          f_n     &f_{n-1}
        \end{mat}
      \end{equation*}
      dengan \( f_i \) bilangan Fibonacci ke-\( i \),
      \( f_i=f_{i-1}+f_{i-2} \) dan \( f_0=0 \), \( f_1=1 \),
      yang kita verifikasi dengan induksi berdasarkan persamaan berikut.
      \begin{equation*}
        \begin{mat}
          f_{i-1}  &f_{i-2} \\
          f_{i-2}  &f_{i-3}
        \end{mat}
        \begin{mat}[r]
          1  &1  \\
          1  &0
        \end{mat}
        =\begin{mat}
           f_i     &f_{i-1}  \\
           f_{i-1} &f_{i-2}
        \end{mat}
      \end{equation*}
    \end{answer}
  \recommended \item
   \definend{Teras}\index{teras}\index{matriks!teras} 
   suatu matriks persegi adalah jumlah
   entri pada diagonalnya (maknanya muncul dalam
   Bab Lima).
   Tunjukkan bahwa \( \trace (GH)=\trace (HG)  \).
   \begin{answer}
     \textit{Bab Lima memberikan alasan yang tidak terlalu bersifat komputasional\Dash
     teras suatu matriks adalah koefisien kedua dalam polinomial
     karakteristiknya\Dash tetapi untuk saat ini kita dapat menggunakan indeks.}
     Kita memiliki
     \begin{align*}
       \trace (GH)
       &=(g_{1,1}h_{1,1}+g_{1,2}h_{2,1}+\dots+g_{1,n}h_{n,1})  \\
       &\text{}\quad+(g_{2,1}h_{1,2}+g_{2,2}h_{2,2}+\dots+g_{2,n}h_{n,2}) \\
       &\text{}\quad
        +\cdots+(g_{n,1}h_{1,n}+g_{n,2}h_{2,n}+\dots+g_{n,n}h_{n,n})
     \end{align*}
     sedangkan
     \begin{align*}
       \trace (HG)
       &=(h_{1,1}g_{1,1}+h_{1,2}g_{2,1}+\dots+h_{1,n}g_{n,1})  \\
       &\text{}\quad+(h_{2,1}g_{1,2}+h_{2,2}g_{2,2}+\dots+h_{2,n}g_{n,2}) \\
       &\text{}\quad
         +\cdots+(h_{n,1}g_{1,n}+h_{n,2}g_{2,n}+\dots+h_{n,n}g_{n,n})
     \end{align*}
     dan keduanya sama.  
    \end{answer}
  \item
    Suatu matriks persegi disebut 
    \definend{segitiga atas}\index{matriks segitiga}%
    \index{matriks!segitiga} 
    jika semua entri taknolnya
    terletak di atas atau pada diagonal.
    Tunjukkan bahwa hasil kali dua matriks segitiga atas juga segitiga atas.
    Apakah hal ini berlaku pula untuk matriks segitiga bawah?
    \begin{answer}
      Suatu matriks segitiga atas jika dan hanya jika entri $i,j$-nya bernilai nol 
      setiap kali \( i>j \).
      Jadi, jika \( G,H \) segitiga atas, maka \( h_{i,j} \) dan
      \( g_{i,j} \) bernilai nol ketika \( i>j \).
      Suatu entri dalam hasil kali
      \( p_{i,j}=g_{i,1}h_{1,j}+\dots+g_{i,n}h_{n,j} \)
      bernilai nol kecuali setidaknya salah satu sukunya taknol, yakni kecuali
      untuk setidaknya salah satu suku
      \( g_{i,r}h_{r,j} \), berlaku sekaligus \( i\leq r \) dan \( r\leq j \).
      Tentu saja, jika \( i>j \), hal ini mustahil terjadi sehingga hasil kali dua
      matriks segitiga atas juga segitiga atas.
      (Argumen serupa berlaku untuk matriks segitiga bawah.) 
   \end{answer}
 \item 
   Suatu matriks persegi disebut \definend{matriks Markov}\index{matriks!Markov} 
   jika setiap entrinya berada di antara nol
   dan satu serta jumlah entri pada setiap barisnya adalah satu.
   Buktikan bahwa hasil kali matriks-matriks Markov juga merupakan matriks Markov.
   \begin{answer}
     Jumlah entri sepanjang baris ke-\( i \) hasil kali adalah sebagai berikut.
     \begin{align*}
       p_{i,1}+\cdots+p_{i,n}
       &=(h_{i,1}g_{1,1}+h_{i,2}g_{2,1}+\dots+h_{i,n}g_{n,1})  \\
       &\text{}\quad+(h_{i,1}g_{1,2}+h_{i,2}g_{2,2}+\dots+h_{i,n}g_{n,2}) \\
       &\text{}\quad
        +\dots+(h_{i,1}g_{1,n}+h_{i,2}g_{2,n}+\dots+h_{i,n}g_{n,n})  \\
       &=h_{i,1}(g_{1,1}+g_{1,2}+\dots+g_{1,n})  \\
       &\text{}\quad+h_{i,2}(g_{2,1}+g_{2,2}+\dots+g_{2,n}) \\
       &\text{}\quad
        +\dots+h_{i,n}(g_{n,1}+g_{n,2}+\dots+g_{n,n})  \\
       &=h_{i,1}\cdot 1+\dots+h_{i,n}\cdot 1  \\
       &=1
     \end{align*}  
    \end{answer}
  \item
    Berikan contoh dua matriks dengan peringkat dan ukuran yang sama,
    tetapi kuadratnya memiliki peringkat berbeda.
    \begin{answer}
      Tetapkan basis $B=\sequence{\vec{\beta}_1, \vec{\beta}_2}$.
      Matriks-matriks yang merepresentasikan peta yang mengirim
      \begin{equation*}
        \vec{\beta}_1 \mapsunder{h} \vec{\beta}_1
        \quad
        \vec{\beta}_2 \mapsunder{h} \zero
      \end{equation*}
      dan
      \begin{equation*}
        \vec{\beta}_1 \mapsunder{g} \vec{\beta}_2
        \quad
        \vec{\beta}_2 \mapsunder{g} \zero
      \end{equation*}
      memenuhi.  
      Sebagai contoh, jika kita menggunakan basis standar, kedua peta di atas
      menghasilkan matriks-matriks berikut.
      \begin{equation*}
        H=
        \begin{mat}
          1 &0 \\
          0 &0
        \end{mat}
        \qquad
        G=
        \begin{mat}
          0 &1 \\
          0 &0
        \end{mat}
      \end{equation*}
      Perhatikan bahwa $H^2$ berperingkat~$1$, sedangkan $G^2$ berperingkat~$0$.
    \end{answer}
  % 2014-Dec-19 This exercise and answer are good; I trimmed to make pages fit better.
  % \item 
  %   Combine the two generalizations of the identity matrix,
  %   the one allowing entries to be other than ones, and the one allowing the
  %   single one in each row and column to be off the diagonal.
  %   What is the action of this type of matrix?
  %   \begin{answer}
  %     The combination is to have all entries of the matrix be zero
  %     except for one (possibly) nonzero entry in each row and column.
  %     We can write such a matrix as the product of a permutation matrix and
  %     a diagonal matrix, e.g.,
  %     \begin{equation*}
  %       \begin{mat}[r]
  %         0  &4  &0  \\
  %         2  &0  &0  \\
  %         0  &0  &-5
  %       \end{mat}
  %       =\begin{mat}[r]
  %         0  &1  &0  \\
  %         1  &0  &0  \\
  %         0  &0  &1
  %       \end{mat}
  %       \begin{mat}[r]
  %         4  &0  &0  \\
  %         0  &2  &0  \\
  %         0  &0  &-5
  %       \end{mat}
  %     \end{equation*}
  %     and its action is thus to rescale the rows and permute them.
  %   \end{answer}
  \item 
    Perkalian matriks sering dilakukan dengan komputer.
    Para peneliti yang berusaha memahami dan meningkatkan kinerjanya
    menghitung banyaknya operasi yang diperlukan.
    \begin{exparts}
      \partsitem \nearbydefinition{def:MatMult} memberikan 
        % a formula for 
        % entry~$i,j$.
        % \begin{equation*}
          $p_{i,j}
          =
          g_{i,1}h_{1,j}+g_{i,2}h_{2,j}+\dots+g_{i,r}h_{r,j}$.
        % \end{equation*}
        Berapa banyak perkalian bilangan real dalam ekspresi tersebut?
        Dengan menggunakannya, berapa banyak yang diperlukan untuk hasil kali matriks
        \( \nbym{m}{r} \) dan matriks \( \nbym{r}{n} \)?
      \partsitem 
        Perkalian matriks bersifat asosiatif, sehingga ketika menghitung
        $H_1H_2H_3H_4$, kita memperoleh 
        jawaban yang sama di mana pun tanda kurung ditempatkan.
        Namun, biaya dalam banyaknya perkalian berbeda-beda.
        Carilah pengelompokan yang memerlukan paling sedikit perkalian bilangan real
        untuk menghitung hasil kali matriks yang terdiri atas
        matriks \( \nbym{5}{10} \),
        matriks \( \nbym{10}{20} \),
        matriks \( \nbym{20}{5} \), dan
        matriks \( \nbym{5}{1} \).
        Gunakan rumus yang sama seperti pada bagian sebelumnya.
      \partsitem \textit{(Sangat sulit.)}
        Carilah cara mengalikan dua matriks \( \nbyn{2} \) dengan hanya tujuh
        perkalian, bukan delapan seperti yang disarankan pendekatan sebelumnya.
    \end{exparts}
    \begin{answer}
      \begin{exparts}
        \partsitem Setiap entri \(p_{i,j}=g_{i,1}h_{1,j}+\dots+g_{i,r}h_{r,j}  \)
          memerlukan \( r \) perkalian, dan terdapat \( m\cdot n \) entri.
          Jadi terdapat \( m\cdot n\cdot r \) perkalian.
        \partsitem Misalkan \( H_1 \) berukuran \( \nbym{5}{10} \),
           \( H_2 \) berukuran \( \nbym{10}{20} \),
           \( H_3 \) berukuran \( \nbym{20}{5} \), dan
           \( H_4 \) berukuran \( \nbym{5}{1} \).
           Maka
           \begin{center}
             \begin{tabular}{l|l}
               \multicolumn{1}{c}{\textit{pengelompokan ini}}  
                &\multicolumn{1}{c}{\textit{menggunakan perkalian sebanyak ini}}\\ 
               \hline
               \( ((H_1H_2)H_3)H_4 \)     &\( 1000+500+25=1525 \)  \\
               \( (H_1(H_2H_3))H_4 \)     &\( 1000+250+25=1275 \)  \\
               \( (H_1H_2)(H_3H_4) \)     &\( 1000+100+100=1200 \)  \\
               \( H_1(H_2(H_3H_4)) \)     &\( 100+200+50=350    \)  \\
               \( H_1((H_2H_3)H_4) \)     &\( 1000+50+50=1100    \)
             \end{tabular}
           \end{center}
           memperlihatkan pengelompokan yang paling murah.
         \partsitem Berikut adalah penyempurnaan oleh S.~Winograd atas rumus dari
           V.~Strassen:
           misalkan \( w=aA-(a-c-d)(A-C+D) \), maka
           \begin{equation*}
               \begin{mat}
                 a  &b  \\
                 c  &d
               \end{mat}
               \begin{mat}
                 A  &C  \\
                 B  &D
                 \end{mat}    \\
               =
               \begin{mat}
                 \alpha  &\beta  \\
                 \gamma  &\delta
               \end{mat}
               \end{equation*}
             dengan $\alpha=aA+bB$,
             $\beta=w+(c+d)(C-A)+(a+b-c-d)D$,
             $\gamma=w+(a-c)(D-C)-d(A-B-C+D)$, dan
             $\delta=w+(a-c)(D-C)+(c+d)(C-A)$.
               % \begin{mat}
               %    aA+bB
               %    &w+(c+d)(C-A)+(a+b-c-d)D \\
               %    w+(a-c)(D-C)-d(A-B-C+D)
               %    &w+(a-c)(D-C)+(c+d)(C-A)
               % \end{mat}
           Cara ini memerlukan tujuh perkalian dan lima belas penjumlahan (simpan
           hasil-hasil antaranya).
      \end{exparts}  
    \end{answer}
  \puzzle \item  
    \cite{Putnam90A5}
    Jika \( A \) dan \( B \) adalah matriks persegi berukuran
    sama sedemikian sehingga \( ABAB=0 \), apakah mesti \( BABA=0 \)?
    \begin{answer}
      \answerasgiven 
      Tidak.
      Misalkan \( A \) dan \( B \), terhadap basis standar, merepresentasikan
      transformasi-transformasi \( \Re^3 \) berikut.
      \begin{equation*}
        \colvec{x \\ y \\ z}\mapsunder{a}\colvec{x \\ y \\ 0}
        \qquad
        \colvec{x \\ y \\ z}\mapsunder{b}\colvec{0 \\ x \\ y}
      \end{equation*}
      Perhatikan bahwa
      \begin{equation*}
        \colvec{x \\ y \\ z}\mapsunder{abab}\colvec[r]{0 \\ 0 \\ 0}
        \quad\text{tetapi}\quad
        \colvec{x \\ y \\ z}\mapsunder{baba}\colvec{0 \\ 0 \\ x}.
      \end{equation*} 
    \end{answer}
  \item  
    \cite{Monthly66p1114}
    Buktikan keempat pernyataan berikut untuk memperoleh
    bukti alternatif bahwa peringkat kolom sama dengan peringkat baris.
    \begin{exparts}
      \partsitem \( \vec{y}\cdot\vec{y}=0 \) jika dan hanya jika \( \vec{y}=\zero \).
      \partsitem \( A\vec{x}=\zero \) jika dan hanya jika \( \trans{A}A\vec{x}=\zero \).
      \partsitem \( \dim(\rangespace{A})=\dim(\rangespace{\trans{A}A}) \).
      \partsitem \( \text{peringkat kolom}(A)=\text{peringkat kolom}(\trans{A})
        =\text{peringkat baris}(A) \).
    \end{exparts}
    \begin{answer}
      \answerasgiven
      \begin{exparts}
        \partsitem Jelas.
        \partsitem Jika \( \trans{A}A\vec{x}=\zero \), maka 
          \( \vec{y}\cdot\vec{y}=0 \)
          dengan \( \vec{y}=A\vec{x} \).
          Jadi \( \vec{y}=\zero \) berdasarkan (a).

          Kebalikannya jelas.
        \partsitem Berdasarkan (b), \( A\vec{x}_1 \),\ldots,\( A\vec{x}_n \) 
          bebas
          linear jika dan hanya jika \( \trans{A}A\vec{x}_1 \),\ldots,
          \( \trans{A}A\vec{x}_n \) bebas linear.
        \partsitem Kita memiliki 
          \begin{multline*}
            \text{peringkat kolom}(A)=\text{peringkat kolom}(\trans{A}A)
            =\dim\set{\trans{A}(A\vec{x})\suchthat \text{semua \(\vec{x}\)} }  \\
            \leq \dim\set{\trans{A}\vec{y}\suchthat \text{semua \(\vec{y}\)} }
            =\text{peringkat kolom}(\trans{A}).
          \end{multline*}
          Jadi juga \( \text{peringkat kolom}(\trans{A})
                       \leq\text{peringkat kolom}(\trans{\trans{A}}) \),
          sehingga \( \text{peringkat kolom}(A)=\text{peringkat kolom}(\trans{A})
                    =\text{peringkat baris}(A) \).
      \end{exparts} 
    \end{answer}
  \item  
    \cite{Monthly55p721}
    Buktikan
    (dengan \( A \) matriks \( \nbyn{n} \), sehingga mendefinisikan
    transformasi pada sebarang ruang berdimensi \( n \), \( V \),
    terhadap \( B,B \), dengan \( B \) suatu basis) bahwa
    \( \dim(\rangespace{A}\intersection\nullspace{A})
           =\dim(\rangespace{A})-\dim(\rangespace{A^2}) \).
    Simpulkan
    \begin{exparts}
      \partsitem \( \nullspace{A}\subset\rangespace{A} \)
            jika dan hanya jika
            \( \dim(\nullspace{A})=\dim(\rangespace{A})
                                   -\dim(\rangespace{A^2}) \);
      \partsitem \( \rangespace{A}\subseteq\nullspace{A} \)
            jika dan hanya jika
            \( A^2=0 \);
      \partsitem \( \rangespace{A}=\nullspace{A} \)
            jika dan hanya jika
            \( A^2=0 \)
            dan \( \dim(\nullspace{A})=\dim(\rangespace{A}) \) ;
      \partsitem \( \dim(\rangespace{A}\intersection\nullspace{A})=0 \)
            jika dan hanya jika
            \( \dim(\rangespace{A})=\dim(\rangespace{A^2}) \) ;
      \partsitem \textit{(Memerlukan subbagian Jumlah Langsung, 
            yang bersifat opsional.)}
            \( V=\rangespace{A}\directsum\nullspace{A} \)
            jika dan hanya jika
            \( \dim(\rangespace{A})=\dim(\rangespace{A^2}) \).
    \end{exparts}
    \begin{answer}
      \answerasgiven
      Misalkan \( \sequence{\vec{z}_1,\dots,\vec{z}_k} \) adalah basis bagi
      \( \rangespace{A}\intersection\nullspace{A} \)
      (\( k \) mungkin bernilai \( 0 \)).
      Pilih \( \vec{x}_1,\dots,\vec{x}_k\in V \) sedemikian sehingga
      \( A\vec{x}_i=\vec{z}_i \).
      Perhatikan bahwa \( \set{A\vec{x}_1,\dots,A\vec{x}_k} \) bebas linear,
      dan perluaslah menjadi basis bagi \( \rangespace{A} \):
      \( A\vec{x}_1,\ldots,A\vec{x}_k,A\vec{x}_{k+1},\dots,A\vec{x}_{r_1} \)
      dengan \( r_1=\dim(\rangespace{A}) \).

      Sekarang ambil \( \vec{x}\in V \).
      Tuliskan
      \begin{equation*}
        A\vec{x}=a_1(A\vec{x}_1)+\dots+a_{r_1}(A\vec{x}_{r_1})
      \end{equation*}
      sehingga
      \begin{equation*}
        A^2\vec{x}=a_1(A^2\vec{x}_1)+\dots+a_{r_1}(A^2\vec{x}_{r_1}).
      \end{equation*}
      Namun, \( A\vec{x}_1,\dots,A\vec{x}_k\in\nullspace{A} \), sehingga
      \( A^2\vec{x}_1=\zero,\dots,A^2\vec{x}_k=\zero \), dan kini kita mengetahui bahwa
      \begin{equation*}
        A^2\vec{x}_{k+1},\dots,A^2\vec{x}_{r_1}
      \end{equation*}
      merentang \( \rangespace{A^2} \).

      Untuk melihat bahwa \( \set{A^2\vec{x}_{k+1},\dots,A^2\vec{x}_{r_1}} \)
      bebas linear, tuliskan
      \begin{align*}
        b_{k+1}A^2\vec{x}_{k+1}+\dots+b_{r_1}A^2\vec{x}_{r_1}
        &=\zero                                                 \\
        A[b_{k+1}A\vec{x}_{k+1}+\dots+b_{r_1}A\vec{x}_{r_1}]
        &=\zero
      \end{align*}
      dan, karena
      \( b_{k+1}A\vec{x}_{k+1}+\dots+b_{r_1}A\vec{x}_{r_1}\in\nullspace{A} \)
      kita memperoleh kontradiksi kecuali ekspresi itu sama dengan \( \zero \) (jelas bahwa ekspresi itu berada dalam
      \( \rangespace{A} \), sedangkan \( A\vec{x}_1,\ldots,A\vec{x}_k \) adalah basis
      bagi \( \rangespace{A}\intersection\nullspace{A} \)).

      Oleh karena itu, \( \dim(\rangespace{A^2})=r_1-k=\dim(\rangespace{A})-
                \dim(\rangespace{A}\intersection\nullspace{A}) \).  
    \end{answer}
\end{exercises}























\subsection{Invers}
\index{matriks!invers|(}
%<*FunctionInverseReview0>
Kita menutup bagian ini dengan mempertimbangkan cara merepresentasikan 
invers suatu peta linear.
%</FunctionInverseReview0>
%<*FunctionInverseReview1>
Pertama-tama, kita ingat kembali beberapa hal tentang
invers. % \appendrefs{function inverses}\spacefactor=1000 
% Some functions have no inverse, or have an inverse on the left side 
%or right side only.
% 
% \begin{example}  \label{ex:ProjLeftInvOfEmbed}
Misalkan 
\( \map{\pi}{\Re^3}{\Re^2} \) adalah peta proyeksi 
dan \( \map{\iota}{\Re^2}{\Re^3} \) adalah penyisipan  
\begin{equation*}
  \colvec{x \\ y \\ z}
    \mapsunder{\pi}
  \colvec{x \\ y}
  \qquad
  \colvec{x \\ y}
    \mapsunder{\iota}
  \colvec{x \\ y \\ 0}
\end{equation*}
maka komposisi $\composed{\pi}{\iota}$ adalah peta identitas
$\composed{\pi}{\iota}=\identity$ pada $\Re^2$.
\begin{equation*}
  %\composed{\pi}{\eta}:\qquad
  \colvec{x \\ y}
    \mapsunder{\iota}
  \colvec{x \\ y \\ 0}
    \mapsunder{\pi}
  \colvec{x \\ y}
\end{equation*}
Kita mengatakan bahwa
$\iota$ adalah \definend{invers kanan}\index{fungsi!invers kanan} 
dari $\pi$,
atau, yang sama artinya, 
bahwa $\pi$ adalah \definend{invers kiri}\index{fungsi!invers kiri} 
dari $\iota$.
%</FunctionInverseReview1>
%<*FunctionInverseReview2>
Namun, komposisi dalam urutan sebaliknya, $\composed{\iota}{\pi}$, 
tidak menghasilkan peta identitas\Dash berikut adalah vektor yang tidak 
dikirim ke dirinya sendiri oleh $\composed{\iota}{\pi}$. 
\begin{equation*}
  %\composed{\eta}{\pi}:\qquad
  \colvec[r]{0 \\ 0 \\ 1}
    \mapsunder{\pi}
  \colvec[r]{0 \\ 0}
    \mapsunder{\iota}
  \colvec[r]{0 \\ 0 \\ 0}
\end{equation*}
%</FunctionInverseReview2>
%<*FunctionInverseReview3>
Bahkan,
$\pi$ sama sekali tidak memiliki invers kiri.
Sebab, jika $f$ merupakan invers kiri dari $\pi$,
maka kita akan memperoleh
\begin{equation*}
  %\composed{f}{\pi}:\qquad
  \colvec{x \\ y \\ z}
    \mapsunder{\pi}
  \colvec{x \\ y}
    \mapsunder{f}
  \colvec{x \\ y \\ z}
\end{equation*}
untuk takhingga banyaknya nilai $z$.
Namun, fungsi~$f$ tidak dapat mengirim satu argumen $\binom{x}{y}$ ke 
lebih dari satu nilai.
% \end{example}

Jadi suatu fungsi dapat memiliki invers kanan tetapi tidak memiliki invers kiri, atau memiliki invers kiri
tetapi tidak memiliki invers kanan.
Suatu fungsi juga dapat tidak memiliki invers dari kedua sisi; salah satu contohnya 
adalah transformasi nol pada $\Re^2$.
%</FunctionInverseReview3>


%<*FunctionInverseReview4>
Beberapa fungsi memiliki 
\definend{invers dua sisi},\index{fungsi!invers dua sisi} 
yaitu fungsi lain
yang menjadi invers baik dari kiri maupun dari kanan.
Sebagai contoh, transformasi 
$\vec{v}\mapsto 2\cdot \vec{v}$ memiliki invers dua sisi 
$\vec{v}\mapsto (1/2)\cdot\vec{v}$.  
% In this subsection we will focus on two-sided inverses.
% 
Lampiran menunjukkan bahwa suatu fungsi
memiliki invers dua sisi jika dan hanya jika fungsi itu satu-satu sekaligus surjektif.
Lampiran juga menunjukkan bahwa jika fungsi $f$ memiliki invers dua sisi, invers tersebut 
tunggal, sehingga kita menyebutnya 
`invers'~\index{fungsi invers}\index{fungsi!invers}
dan menuliskannya sebagai $f^{-1}$.
%</FunctionInverseReview4>

%<*FunctionInverseReview5>
Selain itu, ingatlah bahwa dalam Teorema~II.\ref{th:OOHomoEquivalence} telah kita tunjukkan
bahwa jika suatu peta linear memiliki invers dua sisi,
invers itu juga linear. 
%</FunctionInverseReview5>

%<*FunctionInverseReview6>
Jadi, tujuan kita dalam subbagian ini adalah, jika peta linear $h$ memiliki invers,
menemukan hubungan antara $\rep{h}{B,D}$ dan $\rep{h^{-1}}{D,B}$.
%</FunctionInverseReview6>

\begin{definition} \label{df:MatrixInverse}
%<*df:MatrixInverse>
Suatu matriks \( G \) adalah \definend{matriks invers kiri}\index{invers!kiri}
dari matriks \( H \) jika \( GH \) merupakan matriks identitas.
Matriks \( G \) adalah \definend{invers kanan}\index{invers!kanan}
jika \( HG \) merupakan matriks identitas.
Matriks $H$ yang memiliki invers dua sisi disebut \definend{matriks invertibel}.
Invers dua sisi tersebut
dinotasikan dengan \( H^{-1} \).\index{invers}\index{matriks!invers, definisi}
%</df:MatrixInverse>
\end{definition}

Karena adanya korespondensi antara peta linear dan matriks,
pernyataan tentang invers peta dapat diterjemahkan menjadi pernyataan tentang invers matriks.

\begin{lemma}     \label{le:LeftAndRightInvEqual}
%<*lm:LeftAndRightInvEqual>
Jika suatu matriks memiliki invers kiri sekaligus invers kanan, keduanya sama.
%</lm:LeftAndRightInvEqual>
\end{lemma}

\begin{theorem}   \label{th:MatrixInvertibleIffNonsingular}
\index{invers!keberadaan}\index{nonsingular}
%<*th:MatrixInvertibleIffNonsingular>
Suatu matriks invertibel jika dan hanya jika matriks itu nonsingular.
%</th:MatrixInvertibleIffNonsingular>
\end{theorem}

\begin{proof}
%<*pf:MatrixInvertibleIffNonsingular>
\textit{(Untuk kedua hasil.)}
Diberikan suatu matriks $H$, tetapkan ruang-ruang berdimensi sesuai sebagai domain
dan kodomain, lalu
tetapkan basis bagi ruang-ruang tersebut.
Terhadap basis-basis ini, $H$ merepresentasikan suatu peta $h$.
Pernyataan-pernyataan tersebut berlaku bagi peta, sehingga berlaku pula bagi
matriks.
%</pf:MatrixInvertibleIffNonsingular>
\end{proof}

\begin{lemma} \label{lem:ProdInvIsInv}
%<*lm:ProdInvIsInv>
Hasil kali matriks-matriks invertibel bersifat invertibel:~jika
\( G \) dan \( H \) invertibel serta \( GH \) terdefinisi, maka
\( GH \) invertibel dan \( (GH)^{-1}=H^{-1}G^{-1} \).
%</lm:ProdInvIsInv>
\end{lemma}

\begin{proof}
% \textit{(This is just like the prior proof except that it requires two maps.)}
%<*pf:ProdInvIsInv0>
Karena kedua matriks tersebut invertibel, keduanya persegi, dan
karena hasil kalinya terdefinisi, keduanya harus berukuran
$\nbyn{n}$.
Tetapkan ruang dan basis\Dash misalnya, $\Re^n$ dengan basis-basis standar\Dash
untuk memperoleh peta-peta
$\map{g,h}{\Re^n}{\Re^n}$ yang terkait dengan matriks-matriks tersebut,
$G=\rep{g}{\stdbasis_n,\stdbasis_n}$ dan
$H=\rep{h}{\stdbasis_n,\stdbasis_n}$.
%</pf:ProdInvIsInv0>

%<*pf:ProdInvIsInv1>
Perhatikan $h^{-1}g^{-1}$.
Berdasarkan paragraf sebelumnya, komposisi ini terdefinisi.
Peta ini merupakan invers dua sisi dari \( gh \) karena
\(
  (h^{-1}g^{-1})(gh)
  =
  h^{-1}(\identity)h
  =h^{-1}h
  =\identity
\)
dan
\(
  (gh)(h^{-1}g^{-1})
  =
  g(\identity)g^{-1}
  =gg^{-1}
  =\identity
\).
Matriks-matriks yang merepresentasikan peta-peta tersebut mencerminkan kesamaan ini.
%</pf:ProdInvIsInv1>
\end{proof}

Berikut adalah diagram panah\index{diagram panah} yang menunjukkan hubungan
antara invers peta dan invers matriks.
Diagram ini merupakan kasus khusus
dari diagram yang menghubungkan komposisi fungsi dengan perkalian matriks.
\smallskip
\begin{center}
    \includegraphics{map/mp/ch3.21}
\end{center}

Selain perannya dalam upaya kita untuk
melihat cara merepresentasikan operasi peta,
alasan lain ketertarikan kita pada invers berasal dari
sistem linear.
Suatu sistem linear setara dengan persamaan matriks, seperti berikut.
\begin{equation*}
  \begin{linsys}{2}
    x_1  &+  &x_2  &=  &3  \\
   2x_1  &-  &x_2  &=  &2  
  \end{linsys}
  \quad\Longleftrightarrow\quad
  \begin{mat}[r]
       1  &1  \\
       2  &-1
    \end{mat}
  \colvec{x_1 \\ x_2}
  =
  \colvec[r]{3 \\ 2}
 % \tag*{($*$)}
\end{equation*}
Dengan menetapkan ruang dan basis
(misalnya, $\Re^2,\Re^2$ dengan basis-basis standar),
kita mengambil matriks $H$ sebagai representasi suatu peta $h$.
Persamaan matriks tersebut kemudian menjadi persamaan peta linear berikut.
\begin{equation*}
  h(\vec{x})=\vec{d}
  % \tag*{($**$)}
\end{equation*}
Jika kita memiliki peta invers kiri~$g$, maka
kita dapat menerapkannya pada kedua ruas,
$\composed{g}{h}(\vec{x})=g(\vec{d})$, untuk memperoleh
$\vec{x}=g(\vec{d})$.
Dengan menyatakannya kembali dalam matriks,
kita ingin mengalikan matriks invers dengan $\rep{\vec{d}}{C}$, yaitu
$\rep{g}{C,B}\cdot\rep{\vec{d}}{C}$, untuk memperoleh $\rep{\vec{x}}{B}$.

\begin{example}  \label{ex:InverseByLinSys}
Kita dapat mencari invers kiri bagi matriks yang baru saja diberikan
\begin{equation*}
    \begin{mat}
       m  &n  \\
       p  &q
    \end{mat}
    \begin{mat}[r]
       1  &1  \\
       2  &-1
    \end{mat}
  =
    \begin{mat}[r]
       1  &0  \\
       0  &1
    \end{mat}
\end{equation*}
dengan menggunakan Metode Gauss untuk menyelesaikan sistem linear yang dihasilkan.
\begin{equation*}
  \begin{linsys}{4}
     m  &+  &2n  &    &   &   &    &=  &1     \\
     m  &-  &n   &    &   &   &    &=  &0     \\
        &   &    &    &p  &+  &2q  &=  &0     \\
        &   &    &    &p  &-  &q   &=  &1     
     \end{linsys}
\end{equation*}
Jawaban: \( m=1/3 \), \( n=1/3 \), \( p=2/3 \), dan \( q=-1/3 \).
(Matriks ini sebenarnya merupakan invers dua sisi dari $H$;
pemeriksaannya mudah.)
Dengan matriks ini, kita dapat menyelesaikan sistem dari contoh sebelumnya.
\begin{equation*}
  \colvec{x \\ y}
  =\begin{mat}[r]
       1/3  &1/3  \\
       2/3  &-1/3
    \end{mat}
  \colvec[r]{3 \\ 2}          
  =\colvec[r]{5/3 \\ 4/3}
\end{equation*}
\end{example}

\begin{remark}
Mengapa menyelesaikan sistem dengan matriks invers jika kita memiliki
Metode Gauss?
% takes less arithmetic
% (this assertion can be made precise by counting the 
% number of arithmetic operations,
% as computer algorithm designers do)? 
Selain daya tarik konseptual dari merepresentasikan operasi invers peta,
menyelesaikan sistem linear dengan cara ini memiliki
dua kelebihan.

Pertama, setelah kita melakukan pekerjaan untuk mencari invers,
penyelesaian sistem dengan
koefisien yang sama tetapi konstanta yang berbeda dapat dilakukan dengan cepat:~jika
kita mengubah konstanta di ruas kanan sistem di atas,
kita memperoleh masalah terkait
\begin{equation*}
    \begin{mat}[r]
       1  &1  \\
       2  &-1
    \end{mat}
  \colvec{x \\ y}
  =
  \colvec[r]{5 \\ 1}
\end{equation*}
yang dapat diselesaikan dengan cepat oleh metode invers kita.
\begin{equation*}
  \colvec{x \\ y}
  =
    \begin{mat}[r]
       1/3  &1/3  \\
       2/3  &-1/3
    \end{mat}
  \colvec[r]{5 \\ 1}
  =
  \colvec[r]{2 \\ 3}
\end{equation*}
% In applications we often must solve many systems with the same matrix of
% coefficients.

Kelebihan lain invers adalah kita dapat
menyelidiki sensitivitas suatu sistem terhadap perubahan konstanta.
Sebagai contoh, ubah sedikit angka $3$ di ruas kanan sistem pada contoh sebelumnya menjadi
\begin{equation*}
    \begin{mat}[r]
       1  &1  \\
       2  &-1
    \end{mat}
  \colvec{x_1 \\ x_2}
  =
  \colvec[l]{3.01 \\ 2}
\end{equation*}
dan menyelesaikannya dengan invers
\begin{equation*}
    \begin{mat}[r]
       1/3  &1/3  \\
       2/3  &-1/3
    \end{mat}
  \colvec[l]{3.01 \\ 2}
  =
  \colvec{(1/3)(3.01)+(1/3)(2) \\ (2/3)(3.01)-(1/3)(2)}
\end{equation*}
menunjukkan bahwa komponen pertama solusi
berubah sebesar \( 1/3 \) dari perubahan tersebut, sedangkan
komponen kedua berubah sebesar \( 2/3 \) dari perubahan tersebut.
Ini disebut \definend{analisis sensitivitas}\index{analisis sensitivitas}.
Kita dapat menggunakannya untuk menentukan seberapa teliti
data dalam suatu model linear harus dinyatakan agar solusi memiliki
ketelitian yang diinginkan.
\end{remark}

% We finish by describing the computational procedure
% that we shall use to find the inverse matrix.

\begin{lemma} \label{lem:ComputeInvMat}
%<*lm:ComputeInvMat>
Suatu matriks $H$ invertibel jika dan hanya jika dapat dituliskan sebagai hasil kali
matriks-matriks reduksi elementer.
Kita dapat menghitung inversnya dengan menerapkan pada
matriks identitas langkah-langkah baris yang sama, dalam urutan yang sama, yang
mereduksi $H$ dengan Gauss--Jordan.
%</lm:ComputeInvMat>
\end{lemma}

\begin{proof}
%<*pf:ComputeInvMat0>
Matriks $H$ invertibel jika dan hanya jika matriks itu nonsingular dan dengan demikian
tereduksi menjadi identitas oleh Gauss--Jordan.
Berdasarkan \nearbycorollary{cor:ReducViaMatrices}, kita dapat melakukan reduksi ini
dengan matriks-matriks elementer.
\begin{equation*} 
   R_r\cdot R_{r-1}\dots R_1\cdot H=I 
   \tag{$*$}
\end{equation*}
%</pf:ComputeInvMat0>

%<*pf:ComputeInvMat1>
Untuk kalimat pertama hasil tersebut, perhatikan bahwa
matriks-matriks elementer bersifat invertibel karena operasi baris elementer
dapat dibalik, dan inversnya juga elementer.
Terapkan $R_r^{-1}$ dari kiri pada kedua ruas ($*$).
Kemudian terapkan
$R_{r-1}^{-1}$, dan seterusnya.
Hasilnya memberikan $H$ sebagai hasil kali
matriks-matriks elementer $H=R_1^{-1}\cdots R_r^{-1}\cdot I$.
($I$ di sana mencakup kasus $r=0$.)
%</pf:ComputeInvMat1>

%<*pf:ComputeInvMat2>
Untuk kalimat kedua,
kelompokkan~($*$) sebagai $(R_r\cdot R_{r-1}\dots R_1)\cdot H=I$
dan kenali bahwa isi tanda kurung adalah invers
$H^{-1}=R_r\cdot R_{r-1}\dots R_1\cdot I$.
Dengan kata lain: menerapkan $R_1$ pada identitas,
diikuti oleh $R_2$, dan seterusnya, menghasilkan invers $H$.
%</pf:ComputeInvMat2>
\end{proof}

\begin{example}
Untuk mencari invers dari
\begin{equation*}
    \begin{mat}[r]
       1  &1  \\
       2  &-1
    \end{mat}
\end{equation*}
lakukan reduksi Gauss--Jordan sambil melakukan operasi yang sama pada
identitas.
Agar pencatatannya mudah, kita tuliskan matriks tersebut dan identitas secara berdampingan
dan lakukan langkah-langkah reduksi pada keduanya sekaligus.
\begin{align*}
    \begin{pmat}{rr|rr}
       1  &1   &1  &0  \\
       2  &-1  &0  &1
    \end{pmat}
  &\grstep{-2\rho_1+\rho_2}
  \begin{pmat}{rr|rr}
        1  &1   &1  &0  \\
        0  &-3  &-2 &1
     \end{pmat}                 \\
  &\grstep{-1/3\rho_2}
  \begin{pmat}{rr|rr}
        1  &1   &1   &0     \\
        0  &1   &2/3 &-1/3
     \end{pmat}                   \\
  &\grstep{-\rho_2+\rho_1}
  \begin{pmat}{rr|rr}
        1  &0   &1/3 &1/3  \\
        0  &1   &2/3 &-1/3
     \end{pmat}
\end{align*}
Perhitungan ini telah menghasilkan inversnya.
\begin{equation*}
    \begin{mat}[r]
       1  &1  \\
       2  &-1
    \end{mat}^{-1}
  =
    \begin{mat}[r]
       1/3 &1/3  \\
       2/3 &-1/3
    \end{mat}
\end{equation*}
\end{example}

\begin{example} \label{exam:ThreeByThreeMatInv}
Contoh ini kebetulan diawali dengan pertukaran baris.
\begin{align*}
     \begin{pmat}{rrr|rrr}
        0  &3  &-1  &1  &0  &0  \\
        1  &0  &1   &0  &1  &0  \\
        1  &-1 &0   &0  &0  &1
     \end{pmat}
  &\grstep{\rho_1\leftrightarrow\rho_2}
  \begin{pmat}{rrr|rrr}
         1  &0  &1   &0  &1  &0  \\
         0  &3  &-1  &1  &0  &0  \\
         1  &-1 &0   &0  &0  &1
      \end{pmat}                             \\
  &\grstep{-\rho_1+\rho_3}
  \begin{pmat}{rrr|rrr}
         1  &0  &1   &0  &1  &0  \\
         0  &3  &-1  &1  &0  &0  \\
         0  &-1 &-1  &0  &-1 &1
      \end{pmat}                             \\
  &\quad\vdotswithin{\longrightarrow}                        \\
  &\;\grstep{}
  \begin{pmat}{rrr|rrr}
         1  &0  &0   &1/4  &1/4  &3/4  \\
         0  &1  &0   &1/4  &1/4  &-1/4 \\
         0  &0  &1   &-1/4 &3/4  &-3/4
      \end{pmat}
\end{align*}
\end{example}

\begin{example}
Algoritma ini mendeteksi
matriks yang tidak invertibel ketika bagian kirinya tidak dapat
direduksi menjadi identitas.
\begin{equation*}
    \begin{pmat}{rr|rr}
       1  &1   &1  &0  \\
       2  &2   &0  &1
    \end{pmat}
  \grstep{-2\rho_1+\rho_2}
   \begin{pmat}{rr|rr}
       1  &1   &1  &0  \\
       0  &0   &-2 &1
    \end{pmat}
\end{equation*}
\end{example}

Dengan prosedur ini kita dapat memberikan rumus bagi
invers matriks $\nbyn{2}$ umum,
yang patut dihafalkan.
% But larger matrices have more complex formulas so 
% we will wait for more explanation in the next chapter.

\begin{corollary} \label{cor:TwoByTwoInv}
%<*co:TwoByTwoInv>
Invers suatu matriks \( \nbyn{2} \) ada dan sama dengan
\begin{equation*}
  \begin{mat}
    a  &b  \\
    c  &d
  \end{mat}^{-1}
  =
  \frac{1}{ad-bc}
  \begin{mat}
    d  &-b \\
   -c  &a
  \end{mat}
\end{equation*}
jika dan hanya jika \( ad-bc\neq 0 \).
%</co:TwoByTwoInv>
\end{corollary}
\begin{proof}
%<*pf:TwoByTwoInv>
Perhitungan ini merupakan
\nearbyexercise{exer:TwoByTwoInv}.
%</pf:TwoByTwoInv>
\end{proof}

Dalam subbagian ini, seperti dalam subbagian Mekanika Perkalian Matriks,
kita telah melihat cara memanfaatkan korespondensi antara
peta linear dan matriks.
Kita dapat mengkaji baik peta maupun matriks secara produktif, dengan menerjemahkan bolak-balik
untuk menggunakan mana pun yang paling praktis.

Sepanjang
bagian ini kita telah mengembangkan suatu sistem aljabar bagi matriks.
Kita dapat membandingkannya dengan aljabar bilangan real yang sudah dikenal.
Penjumlahan dan pengurangan matriks
bekerja dengan cara yang kurang lebih sama
seperti operasi bilangan real, kecuali bahwa keduanya hanya menggabungkan
matriks-matriks berukuran sama.
Perkalian skalar dalam beberapa hal merupakan perluasan
dari perkalian bilangan real.
Kita juga memiliki operasi perkalian matriks
beserta inversnya
yang agak menyerupai operasi bilangan real yang sudah dikenal
(misalnya, sifat asosiatif dan distributif terhadap penjumlahan), tetapi
terdapat perbedaan (tidak adanya sifat komutatif).
Bagian ini memberi contoh bahwa sistem aljabar
selain sistem bilangan real yang biasa
dapat menarik dan berguna.


\begin{exercises}
  \item 
    Lengkapilah langkah-langkah antara dalam
    \nearbyexample{exam:ThreeByThreeMatInv}.
    \begin{answer}
      Berikut adalah salah satu cara melakukannya.
      Lanjutkan
      \begin{equation*}
       \grstep{\rho_1\leftrightarrow\rho_2}
       \begin{pmat}{rrr|rrr}
              1  &0  &1   &0  &1  &0  \\
              0  &3  &-1  &1  &0  &0  \\
              1  &-1 &0   &0  &0  &1
           \end{pmat}                          
       \grstep{-\rho_1+\rho_3}
       \begin{pmat}{rrr|rrr}
              1  &0  &1   &0  &1  &0  \\
              0  &3  &-1  &1  &0  &0  \\
              0  &-1 &-1  &0  &-1 &1
           \end{pmat}                        
       \end{equation*}
       dengan
       \begin{align*}
         &\grstep{(1/3)\rho_2+\rho_3}
         \begin{pmat}{rrr|rrr}
                1  &0  &1     &0   &1  &0  \\
                0  &3  &-1    &1   &0  &0  \\
                0  &0  &-4/3  &1/3 &-1 &1
             \end{pmat}                                           \\ 
         &\grstep[-(3/4)\rho_3]{(1/3)\rho_2}
         \begin{pmat}{rrr|rrr}
                1  &0  &1     &0    &1   &0    \\
                0  &1  &-1/3  &1/3  &0   &0    \\
                0  &0  &1     &-1/4 &3/4 &-3/4
             \end{pmat}                                  \\
         &\grstep[-\rho_3+\rho_1]{(1/3)\rho_3+\rho_2}
         \begin{pmat}{rrr|rrr}
                1  &0  &0     &1/4  &1/4 &3/4  \\
                0  &1  &0     &1/4  &1/4 &-1/4 \\
                0  &0  &1     &-1/4 &3/4 &-3/4
             \end{pmat}                         
    \end{align*} 
    dan bacalah jawabannya dari ruas kanan.
   \end{answer}
 \recommended \item 
    Gunakan \nearbycorollary{cor:TwoByTwoInv} untuk menentukan apakah setiap matriks
    memiliki invers.
    \begin{exparts*}
      \partsitem
        $\begin{mat}[r]
           2  &1  \\
          -1  &1          
        \end{mat}$
      \partsitem
        $\begin{mat}[r]
          0  &4  \\
          1  &-3
        \end{mat}$
      \partsitem
        $\begin{mat}[r]
          2  &-3  \\
         -4  &6
        \end{mat}$
    \end{exparts*}
    \begin{answer}
      \begin{exparts*}
        \partsitem Ya, matriks ini memiliki invers: $ad-bc=2\cdot 1-1\cdot(-1)\neq 0$.
        \partsitem Ya.
        \partsitem Tidak.
      \end{exparts*}
    \end{answer}
  \recommended \item  
     Untuk setiap matriks invertibel dalam soal sebelumnya, gunakan
     \nearbycorollary{cor:TwoByTwoInv} untuk mencari inversnya.
     \begin{answer}
       \begin{exparts}
         \partsitem 
           $\displaystyle \frac{1}{2\cdot 1-1\cdot (-1)}
            \cdot\begin{mat}[r]
             1  &-1  \\
             1  &2
           \end{mat}
          =\frac{\displaystyle 1}{\displaystyle 3}\cdot\begin{mat}[r]
             1  &-1  \\
             1  &2
           \end{mat}
          =\begin{mat}[r]
             1/3  &-1/3  \\
             1/3  &2/3
           \end{mat}$
         \partsitem 
           $\displaystyle \frac{1}{0\cdot (-3)-4\cdot 1}
           \cdot\begin{mat}[r]
             -3  &-4  \\
             -1  &0
           \end{mat}
           =\begin{mat}[r]
             3/4  &1  \\
             1/4  &0
           \end{mat}$
         \partsitem Soal sebelumnya menunjukkan bahwa inversnya tidak ada.
       \end{exparts}
     \end{answer}
  \recommended \item
    Carilah inversnya, jika ada, dengan menggunakan Metode Gauss--Jordan.
    Periksa jawaban bagi matriks-matriks $\nbyn{2}$
    dengan \nearbycorollary{cor:TwoByTwoInv}.
    \begin{exparts*}
      \partsitem $\begin{mat}[r]
                   3  &1  \\
                   0  &2
                 \end{mat}$
      \partsitem \( \begin{mat}[r]
                 2   &1/2  \\
                 3   &1
               \end{mat} \)
      \partsitem \( \begin{mat}[r]
                 2   &-4   \\
                -1   &2
               \end{mat} \)
      \partsitem \( \begin{mat}[r]
                 1   &1  &3  \\
                 0   &2  &4  \\
                 -1  &1  &0
               \end{mat} \)
      \partsitem \( \begin{mat}[r]
                 0   &1  &5  \\
                 0   &-2 &4  \\
                 2   &3  &-2
               \end{mat} \)
      \partsitem \( \begin{mat}[r]
                 2   &2  &3  \\
                 1   &-2 &-3 \\
                 4   &-2 &-3
               \end{mat} \)
    \end{exparts*}
    \begin{answer}
      \begin{exparts}
        \partsitem Reduksinya rutin.
          \begin{align*}
            \begin{pmat}{rr|rr}
              3  &1  &1  &0 \\
              0  &2  &0  &1
            \end{pmat}
            &\grstep[(1/2)\rho_2]{(1/3)\rho_1}
            \begin{pmat}{rr|rr}
              1  &1/3  &1/3  &0   \\
              0  &1    &0    &1/2
            \end{pmat}                                 \\
            &\grstep{-(1/3)\rho_2+\rho_1}
            \begin{pmat}{rr|rr}
              1  &0    &1/3  &-1/6 \\
              0  &1    &0    &1/2
            \end{pmat}
          \end{align*}
          Jawaban ini sesuai dengan jawaban dari pemeriksaan.
          \begin{equation*}
            \begin{mat}[r]
              3  &1  \\
              0  &2
            \end{mat}^{-1}
            =\frac{1}{3\cdot 2-0\cdot 1}\cdot
            \begin{mat}[r]
              2  &-1  \\
              0  &3
            \end{mat}
            =\frac{1}{6}\cdot
            \begin{mat}[r]
              2  &-1  \\
              0  &3
            \end{mat}
          \end{equation*}
        \partsitem Reduksi ini mudah.
          \begin{multline*}
            \begin{pmat}{rr|rr}
              2  &1/2  &1  &0  \\
              3  &1    &0  &1
            \end{pmat}
            \grstep{-(3/2)\rho_1+\rho_2}
            \begin{pmat}{rr|rr}
              2  &1/2  &1     &0  \\
              0  &1/4  &-3/2  &1
            \end{pmat}                       \\
            \grstep[4\rho_2]{(1/2)\rho_1}
            \begin{pmat}{rr|rr}
              1  &1/4   &1/2   &0  \\
              0  &1     &-6    &4
            \end{pmat}                         
            \grstep{-(1/4)\rho_2+\rho_1}
            \begin{pmat}{rr|rr}
              1  &0     &2     &-1 \\
              0  &1     &-6    &4
            \end{pmat}
          \end{multline*}
          Hasil pemeriksaannya sesuai.
          \begin{equation*}
            \frac{1}{2\cdot 1-3\cdot (1/2)}\cdot
            \begin{mat}[r]
              1  &-1/2  \\
              -3 &2
            \end{mat}
            =2\cdot
            \begin{mat}[r]
              1  &-1/2  \\
              -3 &2
            \end{mat}
          \end{equation*}
        \partsitem Mencoba reduksi Gauss--Jordan
          \begin{equation*}
            \begin{pmat}{rr|rr}
              2  &-4  &1  &0  \\
             -1  &2   &0  &1 
            \end{pmat}
            \grstep{(1/2)\rho_1+\rho_2}
            \begin{pmat}{rr|rr}
              2  &-4  &1   &0  \\
              0  &0   &1/2 &1 
            \end{pmat}
          \end{equation*}
          menunjukkan bahwa ruas kiri tidak dapat direduksi menjadi identitas, sehingga inversnya
          tidak ada.
          Pemeriksaan $ad-bc=2\cdot 2-(-4)\cdot (-1)=0$ memberikan hasil yang sama.
        \partsitem Reduksi ini menghasilkan suatu invers.
          \begin{multline*}
            \begin{pmat}{rrr|rrr}
              1  &1  &3  &1  &0  &0  \\ 
              0  &2  &4  &0  &1  &0  \\
             -1  &1  &0  &0  &0  &1 
            \end{pmat}  
            \grstep{\rho_1+\rho_3}
            \begin{pmat}{rrr|rrr}
              1  &1  &3  &1  &0  &0  \\ 
              0  &2  &4  &0  &1  &0  \\
              0  &2  &3  &1  &0  &1 
            \end{pmat}                                             \\
           \begin{aligned}
            &\grstep{-\rho_2+\rho_3}
            \begin{pmat}{rrr|rrr}
              1  &1  &3  &1  &0  &0  \\ 
              0  &2  &4  &0  &1  &0  \\
              0  &0  &-1 &1  &-1 &1 
            \end{pmat}             
           \grstep[-\rho_3]{(1/2)\rho_2}
            \begin{pmat}{rrr|rrr}
              1  &1  &3  &1  &0   &0  \\ 
              0  &1  &2  &0  &1/2 &0  \\
              0  &0  &1  &-1 &1   &-1
            \end{pmat}                                        \\
            &\grstep[-3\rho_3+\rho_1]{-2\rho_3+\rho_2}
            \begin{pmat}{rrr|rrr}
              1  &1  &0  &4  &-3   &3  \\ 
              0  &1  &0  &2  &-3/2 &2  \\
              0  &0  &1  &-1 &1    &-1
            \end{pmat}                                         \\
            &\grstep{-\rho_2+\rho_1}
            \begin{pmat}{rrr|rrr}
              1  &0  &0  &2  &-3/2 &1  \\ 
              0  &1  &0  &2  &-3/2 &2  \\
              0  &0  &1  &-1 &1    &-1
            \end{pmat}
           \end{aligned} 
          \end{multline*}
       \partsitem Berikut adalah salah satu cara melakukan reduksi.
          \begin{multline*}
            \begin{pmat}{rrr|rrr}
              0  &1  &5  &1  &0  &0  \\
              0  &-2 &4  &0  &1  &0  \\ 
              2  &3  &-2 &0  &0  &1
            \end{pmat}
            \grstep{\rho_3\leftrightarrow\rho_1}\;
            \begin{pmat}{rrr|rrr}
              2  &3  &-2 &0  &0  &1  \\
              0  &-2 &4  &0  &1  &0  \\ 
              0  &1  &5  &1  &0  &0  
            \end{pmat}                                            \\
            \begin{aligned}
              &\grstep{(1/2)\rho_2+\rho_3}
              \begin{pmat}{rrr|rrr}
                2  &3  &-2 &0  &0   &1  \\
                0  &-2 &4  &0  &1   &0  \\ 
                0  &0  &7  &1  &1/2 &0  
              \end{pmat}                                             \\
              &\grstep[-(1/2)\rho_2 \\ (1/7)\rho_3]{(1/2)\rho_1}
              \begin{pmat}{rrr|rrr}
                1  &3/2  &-1 &0    &0     &1/2  \\
                0  &1    &-2 &0    &-1/2  &0    \\ 
                0  &0    &1  &1/7  &1/14  &0  
              \end{pmat}                                   \\
              &\grstep[\rho_3+\rho_1]{2\rho_3+\rho_2}
              \begin{pmat}{rrr|rrr}
                1  &3/2  &0  &1/7  &1/14  &1/2  \\
                0  &1    &0  &2/7  &-5/14 &0    \\ 
                0  &0    &1  &1/7  &1/14  &0  
              \end{pmat}                                   \\
              &\grstep{-(3/2)\rho_2+\rho_1}
              \begin{pmat}{rrr|rrr}
                1  &0    &0  &-2/7 &17/28 &1/2  \\
                0  &1    &0  &2/7  &-5/14 &0    \\ 
                0  &0    &1  &1/7  &1/14  &0  
              \end{pmat}
            \end{aligned}
          \end{multline*}
        \partsitem Inversnya tidak ada.
          \begin{align*}
            \begin{pmat}{rrr|rrr}
              2  &2  &3  &1  &0  &0  \\
              1  &-2 &-3 &0  &1  &0  \\
              4  &-2 &-3 &0  &0  &1
            \end{pmat}
            &\grstep[-2\rho_1+\rho_3]{-(1/2)\rho_1+\rho_2}
            \begin{pmat}{rrr|rrr}
              2  &2  &3    &1     &0  &0  \\
              0  &-3 &-9/2 &-1/2  &1  &0  \\
              0  &-6 &-9   &-2    &0  &1
            \end{pmat}                                          \\
            &\grstep{-2\rho_2+\rho_3}
            \begin{pmat}{rrr|rrr}
              2  &2  &3    &1     &0   &0  \\
              0  &-3 &-9/2 &-1/2  &1   &0  \\
              0  &0  &0    &-1    &-2  &1
            \end{pmat}
          \end{align*}
          Sebagai pemeriksaan,
          perhatikan bahwa kolom ketiga matriks awal adalah $3/2$ kali
          kolom kedua, sehingga matriks itu memang singular dan karenanya tidak memiliki
          invers.
      \end{exparts}
    \end{answer}
  \recommended \item 
    Matriks apakah yang memiliki matriks berikut sebagai inversnya?
    \begin{equation*}
      \begin{mat}[r]
        1  &3  \\
        2  &5
       \end{mat}
     \end{equation*}
     \begin{answer}
       Kita dapat menggunakan \nearbycorollary{cor:TwoByTwoInv}.
       \begin{equation*}
         \frac{1}{1\cdot 5-2\cdot 3}\cdot
         \begin{mat}[r]
           5  &-3  \\
          -2  &1
         \end{mat}
         =\begin{mat}[r]
           -5  &3  \\
            2  &-1
         \end{mat}
       \end{equation*}  
      \end{answer}
  \item 
   Bagaimana operasi invers berinteraksi dengan perkalian skalar
   dan penjumlahan matriks?
   \begin{exparts}
      \partsitem Apa invers dari \( rH \)?
      \partsitem Apakah \( (H+G)^{-1}=H^{-1}+G^{-1} \)?
    \end{exparts}
    \begin{answer}
      \begin{exparts}
        \partsitem Pembuktian bahwa inversnya adalah
          \( r^{-1}H^{-1}=(1/r)\cdot H^{-1} \)
          (tentu saja, asalkan matriksnya invertibel) mudah.
        \partsitem Tidak.
          Pertama, fakta bahwa $H+G$ memiliki invers tidak mengakibatkan bahwa
          $H$ memiliki invers atau $G$ memiliki invers.
          Kedua matriks berikut tidak invertibel, tetapi jumlahnya invertibel.
          \begin{equation*}
            \begin{mat}[r]
              1  &0  \\
              0  &0
            \end{mat}
            \qquad
            \begin{mat}[r]
              0  &0  \\
              0  &1
            \end{mat}
          \end{equation*}
          Hal lain adalah bahwa meskipun $H$ dan $G$ masing-masing memiliki invers,
          tidak berarti $H+G$ memiliki invers; berikut sebuah contoh.
          \begin{equation*}
            \begin{mat}[r]
              1  &0  \\
              0  &1
            \end{mat}
            \qquad
            \begin{mat}[r]
              -1  &0  \\
               0  &-1
            \end{mat}
          \end{equation*}
          Hal ketiga adalah bahwa meskipun kedua matriks memiliki invers
          dan jumlahnya juga memiliki invers, hal itu tidak mengakibatkan persamaan tersebut berlaku:
          \begin{equation*}
            \begin{mat}[r]
              2  &0  \\
              0  &2
            \end{mat}^{-1}
            =\begin{mat}[r]
              1/2  &0  \\
              0    &1/2
            \end{mat}
            \qquad
            \begin{mat}[r]
              3  &0  \\
              0  &3
            \end{mat}^{-1}
            =\begin{mat}[r]
              1/3  &0  \\
              0    &1/3
            \end{mat}
          \end{equation*}
          tetapi
          \begin{equation*}
            \begin{mat}[r]
              5  &0  \\
              0  &5
            \end{mat}^{-1}
            =\begin{mat}[r]
              1/5  &0  \\
              0    &1/5
            \end{mat}
          \end{equation*}
          dan $(1/2)+(1/3)$ tidak sama dengan $1/5$.
      \end{exparts}  
    \end{answer}
  \recommended \item 
    Apakah \( (T^k)^{-1}=(T^{-1})^k \)?
    \begin{answer}
      Ya:
      \( T^k(T^{-1})^k=(TT\cdots T)\cdot (T^{-1}T^{-1}\cdots T^{-1})
         =T^{k-1}(TT^{-1})(T^{-1})^{k-1}=\dots=I \). 
    \end{answer}
  \item 
    Apakah \( H^{-1} \) invertibel?
    \begin{answer}
       Ya, invers dari \( H^{-1} \) adalah \( H \).
    \end{answer}
  \item 
    Untuk setiap bilangan real \( \theta \), misalkan
    \( \map{t_\theta}{\Re^2}{\Re^2} \) direpresentasikan terhadap
    basis-basis standar oleh matriks berikut.
    \begin{equation*}
      \begin{mat}
         \cos\theta  &-\sin\theta  \\
         \sin\theta  &\cos\theta
      \end{mat}
    \end{equation*}
    Tunjukkan bahwa \( t_{\theta_1+\theta_2}=t_{\theta_1}\cdot t_{\theta_2} \).
    Tunjukkan juga bahwa \( {t_{\theta}}^{-1}=t_{-\theta} \).
    \begin{answer}
      Salah satu cara memeriksa bahwa persamaan pertama benar adalah menggunakan
      rumus jumlah sudut dari trigonometri.
      \begin{multline*}
        \begin{mat}
          \cos(\theta_1+\theta_2) &-\sin(\theta_1+\theta_2)  \\
          \sin(\theta_1+\theta_2) &\cos(\theta_1+\theta_2)
        \end{mat}                                                 \\
        \begin{aligned}
        &=\begin{mat}
          \cos\theta_1\cos\theta_2-\sin\theta_1\sin\theta_2
            &-\sin\theta_1\cos\theta_2-\cos\theta_1\sin\theta_2  \\
          \sin\theta_1\cos\theta_2+\cos\theta_1\sin\theta_2
            &\cos\theta_1\cos\theta_2-\sin\theta_1\sin\theta_2
        \end{mat}                                                    \\
        &=\begin{mat}
          \cos\theta_1 &-\sin\theta_1  \\
          \sin\theta_1 &\cos\theta_1
        \end{mat}
        \begin{mat}
          \cos\theta_2 &-\sin\theta_2  \\
          \sin\theta_2 &\cos\theta_2
        \end{mat}
        \end{aligned}
      \end{multline*}
      Pemeriksaan persamaan kedua dengan cara ini serupa.

      Tentu saja, persamaan-persamaan tersebut tidak hanya dapat diperiksa, tetapi juga dipahami
      dengan mengingat bahwa $t_\theta$ adalah peta yang merotasi
      vektor-vektor mengelilingi titik asal dengan sudut \( \theta \)~radian.
    \end{answer}
  \item \label{exer:TwoByTwoInv}
    Kerjakan perhitungan untuk pembuktian \nearbycorollary{cor:TwoByTwoInv}.
    \begin{answer}
      Ada dua kasus.
      Untuk kasus pertama, kita anggap bahwa \( a \) taknol.
      Maka
      \begin{equation*}
        \grstep{-(c/a)\rho_1+\rho_2}
        \begin{pmat}{cc|cc}
           a   &b           &1     &0   \\
           0   &-(bc/a)+d   &-c/a  &1
         \end{pmat}
        =\begin{pmat}{cc|cc}
           a   &b           &1     &0   \\
           0   &(ad-bc)/a   &-c/a  &1
         \end{pmat}
      \end{equation*}
      menunjukkan bahwa matriks tersebut invertibel (dalam kasus \( a\neq 0 \) ini)
      jika dan hanya jika \( ad-bc\neq 0 \).
      Untuk mencari inversnya, kita selesaikan bagian Jordan dari reduksi.
      \begin{align*}
        &\grstep[(a/(ad-bc))\rho_2]{(1/a)\rho_1}
        \begin{pmat}{cc|cc}
           1   &b/a     &1/a         &0   \\
           0   &1       &-c/(ad-bc)  &a/(ad-bc)
         \end{pmat}                                       \\   
        &\grstep{-(b/a)\rho_2+\rho_1}
        \begin{pmat}{cc|cc}
           1   &0       &d/(ad-bc)   &-b/(ad-bc)   \\
           0   &1       &-c/(ad-bc)  &a/(ad-bc)
         \end{pmat}                                
      \end{align*}

      Kasus lainnya adalah kasus \( a=0 \).
      Kita melakukan pertukaran agar $c$ berada di posisi $1,1$.
      \begin{equation*}
        \grstep{\rho_1\leftrightarrow\rho_2}
        \begin{pmat}{cc|cc}
          c  &d  &0  &1  \\
          0  &b  &1  &0
        \end{pmat}
      \end{equation*}
      Matriks ini nonsingular jika dan hanya jika
      \( b \) dan \( c \) keduanya taknol
      (yang, berdasarkan asumsi kasus bahwa \( a=0 \),
      berlaku jika dan hanya jika \( ad-bc\neq 0 \)).
      Untuk mencari inversnya,
      kita kerjakan bagian Jordan.
      \begin{equation*}
        \grstep[(1/b)\rho_2]{(1/c)\rho_1}
        \begin{pmat}{cc|cc}
          1  &d/c  &0       &1/c  \\
          0  &1    &1/b     &0
        \end{pmat}                        
        \grstep{-(d/c)\rho_2+\rho_1}
        \begin{pmat}{cc|cc}
          1  &0  &-d/bc  &1/c  \\
          0  &1  &1/b    &0
        \end{pmat}
      \end{equation*}
      (Perhatikan bahwa inilah yang diperlukan, karena \( a=0 \) memberikan
      \( ad-bc=-bc \)).
    \end{answer}
  \item 
    Tunjukkan bahwa matriks berikut
    \begin{equation*}
      H=\begin{mat}[r]
          1  &0   &1  \\
          0  &1   &0
        \end{mat}
    \end{equation*}
    memiliki takhingga banyaknya invers kanan.
    Tunjukkan juga bahwa matriks ini tidak memiliki invers kiri.
    \begin{answer}
      Karena $H$ adalah matriks $\nbym{2}{3}$,
      dalam mencari matriks $G$ sedemikian sehingga komposisi $HG$
      bertindak sebagai identitas $\nbyn{2}$,
      kita memerlukan $G$ berukuran $\nbym{3}{2}$.
      Menyusun persamaan
      \begin{equation*}
          \begin{mat}[r]
             1  &0   &1  \\
             0  &1   &0
           \end{mat}
          \begin{mat}
             m  &n  \\
             p  &q  \\
             r  &s
          \end{mat}
        =
          \begin{mat}[r]
             1  &0  \\
             0  &1
          \end{mat}
      \end{equation*}
      dan menyelesaikan sistem linear yang dihasilkan
      \begin{equation*}
        \begin{linsys}{6}
           m  &   &   &   &+r &   &=    &1  \\
              &n  &   &   &   &+s &=    &0  \\
              &   &p  &   &   &   &=    &0  \\
              &   &   &q  &   &   &=    &1
         \end{linsys}
      \end{equation*}
      memberikan takhingga banyaknya solusi.
      \begin{equation*}
        \set{\colvec{m \\ n \\ p \\ q \\ r \\ s}
              =\colvec[r]{1 \\ 0 \\ 0 \\ 1 \\ 0 \\ 0}
              +r\cdot \colvec[r]{-1 \\ 0 \\ 0 \\ 0 \\ 1 \\ 0}
              +s\cdot \colvec[r]{0 \\ -1 \\ 0 \\ 0 \\ 0 \\ 1}
             \suchthat r,s\in\Re}
      \end{equation*}
      Jadi $H$ memiliki takhingga banyaknya invers kanan.

      Untuk invers kiri, persamaannya adalah
      \begin{equation*}
         \begin{mat}
            a  &b  \\
            c  &d  \\
            e  &f
         \end{mat}
         \begin{mat}[r]
             1  &0   &1  \\
             0  &1   &0
          \end{mat}
          =\begin{mat}[r]
             1  &0  &0  \\
             0  &1  &0  \\
             0  &0  &1
           \end{mat}
      \end{equation*}
      menghasilkan sistem linear dengan sembilan persamaan dan enam peubah.
      \begin{equation*}
        \begin{linsys}{6}
          a & & & & & & & & & & &= &1 \\
            & &b& & & & & & & & &= &0 \\
          a & & & & & & & & & & &= &0 \\
            & & & &c& & & & & & &= &0 \\
            & & & & & &d& & & & &= &1 \\
            & & & &c& & & & & & &= &0 \\
            & & & & & & & &e& & &= &0 \\
            & & & & & & & & & &f&= &0 \\
            & & & & & & & &e& & &= &1 
        \end{linsys}
      \end{equation*}
      Sistem ini takkonsisten (persamaan pertama bertentangan
      dengan persamaan ketiga, demikian pula persamaan ketujuh dengan kesembilan),
      sehingga tidak ada invers kiri.
    \end{answer}
  \item 
    Dalam ulasan tentang invers % \nearbyexample{ex:ProjLeftInvOfEmbed},
    pada awal subbagian ini, berapa banyak invers kiri yang dimiliki \( \iota \)?
    \begin{answer}
      Terhadap basis-basis standar, kita memiliki
      \begin{equation*}
        \rep{\iota}{\stdbasis_2,\stdbasis_3}
        =\begin{mat}[r]
          1  &0  \\
          0  &1  \\
          0  &0 
        \end{mat}
      \end{equation*}
      dan menyusun persamaan untuk mencari invers matriks
      \begin{equation*}
        \begin{mat}
          a &b &c \\
          d &e &f
        \end{mat}
        \begin{mat}[r]
          1 &0  \\
          0 &1  \\
          0 &0
        \end{mat}
        =\begin{mat}[r]
          1 &0  \\
          0 &1 
        \end{mat}
        =\rep{\identity}{\stdbasis_2,\stdbasis_2}
      \end{equation*}
      menghasilkan suatu sistem linear.
      \begin{equation*}
        \begin{linsys}{6}
          a & & & & & & & & & & &= &1 \\
            & &b& & & & & & & & &= &0 \\
            & & & & & &d& & & & &= &0 \\
            & & & & & & & &e& & &= &1 
        \end{linsys}
      \end{equation*}
      Sistem ini memiliki takhingga banyaknya solusi dalam $a,\ldots,f$
      karena dua dari peubah-peubah tersebut sepenuhnya bebas,
      \begin{equation*}
        \set{\colvec{a \\ b \\ c \\ d \\ e \\ f}
             =\colvec[r]{1 \\ 0 \\ 0 \\ 0 \\ 1 \\ 0}
              +c\cdot \colvec[r]{0 \\ 0 \\ 1 \\ 0 \\ 0 \\ 0}
              +f\cdot \colvec[r]{0 \\ 0 \\ 0 \\ 0 \\ 0 \\ 1}
             \suchthat c,f\in\Re}
      \end{equation*}
      sehingga persamaan matriksnya memiliki takhingga banyaknya solusi.
      \begin{equation*}
        \set{\begin{mat}
               1  &0  &c  \\
               0  &1  &f
             \end{mat}
             \suchthat c,f\in\Re}
      \end{equation*}
      Dengan basis-basis masih ditetapkan sebagai $\stdbasis_2,\stdbasis_2$,
      misalnya mengambil $c=2$ dan
      $f=3$ memberikan matriks yang merepresentasikan peta berikut.
      \begin{equation*}
        \colvec{x \\ y \\ z}
        \;\mapsunder{f_{2,3}}\;
        \colvec{x+2z \\ y+3z}
      \end{equation*}
      Pemeriksaan bahwa $\composed{f_{2,3}}{\iota}$ merupakan peta identitas pada
      $\Re^2$ mudah.
    \end{answer}
  \item  
    Jika suatu matriks memiliki takhingga banyaknya invers kanan, dapatkah matriks itu memiliki
    takhingga banyaknya invers kiri?
    Haruskah demikian?
    \begin{answer}
      Berdasarkan \nearbylemma{le:LeftAndRightInvEqual}, matriks itu tidak mungkin memiliki takhingga banyaknya
      invers kiri, karena
      matriks yang memiliki invers kiri dan kanan hanya memiliki satu dari masing-masing (dan
      masing-masing itu sekaligus merupakan keduanya\Dash matriks invers kiri dan kanan
      sama).
    \end{answer}
  \item Anggap bahwa $\map{g}{V}{W}$ linear.
    Salah satu pernyataan berikut benar, sedangkan yang lain salah.
    Manakah yang benar dan manakah yang salah?
    \begin{exparts}
      \partsitem Jika $\map{f}{W}{V}$ merupakan invers kiri dari $g$, maka $f$
        harus linear.
      \partsitem Jika $\map{f}{W}{V}$ merupakan invers kanan dari $g$, maka $f$
        harus linear.
    \end{exparts}
    \begin{answer}
      \begin{exparts}
        \partsitem Benar.
          Peta itu harus linear,
          sebagaimana ditunjukkan oleh pembuktian Teorema~II.\ref{th:OOHomoEquivalence}.
        \partsitem Salah.
          Peta itu mungkin linear, tetapi tidak harus demikian.
          Perhatikan peta proyeksi $\map{\pi}{\Re^3}{\Re^2}$ yang dijelaskan
          pada awal subbagian ini.
          Definisikan $\map{\eta}{\Re^2}{\Re^3}$ sebagai berikut.
          \begin{equation*}
            \colvec{x  \\  y}\mapsto\colvec{x  \\ y \\ 1}
          \end{equation*}
          Peta itu merupakan invers kanan dari $\pi$ karena $\composed{\pi}{\eta}$
          bekerja sebagai berikut.
          \begin{equation*}
            \colvec{x  \\  y}\mapsto\colvec{x  \\ y \\ 1}\mapsto\colvec{x \\ y}
          \end{equation*}
          Peta itu tidak linear karena tidak memetakan vektor nol ke vektor
          nol.
      \end{exparts}
    \end{answer}
  \recommended \item
    Anggap bahwa \( H \) invertibel dan \( HG \) merupakan matriks nol.
    Tunjukkan bahwa \( G \) merupakan matriks nol.
    \begin{answer}
      Sifat asosiatif perkalian matriks memberikan
      \( H^{-1}(HG)=H^{-1}Z=Z \) dan juga
      \( H^{-1}(HG)=(H^{-1}H)G=IG=G \).
    \end{answer}
  \item 
    Buktikan bahwa jika \( H \) invertibel, maka
    inversnya komutatif dengan suatu matriks \( GH^{-1}=H^{-1}G \)
    jika dan hanya jika $H$ sendiri komutatif dengan matriks tersebut \( GH=HG \).
    \begin{answer}
       Kalikan kedua ruas persamaan pertama dengan $H$.
    \end{answer}
  \recommended \item
    Tunjukkan bahwa jika \( T \) persegi dan \( T^4 \) merupakan matriks nol,
    maka \( (I-T)^{-1}=I+T+T^2+T^3 \).
    Perumumlah.
    \begin{answer}
      Mudah untuk memeriksa bahwa jika $I-T$ dikalikan pada kedua sisi dengan ekspresi tersebut
      (dengan menganggap bahwa $T^4$ merupakan matriks nol), hasilnya adalah
      matriks identitas.
      Perumuman yang jelas adalah bahwa jika \( T^n \) merupakan matriks nol,
      maka \( (I-T)^{-1}=I+T+T^2+\cdots+T^{n-1} \); pemeriksaannya sekali lagi
      mudah.
    \end{answer}
  \recommended \item 
    Misalkan \( D \) diagonal.
    Deskripsikan \( D^2 \), \( D^3 \), \ldots\thinspace, dan seterusnya.
    Deskripsikan \( D^{-1} \), \( D^{-2} \), \ldots\thinspace, dan seterusnya.
    Definisikan \( D^0 \) secara tepat.
    \begin{answer}
      Pangkat-pangkat matriks dibentuk dengan mengambil pangkat dari
      entri-entri diagonalnya.
      Artinya, semua entri \( D^2 \) bernilai nol kecuali entri-entri diagonal
      \( {d_{1,1}}^2 \), \( {d_{2,2}}^2 \), dan seterusnya.
      Hal ini menyarankan agar \( D^0 \) didefinisikan sebagai matriks identitas.
    \end{answer}
  \item
    Buktikan bahwa setiap matriks yang ekuivalen baris dengan suatu matriks invertibel juga
    invertibel.
    \begin{answer}
      Anggap bahwa $B$ ekuivalen baris dengan $A$ dan bahwa $A$ invertibel.
      Karena keduanya ekuivalen baris, terdapat suatu barisan langkah baris
      untuk mereduksi yang satu menjadi yang lain.
      Kita dapat melakukan reduksi tersebut dengan matriks; misalnya, $A$ dapat
      diubah oleh operasi baris menjadi $B$ sebagai $B=R_n\cdots R_1A$.
      Persamaan ini menyatakan $B$ sebagai hasil kali matriks-matriks invertibel dan,
      berdasarkan \nearbylemma{lem:ProdInvIsInv}, $B$ juga invertibel.
    \end{answer}
  \item 
    \textit{Pertanyaan pertama di bawah pernah muncul sebagai
    \nearbyexercise{exer:RankProdLeqRankFacts}.}
    \begin{exparts}
      \partsitem Tunjukkan bahwa peringkat hasil kali dua matriks kurang dari
        atau sama dengan nilai minimum peringkat keduanya.
      \partsitem Tunjukkan bahwa jika \( T \) dan \( S \) persegi, maka \( TS=I \)
         jika dan hanya jika \( ST=I \).
    \end{exparts}
    \begin{answer}
      \begin{exparts}
        \partsitem Lihat jawaban
          \nearbyexercise{exer:RankProdLeqRankFacts}.
        \partsitem Kita akan menunjukkan bahwa kedua syarat tersebut setara dengan
          syarat bahwa kedua matriks itu nonsingular.

          Karena \( T \) dan \( S \) persegi serta hasil kalinya terdefinisi,
          ukurannya sama, katakanlah \( \nbyn{n} \).
          Perhatikan arah \( TS=I \).
          Berdasarkan butir sebelumnya, peringkat \( I \) kurang dari atau sama dengan
          nilai minimum peringkat \( T \) dan peringkat \( S \).
          Namun, peringkat \( I \) adalah \( n \), sehingga peringkat \( T \) dan
          peringkat \( S \) masing-masing harus \( n \).
          Jadi, masing-masing nonsingular.

          Argumen yang sama menunjukkan bahwa \( ST=I \) mengakibatkan
          masing-masing nonsingular.
      \end{exparts}  
    \end{answer}
  \item 
    Tunjukkan bahwa invers matriks permutasi adalah transposenya.
    \begin{answer}
      Invers bersifat tunggal, sehingga kita hanya perlu menunjukkan bahwa transposenya memenuhi.
      Pemeriksaannya muncul di atas sebagai
      \nearbyexercise{exer:PermTimesTransEqId}.  
    \end{answer}
  \item 
    % \textit{The first two parts of this question appeared as
    % \nearbyexercise{exer:TranspAndMult}.}
    \begin{exparts}
      \partsitem Tunjukkan bahwa \( \trans{(GH)}=\trans{H}\trans{G} \).
      \partsitem Suatu matriks persegi disebut {\em simetris\/} jika setiap
        entri \( i,j \) sama dengan
        entri \( j,i \) (artinya, jika matriks tersebut sama dengan transposenya).
        Tunjukkan bahwa
        matriks \( H\trans{H} \) dan \( \trans{H}H \) simetris.
      \partsitem Tunjukkan bahwa invers dari transpos adalah transpos
        dari invers.
      \partsitem Tunjukkan bahwa invers suatu matriks simetris juga simetris.
    \end{exparts}
    \begin{answer}
      \begin{exparts}
        \partsitem Lihat jawaban \nearbyexercise{exer:TranspAndMult}.
        \partsitem Lihat jawaban \nearbyexercise{exer:TranspAndMult}.
        \partsitem Terapkan bagian pertama pada
          \( I=AA^{-1} \) untuk memperoleh
          $I=\trans{I}=\trans{(AA^{-1})}=\trans{(A^{-1})}\trans{A}$.
        \partsitem Terapkan butir sebelumnya dengan \( \trans{A}=A \),
          karena \( A \) simetris.
      \end{exparts}  
    \end{answer}
  \recommended \item
    % \textit{The items starting this question appeared as
    % \nearbyexercise{exer:ZeroDivisor}.}
    \begin{exparts}
      \partsitem Buktikan bahwa komposisi proyeksi
        \( \map{\pi_x,\pi_y}{\Re^3}{\Re^3} \) merupakan peta nol meskipun
        keduanya bukan peta nol.
      \partsitem Buktikan bahwa komposisi turunan
        \( \map{d^2/dx^2,\,d^3/dx^3}{\polyspace_4}{\polyspace_4} \)
        merupakan peta nol meskipun kedua peta bukan peta nol.
      \partsitem Berikan persamaan matriks yang merepresentasikan masing-masing dari dua
        butir sebelumnya.
    \end{exparts}
    Jika dua objek dikalikan dan menghasilkan nol meskipun keduanya tidak nol, masing-masing
    disebut \definend{pembagi nol}.\index{pembagi nol}
    Buktikan bahwa tidak ada pembagi nol yang invertibel.
    \begin{answer}
      Untuk jawaban butir-butir pada bagian pertama, lihat
      \nearbyexercise{exer:ZeroDivisor}.
      Untuk pembuktian pada bagian kedua, anggap bahwa $A$ merupakan pembagi nol sehingga
      terdapat matriks taknol $B$ dengan $AB=Z$
      (atau $BA=Z$; kasus ini serupa).
      Jika $A$ invertibel,
      maka $A^{-1}(AB)=(A^{-1}A)B=IB=B$, tetapi juga
      $A^{-1}(AB)=A^{-1}Z=Z$, bertentangan dengan fakta bahwa $B$ taknol.
    \end{answer}
  \item 
    Dalam aljabar bilangan real,
    persamaan kuadrat memiliki paling banyak dua solusi.
    Aljabar matriks berbeda.
    Tunjukkan bahwa persamaan matriks \( \nbyn{2} \), \( T^2=I \),
    memiliki lebih dari dua solusi.
    \begin{answer}
      Ada takhingga banyaknya matriks $\nbyn{2}$ $T$ yang kuadratnya
      sama dengan identitas.
      Berikut empat di antaranya.
      \begin{equation*}
        \begin{mat}[r]
          1  &0  \\
          0      &1
        \end{mat}
        \quad
        \begin{mat}[r]
          -1  &0  \\
          0   &-1
        \end{mat}
        \quad
        \begin{mat}[r]
          1  &0  \\
          0  &-1
        \end{mat}
        \quad
        \begin{mat}[r]
          -1  &0  \\
          0    &1
        \end{mat}
      \end{equation*}
      Dua lagi adalah
      \begin{equation*}
        \begin{mat}[r]
          0  &1  \\
          1  &0
        \end{mat}
        \quad
        \begin{mat}[r]
          0  &-1  \\
          -1 &0
        \end{mat}      
      \end{equation*}
      dan berikut satu lagi,
      \begin{equation*}
        \frac{1}{5}
        \begin{mat}[r]
          3  &4  \\
          4  &-3
        \end{mat}
      \end{equation*}
      (pada contoh terakhir ini, polanya melibatkan tripel Pythagoras, yaitu bilangan-bilangan
      $r,s,t\in\R$ sedemikian sehingga $r^2+s^2=t^2$).
      \textit{Catatan:}~lihat juga
      \url{https://en.wikipedia.org/wiki/Square_root_of_a_2_by_2_matrix}.
    \end{answer}
  \item 
   Apakah relasi `merupakan invers dua sisi dari' bersifat transitif?
   Refleksif?
   Simetris?
   \begin{answer}
     Relasi ini tidak refleksif karena, sebagai contoh,
     \begin{equation*}
       H=\begin{mat}[r]
         1  &0  \\
         0  &2
       \end{mat}
     \end{equation*}
     bukan invers dua sisi dari dirinya sendiri.
     Contoh yang sama menunjukkan bahwa relasi ini tidak transitif.
     Matriks itu memiliki invers dua sisi berikut.
     \begin{equation*}
       G=\begin{mat}[r]
         1  &0  \\
         0  &1/2
       \end{mat}
     \end{equation*}
     Meskipun \( H \) merupakan invers dua sisi dari \( G \) dan \( G \)
     merupakan invers dua sisi dari \( H \), kita mengetahui bahwa \( H \) bukan invers dua sisi
     dari \( H \).
     Namun, relasi ini simetris:~jika \( G \) merupakan invers dua sisi dari
     \( H \), maka
     \( GH=I=HG \), sehingga \( H \) juga merupakan invers dua sisi
     dari \( G \).
   \end{answer}
  \item  
    \cite{Monthly51p614}
    Buktikan: jika jumlah elemen pada setiap baris suatu matriks persegi
    adalah \( k \), maka jumlah elemen pada setiap baris matriks
    inversnya adalah \( 1/k \).
    \begin{answer}
      \answerasgiven %
      Misalkan \( A \) adalah matriks \( \nbyn{m} \) nonsingular dengan sifat tersebut.
      Misalkan \( B \) inversnya.
      Maka, untuk \( n\leq m \),
      \begin{equation*}
        1
        =\sum_{r=1}^{m}\delta_{nr}
        =\sum_{r=1}^{m}\sum_{s=1}^{m}b_{ns}a_{sr}
        =\sum_{s=1}^{m}\sum_{r=1}^{m}b_{ns}a_{sr}
        =k\sum_{s=1}^{m}b_{ns}
      \end{equation*}
      (\( A \) singular jika \( k=0 \)).
   \end{answer}
\end{exercises}
\index{matriks!invers|)}
